% Fonction logarithme népérien — Mathématiques Tle A — Cours signé OnBuch+
% Programme officiel MINESEC. Compiler avec : tectonic MA06-logarithme-neperien-a.tex
\newcommand{\DOCMATIERE}{Mathématiques}
\newcommand{\DOCNIVEAU}{Tle A}
\newcommand{\DOCMODULE}{Relations et opérations fondamentales dans l'ensemble des nombres réels}
\newcommand{\DOCLECON}{A06}
\newcommand{\DOCTITRE}{Fonction logarithme népérien}
% OnBuch+ — préambule « cours signés » ENRICHI (compilé avec tectonic / XeLaTeX)
% Système visuel « Pop » d'OnBuch+ : fond crème (jamais blanc), bordures encre
% épaisses, ombres « dures » décalées, titres Archivo Black en capitales.
% Version MATHÉMATIQUES : graphes (pgfplots/tikz), boîtes pédagogiques
% (définition, propriété, méthode, attention, le savais-tu, exercice résolu,
% exercice, TP, bilan), page de garde et signature OnBuch+ ancrée en bas.
%
% Macros attendues dans main.tex AVANT \input{preamble} :
%   \DOCMATIERE  \DOCNIVEAU  \DOCMODULE  \DOCLECON (numéro)  \DOCTITRE
\documentclass[11pt]{article}

\usepackage{fontspec}
\usepackage[french]{babel}
\usepackage[a4paper,margin=2cm,top=3.4cm,bottom=2.8cm,headheight=1.4cm,footskip=1.3cm]{geometry}
\usepackage{amsmath,amssymb,amsfonts}
\usepackage{mathtools} % \xmapsto, \coloneqq... souvent utilisés par les modèles
\usepackage{cancel} % simplifications barrées
% \abs{z} (module, valeur absolue) et \norm{u} (norme), utilisés par les modèles
\providecommand{\abs}{}\let\abs\relax\DeclarePairedDelimiter{\abs}{\lvert}{\rvert}
\providecommand{\norm}{}\let\norm\relax\DeclarePairedDelimiter{\norm}{\lVert}{\rVert}
\usepackage[table]{xcolor} % [table] : fournit aussi \rowcolors (alternance de lignes)
\usepackage{pagecolor}
\usepackage{tikz}
\usetikzlibrary{arrows.meta,positioning,calc,shapes.geometric,shapes.misc,shapes.arrows,decorations.pathmorphing,decorations.markings,patterns,shadows,mindmap,trees,fit,backgrounds,matrix,intersections}
\usepackage{pgfplots}
\usepackage{pgf-pie} % diagrammes circulaires \pie (statistiques)
\pgfplotsset{compat=1.18}
\usepgfplotslibrary{fillbetween,groupplots} % aires entre courbes (calcul intégral)
\usepackage{enumitem}
\usepackage{fancyhdr}
% [most] charge toutes les bibliothèques tcolorbox (listings, minted, raster,
% documentation...) et gonfle inutilement la mémoire TeX sur un document long
% (~300 boîtes) : on ne charge que ce qui est réellement utilisé.
\usepackage[skins,breakable]{tcolorbox}
\usepackage{graphicx}
\usepackage{colortbl}
\usepackage{booktabs}
\usepackage{tabularx}
\usepackage{diagbox} % case d'en-tête diagonale des tableaux à double entrée
% Colonnes extensibles centrée / gauche / droite (C, L, R), souvent utilisées par les modèles
\newcolumntype{C}{>{\centering\arraybackslash}X}
\newcolumntype{L}{>{\raggedright\arraybackslash}X}
\newcolumntype{R}{>{\raggedleft\arraybackslash}X}
\usepackage{array}
\usepackage{multicol}
\usepackage{siunitx}
% parse-numbers=false : accepte aussi \SI{2,5 \times 10^{-3}}{...}
\sisetup{locale=FR,output-decimal-marker={,},group-minimum-digits=4,parse-numbers=false}
\DeclareSIUnit\ohm{\ensuremath{\symup{\Omega}}} % U+2126 absent de la police
\usepackage{qrcode}
% \degree : commande fréquemment utilisée par les modèles pour un angle ou une
% température en dehors de \SI{}{} (non fournie par les paquets chargés).
\providecommand{\degree}{^{\circ}}
% \qed (fin de démonstration), utilisé par les modèles sans amsthm
\providecommand{\qed}{\hfill$\square$}
% \barre : double barre des valeurs interdites dans un tableau de variations (tkz-tab)
\providecommand{\barre}{\ensuremath{\|}}
% environnement proof (amsthm non chargé) utilisé par les modèles
\ifdefined\proof\else\newenvironment{proof}[1][Démonstration]{\par\noindent\textit{#1.}\ }{\qed\par}\fi
\usepackage{lastpage}
\usepackage{hyperref}
\hypersetup{hidelinks}

% ============================================================
% Palette « Pop » OnBuch+ — mêmes hex que la classe P (plus_ui.dart)
% ============================================================
\definecolor{poporange}{HTML}{F59321}
\definecolor{popdark}{HTML}{E07A0C}
\definecolor{popcream}{HTML}{FFF6E8}
\definecolor{popink}{HTML}{1C1714}
\definecolor{poppurple}{HTML}{7A5AE0}
\definecolor{popgreen}{HTML}{2E9E6A}
\definecolor{poppink}{HTML}{E0457B}
\definecolor{popblue}{HTML}{2F7DE1}
\definecolor{popgold}{HTML}{C98A12}
\definecolor{popmuted}{HTML}{8A7F76}
\definecolor{popline}{HTML}{EADFD2}
\definecolor{popgray}{HTML}{8A7F76}
\colorlet{popgrey}{popgray}
% Alias tolérants (le modèle nomme parfois les couleurs différemment)
\colorlet{popwhite}{white}
\colorlet{popblack}{popink}
\colorlet{popred}{poppink}
\colorlet{popyellow}{popgold}
% Teintes claires pour fonds de tableaux / zones
\colorlet{popblueL}{popblue!10!white}
\colorlet{poporangeL}{poporange!14!white}
\colorlet{popgreenL}{popgreen!12!white}
\colorlet{poppurpleL}{poppurple!12!white}
\colorlet{poppinkL}{poppink!12!white}
\colorlet{popgoldL}{popgold!14!white}

\pagecolor{popcream}

% ============================================================
% Typographie
% ============================================================
\defaultfontfeatures{Ligatures=TeX}
\setmainfont{PlusJakartaSans-Medium.ttf}[Path=../../../fonts/,BoldFont=PlusJakartaSans-Bold.ttf]
% Maths en sans-serif (Fira Math) avec chiffres et lettres droites de Plus
% Jakarta, pour que les formules se fondent dans le texte.
\usepackage[math-style=ISO,bold-style=ISO]{unicode-math}
\setmathfont{FiraMath-Regular.otf}[Path=../../../fonts/]
\setmathfont{PlusJakartaSans-Medium.ttf}[Path=../../../fonts/,range={up/num,up/{Latin,latin},\mathup}]
\usepackage{icomma} % virgule décimale sans espace en mode math
% Caractères mathématiques Unicode absents des polices de texte (Plus Jakarta,
% Archivo Black) : rendus par leur équivalent mathématique, en texte comme en
% formule — sinon ils s'affichent en carré vide (ex. « Arithmétique dans ℤ »).
\usepackage{newunicodechar}
\newunicodechar{ℝ}{\ensuremath{\mathbb{R}}}
\newunicodechar{ℤ}{\ensuremath{\mathbb{Z}}}
\newunicodechar{ℂ}{\ensuremath{\mathbb{C}}}
\newunicodechar{ℕ}{\ensuremath{\mathbb{N}}}
\newunicodechar{ℚ}{\ensuremath{\mathbb{Q}}}
\newunicodechar{𝔻}{\ensuremath{\mathbb{D}}}
\newunicodechar{α}{\ensuremath{\alpha}}
\newunicodechar{β}{\ensuremath{\beta}}
\newunicodechar{γ}{\ensuremath{\gamma}}
\newunicodechar{δ}{\ensuremath{\delta}}
\newunicodechar{ε}{\ensuremath{\varepsilon}}
\newunicodechar{θ}{\ensuremath{\theta}}
\newunicodechar{λ}{\ensuremath{\lambda}}
\newunicodechar{μ}{\ensuremath{\mu}}
\newunicodechar{σ}{\ensuremath{\sigma}}
\newunicodechar{τ}{\ensuremath{\tau}}
\newunicodechar{φ}{\ensuremath{\varphi}}
\newunicodechar{ψ}{\ensuremath{\psi}}
\newunicodechar{ω}{\ensuremath{\omega}}
\newunicodechar{Σ}{\ensuremath{\Sigma}}
% majuscules produites par \MakeUppercase dans les titres (α-aminés -> Α)
\newunicodechar{Α}{\ensuremath{\alpha}}
\newunicodechar{Β}{\ensuremath{\beta}}
\newunicodechar{Γ}{\ensuremath{\Gamma}}
\newunicodechar{Δ}{\ensuremath{\Delta}}
\newunicodechar{Ε}{\ensuremath{\varepsilon}}
\newunicodechar{Θ}{\ensuremath{\Theta}}
\newunicodechar{Λ}{\ensuremath{\Lambda}}
\newunicodechar{Μ}{\ensuremath{\mu}}
\newunicodechar{Τ}{\ensuremath{\tau}}
\newunicodechar{Φ}{\ensuremath{\Phi}}
\newunicodechar{Ψ}{\ensuremath{\Psi}}
\newunicodechar{Ω}{\ensuremath{\Omega}}
\newunicodechar{↦}{\ensuremath{\mapsto}}
\newunicodechar{∈}{\ensuremath{\in}}
\newunicodechar{∉}{\ensuremath{\notin}}
\newunicodechar{∧}{\ensuremath{\wedge}}
\newunicodechar{⇔}{\ensuremath{\Leftrightarrow}}
\newunicodechar{⟺}{\ensuremath{\Longleftrightarrow}}
\newunicodechar{⇒}{\ensuremath{\Rightarrow}}
\newunicodechar{⟹}{\ensuremath{\Longrightarrow}}
\newunicodechar{∘}{\ensuremath{\circ}}
\newunicodechar{∪}{\ensuremath{\cup}}
\newunicodechar{∩}{\ensuremath{\cap}}
\newunicodechar{≡}{\ensuremath{\equiv}}
\newunicodechar{⊂}{\ensuremath{\subset}}
\newunicodechar{⊥}{\ensuremath{\perp}}
\newunicodechar{∃}{\ensuremath{\exists}}
\newunicodechar{∀}{\ensuremath{\forall}}
\newunicodechar{∛}{\ensuremath{\sqrt[3]{\,}}}
\newunicodechar{∜}{\ensuremath{\sqrt[4]{\,}}}
\newunicodechar{⃗}{\ensuremath{{}^{\to}}}
\newunicodechar{ⁿ}{\ifmmode^{n}\else\textsuperscript{\ensuremath{n}}\fi}
\newunicodechar{ˣ}{\ifmmode^{x}\else\textsuperscript{\ensuremath{x}}\fi}
\newunicodechar{ᵇ}{\ifmmode^{b}\else\textsuperscript{\ensuremath{b}}\fi}
\newunicodechar{ʳ}{\ifmmode^{r}\else\textsuperscript{\ensuremath{r}}\fi}
\newunicodechar{ᵘ}{\ifmmode^{u}\else\textsuperscript{\ensuremath{u}}\fi}
\newunicodechar{ʸ}{\ifmmode^{y}\else\textsuperscript{\ensuremath{y}}\fi}
\newunicodechar{ᵃ}{\ifmmode^{a}\else\textsuperscript{\ensuremath{a}}\fi}
\newunicodechar{ᵉ}{\ifmmode^{e}\else\textsuperscript{\ensuremath{e}}\fi}
\newunicodechar{ᵏ}{\ifmmode^{k}\else\textsuperscript{\ensuremath{k}}\fi}
\newunicodechar{ᵖ}{\ifmmode^{p}\else\textsuperscript{\ensuremath{p}}\fi}
\newunicodechar{ᴺ}{\ifmmode^{N}\else\textsuperscript{\ensuremath{N}}\fi}
\newunicodechar{ᶜ}{\ifmmode^{c}\else\textsuperscript{\ensuremath{c}}\fi}
\newunicodechar{ʰ}{\ifmmode^{h}\else\textsuperscript{\ensuremath{h}}\fi}
\newunicodechar{ᵠ}{\ifmmode^{q}\else\textsuperscript{\ensuremath{q}}\fi}
\newunicodechar{ₙ}{\ifmmode_{n}\else\textsubscript{\ensuremath{n}}\fi}
\newunicodechar{ₐ}{\ifmmode_{a}\else\textsubscript{\ensuremath{a}}\fi}
\newunicodechar{ᵢ}{\ifmmode_{i}\else\textsubscript{\ensuremath{i}}\fi}
\newunicodechar{ᵣ}{\ifmmode_{r}\else\textsubscript{\ensuremath{r}}\fi}
\newunicodechar{ₖ}{\ifmmode_{k}\else\textsubscript{\ensuremath{k}}\fi}
\newunicodechar{ᵦ}{\ifmmode_{\beta}\else\textsubscript{\ensuremath{\beta}}\fi}
\newunicodechar{ₓ}{\ifmmode_{x}\else\textsubscript{\ensuremath{x}}\fi}
\newfontfamily\popdisplay{ArchivoBlack-Regular.ttf}[Path=../../../fonts/]
\newfontfamily\poplogo{SpaceGrotesk-Bold.ttf}[Path=../../../fonts/]
\newfontfamily\popsemi{PlusJakartaSans-SemiBold.ttf}[Path=../../../fonts/]
\newfontfamily\popxbold{PlusJakartaSans-ExtraBold.ttf}[Path=../../../fonts/]
\newfontfamily\popxboldls{PlusJakartaSans-ExtraBold.ttf}[Path=../../../fonts/,LetterSpace=3.0]

\setlist[enumerate]{leftmargin=1.7em,itemsep=5pt,topsep=4pt}
\setlist[itemize]{leftmargin=1.7em,itemsep=4pt,topsep=4pt,label={\color{poporange}\textbullet}}
\setlength{\parindent}{0pt}
\setlength{\parskip}{5pt}
\renewcommand{\arraystretch}{1.25}

% pgfplots : style Pop par défaut
\pgfplotsset{
  every axis/.append style={axis line style={popink,line width=1pt},tick style={popink},
    grid style={popline,line width=0.5pt},label style={font=\small},tick label style={font=\footnotesize}},
  popcurve/.style={poporange,line width=1.8pt,no markers,smooth},
}
% Même style disponible dans un \draw TikZ simple (hors pgfplots)
\tikzset{popcurve/.style={poporange,line width=1.8pt}}

% ============================================================
% Pill / squiggle
% ============================================================
\newcommand{\poppill}[3]{%
  \tikz[baseline=-0.55ex]\node[draw=popink,line width=1.2pt,rounded corners=2.6pt,%
    fill=#2,inner xsep=6.5pt,inner ysep=3pt]{%
    \popxboldls\fontsize{7}{7}\selectfont\color{#3}\MakeUppercase{#1}};%
}
\newcommand{\popsquiggle}[1]{%
  \tikz[baseline=-0.4ex]{
    \draw[#1,line width=2.6pt,line cap=round]
      plot [smooth] coordinates {(0,0.09) (0.34,0.01) (0.68,0.21) (1.02,0.01) (1.36,0.10)};
  }%
}
% Mot-clé surligné (marqueur orange sous le texte)
\newcommand{\cle}[1]{{\popxbold\color{popdark}#1}}

% ============================================================
% En-tête / pied
% ============================================================
\pagestyle{fancy}
\fancyhf{}
\renewcommand{\headrulewidth}{0pt}
\renewcommand{\footrulewidth}{0pt}
\lhead{\raisebox{-0.15ex}{\poplogo\fontsize{13}{13}\selectfont\color{popink}OnBuch\color{poporange}+}}
\rhead{\poppill{\DOCMATIERE\ \textbullet\ \DOCNIVEAU\ \textbullet\ Leçon \DOCLECON}{white}{popink}}
\lfoot{\popxbold\fontsize{7.6}{7.6}\selectfont\color{popmuted}\MakeUppercase{Cours signé OnBuch+}\ \textbullet\ {\popsemi\DOCTITRE}}
\rfoot{\tikz[baseline=-0.6ex]\node[rounded corners=8.5pt,fill=popink,minimum height=17pt,minimum width=17pt,inner xsep=5pt,inner sep=0]{\popxbold\fontsize{8}{8}\selectfont\color{popcream}\thepage\,/\,\pageref*{LastPage}};}

% ============================================================
% Titres
% ============================================================
\newcommand{\chaptitre}[2]{%
  \begin{tcolorbox}[enhanced,colback=poporange,colframe=popink,boxrule=2.6pt,arc=11pt,
      drop shadow=popink,left=15pt,right=15pt,top=12pt,bottom=11pt,before skip=3pt,after skip=18pt]
    \poppill{#2}{popink}{popcream}\par\vspace{9pt}
    {\raggedright\hyphenpenalty=10000\exhyphenpenalty=10000\popdisplay\fontsize{22}{24}\selectfont\color{popcream}\MakeUppercase{#1}\par}
  \end{tcolorbox}
}
% Section numérotée : pastille orange avec le numéro + titre Archivo
\newcounter{popsec}
\newcommand{\coursec}[1]{%
  \stepcounter{popsec}\setcounter{popsub}{0}%
  \par\vspace{16pt}\noindent
  \begin{minipage}{\linewidth}
  \tikz[baseline=-0.7ex]\node[draw=popink,line width=1.6pt,fill=poporange,rounded corners=4pt,minimum size=22pt,inner sep=0]{\popdisplay\fontsize{12}{12}\selectfont\color{popcream}\thepopsec};%
  \hspace{8pt}{\popdisplay\fontsize{15}{17}\selectfont\color{popink}\MakeUppercase{#1}}\par
  \vspace{4pt}\noindent{\color{poporange}\rule{36pt}{3.6pt}}
  \end{minipage}\par\nopagebreak\vspace{8pt}
}
\newcounter{popsub}
\newcommand{\courssub}[1]{%
  \stepcounter{popsub}%
  \par\vspace{9pt}\noindent{\popxbold\fontsize{11.5}{13}\selectfont\color{popink}{\color{poporange}\thepopsec.\thepopsub}\hspace{6pt}#1}\par\nopagebreak\vspace{4pt}
}

% ============================================================
% Boîtes pédagogiques — même recette : fond blanc, bordure épaisse colorée,
% ombre dure encre, badge plein. Titre optionnel [#1].
% ============================================================
\tcbset{popbase/.style={breakable,enhanced,colback=white,boxrule=2.3pt,arc=9pt,
  shadow={3pt}{-3pt}{0pt}{popink},left=14pt,right=14pt,top=11pt,bottom=11pt,before skip=11pt,after skip=15pt,
  fontupper=\popsemi\fontsize{10.6}{14.5}\selectfont\color{popink},
  fontlower=\popsemi\fontsize{10.6}{14.5}\selectfont\color{popink}}}
% Coche dessinée (le glyphe ✓ n'existe pas dans les polices du document)
\newcommand{\popcheck}[1]{\tikz[baseline=-0.5ex]\draw[#1,line width=1.6pt,line cap=round,line join=round](-0.13,0.0)--(-0.04,-0.09)--(0.13,0.1);}
\newcommand{\popboxhead}[3]{\poppill{#1}{#2}{white}\ifx\relax#3\relax\else\hspace{6pt}{\popxbold\fontsize{10.6}{12}\selectfont\color{popink}#3}\fi\par\vspace{7pt}}

\newtcolorbox{aretenir}[1][]{popbase,colframe=popgreen,before upper={\popboxhead{À retenir}{popgreen}{#1}}}
\newtcolorbox{exemplebox}[1][]{popbase,colframe=poporange,before upper={\popboxhead{Exemple}{poporange}{#1}}}
\newtcolorbox{definition}[1][]{popbase,colframe=popblue,before upper={\popboxhead{Définition}{popblue}{#1}}}
\newtcolorbox{propriete}[1][]{popbase,colframe=poppurple,before upper={\popboxhead{Propriété}{poppurple}{#1}}}
\newtcolorbox{methode}[1][]{popbase,colframe=popgold,before upper={\popboxhead{Méthode}{popgold}{#1}}}
\newtcolorbox{attention}[1][]{popbase,colframe=poppink,before upper={\popboxhead{Attention}{poppink}{#1}}}
\newtcolorbox{experience}[1][]{popbase,colframe=popblue,colback=popblueL,before upper={\popboxhead{Expérience}{popblue}{#1}}}
% « Le savais-tu ? » : carte violette pleine (contexte, culture, Cameroun)
\newtcolorbox{savaistu}[1][]{popbase,colback=poppurple,colframe=popink,
  fontupper=\popsemi\fontsize{10.4}{14.2}\selectfont\color{white},
  before upper={\poppill{Le savais-tu\,?}{white}{poppurple}\ifx\relax#1\relax\else\hspace{6pt}{\popxbold\color{white}#1}\fi\par\vspace{7pt}}}
% Exercice résolu : énoncé en haut, \tcblower puis la solution sur fond crème
\newcounter{popexo}\newcounter{popexr}
\newtcolorbox{exoresolu}[1][]{popbase,colframe=poporange,colbacklower=popcream,
  segmentation style={popink,line width=1.2pt,dashed},
  before upper={\stepcounter{popexr}\popboxhead{Exercice résolu \thepopexr}{poporange}{#1}},
  before lower={\poppill{Solution}{popink}{popcream}\par\vspace{6pt}}}
% Exercice (non résolu en place) — la correction est dans la partie Corrigés
\newtcolorbox{exercice}[1][]{popbase,colframe=popink,
  before upper={\stepcounter{popexo}\popboxhead{Exercice \thepopexo}{popink}{#1}}}
% Corrigé
\newtcolorbox{corrige}[1][]{popbase,colframe=popgreen,colback=popgreenL,
  before upper={\popboxhead{Corrigé}{popgreen}{#1}}}
% Objectifs / prérequis / bilan
\newtcolorbox{objectifs}{popbase,colframe=popink,colback=poporangeL,
  before upper={\popboxhead{Objectifs de la leçon}{popink}{}}}
\newtcolorbox{prerequis}{popbase,colframe=popink,colback=popblueL,
  before upper={\popboxhead{Prérequis}{popblue}{}}}
\newtcolorbox{bilan}{popbase,colframe=popink,colback=white,boxrule=2.8pt,
  before upper={\popboxhead{Fiche bilan}{poporange}{L'essentiel à connaître pour l'examen}}}
% Figure encadrée avec légende
\newtcolorbox{popfigure}[1][]{enhanced,breakable=false,colback=white,colframe=popline,boxrule=1.4pt,arc=8pt,
  left=10pt,right=10pt,top=10pt,bottom=8pt,before skip=10pt,after skip=12pt,halign=center,
  lower separated=false,halign lower=center,
  fontlower=\popsemi\fontsize{9}{11.5}\selectfont\color{popmuted},
  #1}
\newcommand{\legende}[1]{\par\vspace{6pt}{\popsemi\fontsize{9}{11.5}\selectfont\color{popmuted}#1}}

% ============================================================
% Signature OnBuch+
% ============================================================
\newcommand{\poplogoinline}[1]{{\poplogo\fontsize{#1}{#1}\selectfont\color{popink}OnBuch\color{poporange}+}}
\newcommand{\signatureonbuch}{%
  \begin{tcolorbox}[enhanced,colback=popink,colframe=popink,boxrule=2.2pt,arc=10pt,
      drop shadow=poporange,left=14pt,right=14pt,top=10pt,bottom=10pt,before skip=0pt,after skip=0pt]
    \tikz[baseline=-0.7ex]\node[circle,fill=popgreen,draw=popcream,line width=1.3pt,minimum size=22pt,inner sep=0]{\popcheck{white}};%
    \hspace{9pt}\begin{minipage}[c]{0.62\linewidth}
      {\popxbold\fontsize{12}{13}\selectfont\color{popcream}Signé\ \textbullet\ {\poplogo OnBuch\color{poporange}+}}\par\vspace{2pt}
      {\raggedright\popsemi\fontsize{8}{10.5}\selectfont\color{popline}\MakeUppercase{Équipe pédagogique OnBuch+}\\\MakeUppercase{Conforme au programme officiel MINESEC}\par}
    \end{minipage}\hfill
    {\poplogo\fontsize{22}{22}\selectfont\color{popcream}OnBuch\color{poporange}+}
  \end{tcolorbox}%
}

% ============================================================
% Page de garde
% #1 titre · #2 matière · #3 niveau · #4 module · #5 n° leçon · #6 durée ·
% #7 description / objectifs (texte ou itemize)
% ============================================================
\newcommand{\pagegarde}[7]{%
  \thispagestyle{empty}
  \newgeometry{a4paper,margin=1.7cm,top=1.6cm,bottom=1.4cm}%
  \begin{tikzpicture}[remember picture,overlay]
    \fill[poporange!18] ([xshift=-3.2cm,yshift=-3.6cm]current page.north east) circle (4.6cm);
    \fill[poppurple!14] ([xshift=2.4cm,yshift=7.6cm]current page.south west) circle (3.8cm);
  \end{tikzpicture}%
  {\poplogo\fontsize{30}{30}\selectfont\color{popink}OnBuch\color{poporange}+}\hfill
  \raisebox{0.6ex}{\poppill{Cours signé \textbullet\ Premium}{popink}{popcream}}\par
  \vspace{1pt}\popsquiggle{poppurple}\par
  \vspace{8pt}
  \begin{tcolorbox}[enhanced,colback=poporange,colframe=popink,boxrule=2.8pt,arc=13pt,
      drop shadow=popink,left=18pt,right=18pt,top=12pt,bottom=13pt,after skip=10pt]
    \poppill{#2 \textbullet\ #3}{popink}{popcream}\hspace{5pt}\poppill{#4}{white}{popink}\par\vspace{8pt}
    {\popdisplay\fontsize{14}{14}\selectfont\color{popink}LEÇON #5}\par\vspace{4pt}
    {\raggedright\hyphenpenalty=10000\exhyphenpenalty=10000\popdisplay\fontsize{25}{27}\selectfont\color{popcream}\MakeUppercase{#1}\par}
    \vspace{8pt}
    {\popxbold\fontsize{9.5}{9.5}\selectfont\color{popink}\MakeUppercase{Durée conseillée : #6}}
  \end{tcolorbox}
  \begin{tcolorbox}[enhanced,colback=white,colframe=popink,boxrule=2.2pt,arc=10pt,
      drop shadow=popink,left=16pt,right=16pt,top=10pt,bottom=10pt,after skip=8pt]
    \poppill{Dans cette leçon}{poppurple}{white}\par\vspace{5pt}
    \popsemi\fontsize{9.3}{12.6}\selectfont\color{popink}#7
  \end{tcolorbox}
  \noindent\begin{minipage}[t]{0.487\linewidth}
    \begin{tcolorbox}[enhanced,colback=popgreen,colframe=popink,boxrule=2.2pt,arc=10pt,
        drop shadow=popink,left=11pt,right=11pt,top=10pt,bottom=10pt,height=3.1cm,valign=center]
      \tcbox[colback=white,colframe=popink,boxrule=1.3pt,arc=5pt,boxsep=0pt,nobeforeafter,
        left=3pt,right=3pt,top=3pt,bottom=3pt,valign=center]{\qrcode[height=1.9cm]{https://play.google.com/store/apps/details?id=com.luvvix.onbuch&hl=fr}}%
      \hspace{8pt}\begin{minipage}[c]{0.5\linewidth}
        {\popdisplay\fontsize{11}{12.5}\selectfont\color{white}\MakeUppercase{Télécharge OnBuch}}\par\vspace{4pt}
        {\raggedright\popsemi\fontsize{8.4}{11}\selectfont\color{white}Gratuit \textbullet\ Google Play\\Tuteur IA, annales, concours, résultats\par}
      \end{minipage}
    \end{tcolorbox}
  \end{minipage}\hfill
  \begin{minipage}[t]{0.487\linewidth}
    \begin{tcolorbox}[enhanced,colback=poppurple,colframe=popink,boxrule=2.2pt,arc=10pt,
        drop shadow=popink,left=11pt,right=11pt,top=10pt,bottom=10pt,height=3.1cm,valign=center]
      \tikz[baseline=-0.7ex]\node[circle,fill=white,draw=popink,line width=1.3pt,minimum size=1.6cm,inner sep=0]{\popdisplay\fontsize{24}{24}\selectfont\color{poppurple}+};%
      \hspace{8pt}\begin{minipage}[c]{0.6\linewidth}
        {\popdisplay\fontsize{11}{12.5}\selectfont\color{white}\MakeUppercase{Passe à OnBuch+}}\par\vspace{4pt}
        {\raggedright\popsemi\fontsize{8.4}{11}\selectfont\color{white}Tous les cours signés, Tuteur IA illimité, fiches PDF — onglet OnBuch+ de l'app\par}
      \end{minipage}
    \end{tcolorbox}
  \end{minipage}
  \vspace{8pt}
  \popcontenu
  \vfill
  \signatureonbuch
  \restoregeometry
  \newpage
}

% Rangée « Ce cours contient » (page de garde)
\newcommand{\poptile}[3]{% #1 couleur #2 gros texte #3 légende
  \begin{tcolorbox}[enhanced,colback=#1,colframe=popink,boxrule=2pt,arc=9pt,shadow={2.5pt}{-2.5pt}{0pt}{popink},
      left=6pt,right=6pt,top=9pt,bottom=9pt,halign=center,valign=center,height=1.75cm,width=\linewidth,nobeforeafter]
    {\popdisplay\fontsize{12.5}{13}\selectfont\color{white}\MakeUppercase{#2}}\par\vspace{3pt}
    {\centering\popsemi\fontsize{7.8}{9.5}\selectfont\color{white}#3\par}
  \end{tcolorbox}}
\newcommand{\popcontenu}{%
  \par\noindent{\popxboldls\fontsize{7.5}{7.5}\selectfont\color{popmuted}CE COURS SIGNÉ CONTIENT}\par\vspace{7pt}
  \noindent\begin{minipage}[t]{0.235\linewidth}\poptile{popblue}{Cours}{complet, illustré, pas à pas}\end{minipage}\hfill
  \begin{minipage}[t]{0.235\linewidth}\poptile{poporange}{Méthodes}{et exercices résolus}\end{minipage}\hfill
  \begin{minipage}[t]{0.235\linewidth}\poptile{poppink}{Exercices}{type examen + corrigés}\end{minipage}\hfill
  \begin{minipage}[t]{0.235\linewidth}\poptile{popgreen}{Fiche bilan}{l'essentiel à retenir}\end{minipage}\par}

% Sommaire
\newtcolorbox{sommaire}{enhanced,colback=white,colframe=popink,boxrule=2.2pt,arc=10pt,
  drop shadow=popink,left=18pt,right=18pt,top=13pt,bottom=13pt,before skip=6pt,after skip=20pt,
  before upper={\poppill{Sommaire}{poppurple}{white}\par\vspace{9pt}\popsemi\fontsize{10.6}{14.5}\selectfont\color{popink}}}

% Signature de fin de cours — poussée tout en bas de la dernière page
\newcommand{\findecours}{%
  \par\vfill\nopagebreak
  \begin{center}\popsquiggle{poporange}\end{center}\vspace{4pt}
  \signatureonbuch
}
\AtBeginDocument{\def\llbracket{\mathopen{[\mkern-3.5mu[}}\def\rrbracket{\mathclose{]\mkern-3.5mu]}}} % intervalles d'entiers (glyphe ⟦ absent de la police)
% Glyphes absents de Fira Math : redéfinis à partir de glyphes disponibles
\AtBeginDocument{%
  \def\setminus{\mathbin{\backslash}}\let\smallsetminus\setminus
  \def\vdots{\mathord{\vbox{\baselineskip3.2pt\lineskiplimit0pt\kern4pt\hbox{.}\hbox{.}\hbox{.}}}}%
  \def\ddots{\mathinner{\mkern1mu\raise6.5pt\hbox{.}\mkern2mu\raise3.6pt\hbox{.}\mkern2mu\raise0.7pt\hbox{.}\mkern1mu}}%
  \def\triangle{\mathord{\scalebox{0.9}{$\symup{\Delta}$}}}%
  \def\nabla{\mathord{\raisebox{\depth}{\scalebox{1}[-1]{$\symup{\Delta}$}}}}%
  \def\checkmark{\popcheck{popdark}}%
  \def\star{\mathbin{\ast}}\def\bigstar{\mathord{\ast}}%
  \def\diamond{\mathbin{\scalebox{0.6}{$\Diamond$}}}%
}

%% FIGFIX-BEGIN v5 — correctifs globaux des figures (docs/onbuch-generation/tools/figfix)
\usepackage{adjustbox}\usepackage{etoolbox}\tcbuselibrary{raster}
% toute figure TikZ/pgfplots plus large que la ligne est réduite à la largeur disponible
\BeforeBeginEnvironment{tikzpicture}{\begin{adjustbox}{max width=\linewidth}}
\AfterEndEnvironment{tikzpicture}{\end{adjustbox}}
% légendes pgfplots lisibles par-dessus les courbes
\pgfplotsset{every axis/.append style={legend style={fill=white,fill opacity=0.92,text opacity=1,draw=black!15,inner sep=3pt}}}
% étiquettes d'arêtes (syntaxe edge["..."]) avec fond blanc
\tikzset{every edge quotes/.append style={fill=white,inner sep=1.2pt}}
%% FIGFIX-END
\begin{document}

\pagegarde{Fonction logarithme népérien}{Mathématiques}{Tle A}{Relations et opérations fondamentales dans l'ensemble des nombres réels}{A06}{6 h}{Cette leçon introduit la fonction logarithme népérien ln, fonction réciproque de l'exponentielle, et en établit les propriétés algébriques, les limites et la dérivée. Elle fournit les outils pour résoudre équations et inéquations, calculer des primitives de type a/(bx+c) et étudier complètement des fonctions du type ln(ax+b).
\vspace{4pt}
\begin{multicols}{2}\small
\begin{itemize}
\item À la fin de cette leçon, je sais définir la fonction ln comme réciproque de l'exponentielle et donner son ensemble de définition
\item je sais utiliser les propriétés algébriques de ln (produit, quotient, puissance, racine) pour transformer une expression
\item je sais résoudre des équations et inéquations contenant ln, en déterminant d'abord l'ensemble de validité
\item je sais résoudre des équations et inéquations du second degré en posant l'inconnue auxiliaire X = ln x
\item je sais déterminer les limites usuelles de ln (en 0+, en +∞, et les croissances comparées)
\item je sais calculer la dérivée de ln(ax+b) et de fonctions composées avec ln
\end{itemize}\end{multicols}}

\begin{sommaire}
\begin{multicols}{2}
\begin{enumerate}
\item Définition et premières propriétés de ln
\item Propriétés algébriques de ln
\item Équations, inéquations et inconnue auxiliaire
\item Limites usuelles et croissances comparées
\item Dérivation et primitives liées à ln
\item Étude complète de fonctions du type ln(ax + b)
\item Activité d'intégration
\item Méthodes et exercices résolus
\item Exercices
\item Corrigés des exercices
\item Fiche bilan
\end{enumerate}
\end{multicols}
\end{sommaire}

\begin{objectifs}
\begin{itemize}
\item À la fin de cette leçon, je sais définir la fonction ln comme réciproque de l'exponentielle et donner son ensemble de définition
\item je sais utiliser les propriétés algébriques de ln (produit, quotient, puissance, racine) pour transformer une expression
\item je sais résoudre des équations et inéquations contenant ln, en déterminant d'abord l'ensemble de validité
\item je sais résoudre des équations et inéquations du second degré en posant l'inconnue auxiliaire X = ln x
\item je sais déterminer les limites usuelles de ln (en 0+, en +∞, et les croissances comparées)
\item je sais calculer la dérivée de ln(ax+b) et de fonctions composées avec ln
\item je sais déterminer une primitive de a/(bx+c) et de fonctions faisant intervenir ln
\item je sais étudier complètement une fonction du type ln(ax+b) et tracer sa courbe avec son asymptote verticale
\end{itemize}
\end{objectifs}

\begin{prerequis}
\begin{itemize}
\item Fonction exponentielle exp : définition, propriétés, courbe (notion de début d'année de Terminale)
\item Fonctions réciproques et courbes symétriques par rapport à la droite y = x
\item Dérivation : dérivées usuelles, dérivée d'une composée (u∘v)', opérations sur les dérivées
\item Primitives : définition, primitives usuelles, linéarité
\item Limites de fonctions et asymptotes (notions de Première)
\end{itemize}
\end{prerequis}

\newpage

\coursec{Définition et premières propriétés de ln}

Un commerçant du marché Mokolo à Douala place \num{100000} FCFA à intérêts composés de $10\,\%$ l'an. Chaque année, son capital est multiplié par $1{,}1$ : après $n$ années, il dispose de $100\,000 \times (1{,}1)^n$ FCFA. Il se demande : \emph{au bout de combien d'années mon capital aura-t-il doublé ?} Il doit résoudre $(1{,}1)^n = 2$, c'est-à-dire trouver l'\emph{exposant} connaissant le résultat de l'exponentiation.

Un ingénieur de la SONARA mesure l'atténuation du son dans une salle de spectacle : l'intensité perçue décroît de façon exponentielle avec la distance. Un biologiste de l'Université de Yaoundé I suit une population de bactéries qui double toutes les 20 minutes et cherche l'instant où elle atteindra un million d'individus. Dans les trois cas, le problème est le même : il faut \cle{inverser} une exponentielle.

C'est exactement le rôle de la \cle{fonction logarithme népérien}, nouvel outil que nous allons construire rigoureusement à partir de la fonction exponentielle, puis étudier et utiliser pour résoudre équations et inéquations.

\textbf{Problématique :} la fonction exponentielle associe à tout réel $x$ un nombre strictement positif $e^x$. Peut-on « remonter » : étant donné un nombre strictement positif $a$, existe-t-il un unique réel $x$ tel que $e^x = a$ ? Si oui, comment définir et étudier cette « fonction retour » ?

\courssub{La fonction exponentielle est bijective de $\mathbb{R}$ sur $]0\,;+\infty[$}

Commençons par rappeler ce que tu sais déjà de la fonction exponentielle, étudiée dans la leçon précédente.

\begin{propriete}[Rappel sur la fonction exponentielle]
La fonction exponentielle, notée $\exp$ ou $x \mapsto e^x$, est :
\begin{itemize}
\item définie et \cle{continue} sur $\mathbb{R}$ ;
\item \cle{strictement croissante} sur $\mathbb{R}$ ;
\item à valeurs dans $]0\,;+\infty[$, avec $\lim\limits_{x \to -\infty} e^x = 0$ et $\lim\limits_{x \to +\infty} e^x = +\infty$.
\end{itemize}
\end{propriete}

Ces trois ingrédients — continuité, stricte monotonie, et le fait que l'exponentielle « balaie » tout l'intervalle $]0\,;+\infty[$ — sont exactement les hypothèses du théorème suivant, que nous admettons.

\begin{propriete}[Théorème de la bijection (admis)]
Soit $f$ une fonction \cle{continue} et \cle{strictement monotone} sur un intervalle $I$. Alors $f$ réalise une \cle{bijection} de $I$ sur l'intervalle image $J = f(I)$ : pour tout $y \in J$, l'équation $f(x) = y$ admet une \textbf{unique} solution $x$ dans $I$.

On peut donc définir la \cle{fonction réciproque} $f^{-1} : J \to I$, qui à chaque $y \in J$ associe l'unique $x \in I$ tel que $f(x) = y$. De plus, $f^{-1}$ est continue et strictement monotone sur $J$, de même sens de variation que $f$.
\end{propriete}

\begin{attention}[Bijection : un mot à bien comprendre]
Dire que $f$ est \emph{bijective} de $I$ sur $J$ signifie deux choses à la fois : chaque élément de $J$ est atteint (\emph{surjection}) et il est atteint une seule fois (\emph{injection}). C'est l'unicité qui permet de « revenir en arrière » sans ambiguïté. Sans la stricte monotonie, une équation comme $x^2 = 4$ aurait deux solutions, et on ne saurait pas laquelle choisir pour définir la réciproque.
\end{attention}

Appliquons ce théorème à l'exponentielle : elle est continue et strictement croissante sur $\mathbb{R}$, et $f(\mathbb{R}) = \,]0\,;+\infty[$. Elle réalise donc une bijection de $\mathbb{R}$ sur $]0\,;+\infty[$. Sa réciproque existe : c'est elle que nous baptisons \emph{logarithme népérien}.

\begin{popfigure}
\begin{center}
\begin{tikzpicture}[scale=1.05, >=Stealth]
% Ensembles
\draw[thick, popblue] (-4.2,0) ellipse (1.5 and 2.1);
\draw[thick, poporange] (4.2,0) ellipse (1.5 and 2.1);
\node[popblue, font=\bfseries] at (-4.2,2.7) {$\mathbb{R}$};
\node[poporange, font=\bfseries] at (4.2,2.7) {$]0\,;+\infty[$};
% Points
\fill[popblue] (-4.2,0.9) circle (2pt) node[left, popdark] {$x$};
\fill[poporange] (4.2,0.9) circle (2pt) node[right, popdark] {$e^x$};
\fill[popblue] (-4.2,-0.9) circle (2pt) node[left, popdark] {$\ln a$};
\fill[poporange] (4.2,-0.9) circle (2pt) node[right, popdark] {$a$};
% Flèches
\draw[->, very thick, popblue] (-2.6,1.3) to[bend left=20] node[midway, above, popdark]{$\exp$} (2.6,1.3);
\draw[->, very thick, poporange] (2.6,-1.3) to[bend left=20] node[midway, below, popdark]{$\ln$} (-2.6,-1.3);
\draw[->, thick, popblue!70] (-3.9,0.9) to[bend left=12] (3.9,0.9);
\draw[->, thick, poporange!80] (3.9,-0.9) to[bend left=12] (-3.9,-0.9);
\end{tikzpicture}
\end{center}
\legende{La fonction $\exp$ réalise une bijection de $\mathbb{R}$ sur $]0\,;+\infty[$ ; sa réciproque $\ln$ fait le chemin inverse.}
\end{popfigure}

\courssub{Définition du logarithme népérien}

\begin{definition}[Fonction logarithme népérien]
La fonction \cle{logarithme népérien}, notée $\ln$, est la \cle{fonction réciproque} de la fonction exponentielle. Elle est donc définie sur $]0\,;+\infty[$, à valeurs dans $\mathbb{R}$, par l'équivalence fondamentale :
\[ \text{Pour tout } x > 0 \text{ et tout } y \in \mathbb{R} : \qquad y = \ln x \iff x = e^y. \]
\end{definition}

Autrement dit, $\ln x$ est \textbf{l'unique nombre réel dont l'exponentielle vaut $x$}. Retiens cette phrase : elle contient toute la définition. Quand tu lis $\ln 5$, pense : « c'est le nombre $y$ tel que $e^y = 5$ ».

\begin{propriete}[Conséquences immédiates de la définition]
\begin{itemize}
\item $\ln 1 = 0$ \quad (car $e^0 = 1$) ;
\item $\ln e = 1$ \quad (car $e^1 = e$), où $e \approx 2{,}718$ ;
\item $\ln$ est \cle{continue} et \cle{strictement croissante} sur $]0\,;+\infty[$ (hérité de la stricte croissance de $\exp$, d'après le théorème de la bijection).
\end{itemize}
\end{propriete}

\begin{attention}[Le piège numéro 1 : le domaine de définition]
Le symbole $\ln x$ \textbf{n'a de sens que si $x > 0$}. Écrire $\ln(-3)$ ou $\ln 0$ est une erreur grave : comme $e^y > 0$ pour tout réel $y$, aucun nombre n'a pour exponentielle $-3$ ou $0$. Avant toute résolution d'équation ou d'inéquation contenant un logarithme, commence \emph{toujours} par chercher les conditions d'existence. Par exemple, $\ln(2x-6)$ n'existe que si $2x - 6 > 0$, c'est-à-dire $x > 3$.
\end{attention}

\courssub{Les relations fondamentales entre $\ln$ et $\exp$}

Dire que $\ln$ et $\exp$ sont réciproques l'une de l'autre, c'est dire que chacune « annule » l'action de l'autre. Précisons cela.

\begin{propriete}[Relations fondamentales]
\begin{itemize}
\item Pour tout réel $x > 0$ : \quad $e^{\ln x} = x$.
\item Pour tout réel $x$ : \quad $\ln\left(e^x\right) = x$.
\end{itemize}
\end{propriete}

\begin{experience}[Démonstration guidée des relations fondamentales]
\textbf{Première relation.} Soit $x > 0$. Posons $y = \ln x$. Par définition du logarithme (l'équivalence $y = \ln x \iff x = e^y$), on a $x = e^y$. En remplaçant $y$ par $\ln x$ dans cette égalité, on obtient $x = e^{\ln x}$. \hfill $\square$

\textbf{Seconde relation.} Soit $x \in \mathbb{R}$. Le nombre $e^x$ est strictement positif, donc $\ln(e^x)$ existe. Posons $z = \ln(e^x)$. Par définition, $z$ est l'unique réel tel que $e^z = e^x$. Or $x$ vérifie évidemment $e^x = e^x$ ; par \emph{unicité} (l'exponentielle est bijective), on conclut $z = x$, c'est-à-dire $\ln(e^x) = x$. \hfill $\square$

Retiens l'idée directrice : dans les deux cas, on \emph{nomme} la quantité inconnue, on applique la définition, puis on invoque l'\textbf{unicité} garantie par la bijection.
\end{experience}

\begin{exemplebox}[Premiers calculs avec $\ln$]
Résolvons l'équation $e^x = 5$. Comme $5 > 0$, l'équivalence fondamentale donne :
\[ e^x = 5 \iff x = \ln 5. \]
La calculatrice fournit $\ln 5 \approx 1{,}609$. Vérifions : $e^{1{,}609} \approx 5{,}00$. \checkmark

De même, sans calculatrice : $\ln(e^3) = 3$ par la seconde relation fondamentale, et $e^{\ln 7} = 7$ par la première.
\end{exemplebox}

\begin{exoresolu}[Le capital du commerçant de Douala]
Le commerçant veut savoir au bout de combien d'années $n$ son capital aura doublé, c'est-à-dire résoudre $(1{,}1)^n = 2$.
\tcblower
\textbf{Solution.} Astuce : écrivons $1{,}1$ comme une exponentielle. Posons $1{,}1 = e^{a}$ avec $a = \ln(1{,}1) \approx 0{,}0953$. Alors :
\[ (1{,}1)^n = 2 \iff e^{na} = 2 \iff na = \ln 2 \iff n = \frac{\ln 2}{\ln(1{,}1)} \approx \frac{0{,}693}{0{,}0953} \approx 7{,}27. \]
Comme les intérêts sont versés annuellement, le capital aura doublé au bout de \textbf{8 ans}. (Vérification : $(1{,}1)^8 \approx 2{,}14 > 2$, tandis que $(1{,}1)^7 \approx 1{,}95 < 2$.) Nous reverrons cette technique en détail dans la section sur les équations.
\end{exoresolu}

\courssub{Sens de variation et signe de $\ln x$}

La stricte croissance de $\ln$ n'est pas un détail : c'est l'outil qui permettra de résoudre les \emph{inéquations}. Elle signifie concrètement que $\ln$ « conserve l'ordre » :
\[ \text{Pour tous } a, b > 0 : \qquad a < b \iff \ln a < \ln b. \]

\begin{propriete}[Signe de $\ln x$]
\begin{itemize}
\item Si $0 < x < 1$, alors $\ln x < 0$ ;
\item si $x = 1$, alors $\ln x = 0$ ;
\item si $x > 1$, alors $\ln x > 0$.
\end{itemize}
\end{propriete}

\begin{experience}[Justification du signe de $\ln x$]
Nous savons que $\ln 1 = 0$ et que $\ln$ est \emph{strictement croissante} sur $]0\,;+\infty[$.

\textbf{Cas 1 :} soit $x > 1$. Par stricte croissance, $x > 1 \Rightarrow \ln x > \ln 1 = 0$.

\textbf{Cas 2 :} soit $0 < x < 1$. Par stricte croissance, $x < 1 \Rightarrow \ln x < \ln 1 = 0$.

\textbf{Cas 3 :} $x = 1$ donne $\ln x = 0$. \hfill $\square$

Tout repose donc sur deux faits : la valeur $\ln 1 = 0$ et la croissance. C'est une démonstration courte mais exemplaire de la méthode à rédiger au Bac.
\end{experience}

\begin{exemplebox}[Lecture de signes]
Sans calculatrice : $\ln 3 > 0$ car $3 > 1$ ; $\ln(0{,}4) < 0$ car $0 < 0{,}4 < 1$ ; $\ln\left(\frac{1}{2}\right) < 0$ car $\frac{1}{2} < 1$. De plus, comme $2 < 3$, on a $\ln 2 < \ln 3$ : effectivement, $\ln 2 \approx 0{,}693$ et $\ln 3 \approx 1{,}099$.
\end{exemplebox}

\begin{attention}[Deux confusions classiques]
\begin{itemize}
\item $\ln x$ n'est \textbf{pas} toujours positif ! Il est négatif sur $]0\,;1[$. Dire « un logarithme est positif » est faux en général.
\item Ne confonds pas $\ln x < 0$ avec « $\ln x$ n'existe pas » : $\ln(0{,}5) \approx -0{,}693$ existe parfaitement, il est simplement négatif. Ce qui n'existe pas, c'est $\ln$ d'un nombre négatif ou nul.
\end{itemize}
\end{attention}

\courssub{Représentation graphique : la symétrie avec l'exponentielle}

Il existe un moyen géométrique élégant d'obtenir la courbe de $\ln$ sans aucun calcul supplémentaire. Rappelons un résultat général : lorsque deux fonctions $f$ et $f^{-1}$ sont réciproques, leurs courbes sont \cle{symétriques par rapport à la droite d'équation $y = x$} (la première bissectrice). En effet, le point $M(a\,;b)$ appartient à la courbe de $f$ si et seulement si $b = f(a)$, c'est-à-dire $a = f^{-1}(b)$, ce qui signifie que le point $M'(b\,;a)$ — le symétrique de $M$ par rapport à $y = x$ — appartient à la courbe de $f^{-1}$.

\begin{propriete}[Courbe représentative de $\ln$]
Dans un repère orthonormé, la courbe représentative de la fonction $\ln$ est la \cle{symétrique} de celle de la fonction $\exp$ par rapport à la droite d'équation $y = x$.
\end{propriete}

Observons les points remarquables : la courbe de $\exp$ passe par $(0\,;1)$ et $(1\,;e)$ ; par symétrie, celle de $\ln$ passe par $(1\,;0)$ et $(e\,;1)$. L'axe des ordonnées, asymptote horizontale « à gauche » pour $\exp$... devient, après échange des rôles de $x$ et $y$, une \emph{asymptote verticale} pour $\ln$ en $0$ : la courbe de $\ln$ « plonge » vers $-\infty$ lorsque $x$ s'approche de $0$ par valeurs positives. Nous démontrerons ces limites dans la section 4.

\begin{popfigure}
\begin{center}
\begin{tikzpicture}
\begin{axis}[
  width=0.86\linewidth,
  axis lines=middle,
  xlabel={$x$}, ylabel={$y$},
  xmin=-3.2, xmax=4.6, ymin=-3.2, ymax=4.6,
  xtick={-2,-1,1,2,3,4}, ytick={-2,-1,1,2,3,4},
  grid=both, grid style={popline!40, very thin},
  legend style={at={(0.03,0.97)}, anchor=north west, draw=popline, fill=popcream},
  samples=200, domain=-3:4.6
]
\addplot[popcurve, popblue, domain=-3:1.6] {exp(x)};
\addlegendentry{$y = e^x$}
\addplot[popcurve, poporange, domain=0.02:4.6] {ln(x)};
\addlegendentry{$y = \ln x$}
\addplot[dashed, popmuted, domain=-3:4.6] {x};
\addlegendentry{$y = x$}
% Points remarquables et symétrie
\addplot[only marks, mark=*, popdark, mark size=1.6pt] coordinates {(0,1) (1,0) (1,2.718) (2.718,1)};
\node[popdark, anchor=south west, font=\small] at (axis cs:0.05,1.05) {$(0\,;1)$};
\node[popdark, anchor=north west, font=\small] at (axis cs:1.05,-0.05) {$(1\,;0)$};
\node[popdark, anchor=south, font=\small] at (axis cs:1,2.718) {$(1\,;e)$};
\node[popdark, anchor=west, font=\small] at (axis cs:2.718,0.9) {$(e\,;1)$};
\draw[<->, poppink, thick] (axis cs:1.35,2.35) -- (axis cs:2.35,1.35);
\end{axis}
\end{tikzpicture}
\end{center}
\legende{Les courbes de $\exp$ et de $\ln$ sont symétriques par rapport à la première bissectrice $y = x$. Les points $(0\,;1)$ et $(1\,;0)$, ainsi que $(1\,;e)$ et $(e\,;1)$, se correspondent par cette symétrie.}
\end{popfigure}

\begin{methode}[Exploiter la symétrie pour lire la courbe de $\ln$]
\begin{enumerate}
\item Repère les points clés de $\exp$ : $(0\,;1)$, $(1\,;e)$.
\item Échange les coordonnées : tu obtiens $(1\,;0)$ et $(e\,;1)$ sur la courbe de $\ln$.
\item Traduis les comportements aux bornes : $\exp$ tend vers $0$ en $-\infty$ $\Rightarrow$ $\ln$ tend vers $-\infty$ en $0^+$ (asymptote verticale $x = 0$) ; $\exp$ tend vers $+\infty$ en $+\infty$ $\Rightarrow$ $\ln$ tend vers $+\infty$ en $+\infty$.
\item Vérifie la cohérence : la courbe de $\ln$ est strictement croissante, négative avant $1$, positive après $1$.
\end{enumerate}
\end{methode}

\begin{exoresolu}[Lecture graphique et équivalence fondamentale]
Sur la figure ci-dessus, on lit que la courbe de $\ln$ coupe l'axe des abscisses en un seul point. Déterminer ce point, puis résoudre graphiquement et algébriquement l'équation $\ln x = 2$.
\tcblower
\textbf{Solution.} La courbe de $\ln$ coupe l'axe des abscisses au point $(1\,;0)$, ce qui confirme $\ln 1 = 0$.

Pour $\ln x = 2$ : graphiquement, on cherche l'abscisse du point de la courbe d'ordonnée $2$ ; on lit $x \approx 7{,}4$. Algébriquement, l'équivalence fondamentale donne :
\[ \ln x = 2 \iff x = e^2 \approx 7{,}389. \]
La solution est $S = \{e^2\}$. Vérification : $\ln(e^2) = 2$ par la seconde relation fondamentale. \checkmark
\end{exoresolu}

\begin{aretenir}[L'essentiel de la section]
\begin{itemize}
\item $\ln$ est la \cle{réciproque} de $\exp$ ; elle est définie sur $]0\,;+\infty[$ uniquement.
\item Équivalence fondamentale : $y = \ln x \iff x = e^y$ (pour $x > 0$).
\item Relations fondamentales : $e^{\ln x} = x$ pour $x > 0$ et $\ln(e^x) = x$ pour tout réel $x$.
\item Valeurs de référence : $\ln 1 = 0$ et $\ln e = 1$.
\item $\ln$ est continue et strictement croissante sur $]0\,;+\infty[$ : elle conserve l'ordre.
\item Signe : $\ln x < 0$ si $0 < x < 1$ ; $\ln x > 0$ si $x > 1$.
\item Les courbes de $\ln$ et $\exp$ sont symétriques par rapport à la droite $y = x$.
\end{itemize}
\end{aretenir}

\begin{savaistu}[D'où vient le mot « logarithme » ?]
Le mot \emph{logarithme} vient du grec \emph{logos} (raison, rapport) et \emph{arithmos} (nombre) : littéralement, le « nombre du rapport ». Les logarithmes furent inventés au début du XVII\textsuperscript{e} siècle par le mathématicien écossais John \textbf{Neper} (1550--1617) — d'où l'adjectif \emph{népérien} — pour transformer les multiplications fastidieuses en simples additions, grâce à la propriété $\ln(ab) = \ln a + \ln b$ que tu découvriras dans la prochaine section. Les astronomes de l'époque, écrasés par les calculs, y gagnèrent des années de vie : Laplace dira que Neper « a doublé la vie des astronomes ». Aujourd'hui encore, les logarithmes sont partout : l'échelle des décibels pour le son (comme à la SONARA !), l'échelle de Richter pour les séismes, le pH en chimie, et les calculs d'intérêts composés de notre commerçant de Douala.
\end{savaistu}

\coursec{Propriétés algébriques de ln}

Dans la section précédente, nous avons défini le logarithme népérien comme la \emph{fonction réciproque} de l'exponentielle : pour tout réel $a > 0$, $\ln a$ est l'unique réel $x$ tel que $e^x = a$. Nous avons aussi retenu les identités fondamentales $\ln 1 = 0$, $\ln e = 1$ et $e^{\ln a} = a$. Cette relation de réciprocité va maintenant porter ses fruits : l'exponentielle transforme les \emph{additions} en \emph{multiplications} (c'est la célèbre formule $e^{x+y} = e^x \times e^y$). Par réciprocité, le logarithme va faire le chemin inverse : il transforme les \emph{multiplications} en \emph{additions}. C'est cette idée simple, mais d'une puissance extraordinaire, qui fait toute la richesse des propriétés algébriques de $\ln$ que nous allons établir dans cette section. Ces formules sont les outils de base de la résolution d'équations et d'inéquations au Baccalauréat : il faut les connaître par cœur, avec leurs conditions de validité.

\courssub{La propriété fondamentale : le logarithme d'un produit}

Commençons par une intuition. Prenons deux nombres strictement positifs, par exemple $a = 8$ et $b = 4$. Leur produit est $ab = 32$. Que vaut $\ln 32$ par rapport à $\ln 8$ et $\ln 4$ ? Avec une calculatrice : $\ln 8 \approx 2{,}079$, $\ln 4 \approx 1{,}386$, et $\ln 32 \approx 3{,}466$. Or $2{,}079 + 1{,}386 = 3{,}465 \approx 3{,}466$ (à la précision des arrondis près). Il semble donc que $\ln 32 = \ln 8 + \ln 4$. Ce n'est pas un hasard : c'est la propriété fondamentale du logarithme, appelée \cle{relation fonctionnelle}.

\begin{propriete}[Relation fonctionnelle du logarithme népérien (démonstration exigible)]
Pour tous réels $a > 0$ et $b > 0$ :
\[ \ln(ab) = \ln a + \ln b. \]
\end{propriete}

\begin{experience}[Démonstration de la relation fonctionnelle]
Soient $a > 0$ et $b > 0$. L'idée est de passer par l'exponentielle, qui est la réciproque de $\ln$.
\begin{enumerate}
\item \textbf{Étape 1 : écrire $a$ et $b$ sous forme exponentielle.} Puisque $a > 0$, posons $x = \ln a$, de sorte que $a = e^x$. De même, posons $y = \ln b$, de sorte que $b = e^y$.
\item \textbf{Étape 2 : calculer le produit.} On a alors
\[ ab = e^x \times e^y = e^{x+y}, \]
d'après la propriété fondamentale de l'exponentielle.
\item \textbf{Étape 3 : appliquer $\ln$ aux deux membres.} Comme $ab > 0$, on peut prendre le logarithme :
\[ \ln(ab) = \ln\left(e^{x+y}\right) = x + y, \]
car $\ln(e^t) = t$ pour tout réel $t$.
\item \textbf{Étape 4 : conclure.} En remplaçant $x$ et $y$ par leurs valeurs :
\[ \ln(ab) = \ln a + \ln b. \qquad \blacksquare \]
\end{enumerate}
\end{experience}

\begin{aretenir}[L'idée à retenir]
Le logarithme népérien \emph{transforme les produits en sommes}. C'est exactement le comportement inverse de l'exponentielle, qui transforme les sommes en produits. Retenir cette phrase, c'est retenir toutes les formules de cette section.
\end{aretenir}

\begin{exemplebox}[Première application]
Sachant que $\ln 2 \approx 0{,}693$ et $\ln 3 \approx 1{,}099$, calculons $\ln 6$ sans calculatrice :
\[ \ln 6 = \ln(2 \times 3) = \ln 2 + \ln 3 \approx 0{,}693 + 1{,}099 = 1{,}792. \]
De même, $\ln 12 = \ln(4 \times 3) = \ln 4 + \ln 3 = 2\ln 2 + \ln 3 \approx 2 \times 0{,}693 + 1{,}099 = 2{,}485$.
\end{exemplebox}

\courssub{Logarithme de l'inverse et logarithme d'un quotient}

\begin{propriete}[Logarithme de l'inverse]
Pour tout réel $b > 0$ :
\[ \ln\left(\frac{1}{b}\right) = -\ln b. \]
\end{propriete}

\begin{experience}[Démonstration guidée]
Partons d'une égalité évidente : $b \times \dfrac{1}{b} = 1$. Appliquons $\ln$ aux deux membres (les deux membres sont strictement positifs) :
\[ \ln\left(b \times \frac{1}{b}\right) = \ln 1. \]
Le membre de gauche vaut $\ln b + \ln\left(\dfrac{1}{b}\right)$ par la relation fonctionnelle, et le membre de droite vaut $0$. Donc :
\[ \ln b + \ln\left(\frac{1}{b}\right) = 0 \quad \Longrightarrow \quad \ln\left(\frac{1}{b}\right) = -\ln b. \qquad \blacksquare \]
\end{experience}

\begin{propriete}[Logarithme d'un quotient]
Pour tous réels $a > 0$ et $b > 0$ :
\[ \ln\left(\frac{a}{b}\right) = \ln a - \ln b. \]
\end{propriete}

La démonstration est immédiate en combinant les deux propriétés précédentes : un quotient n'est rien d'autre qu'un produit par l'inverse. Elle est proposée en exercice corrigé ci-dessous.

\begin{exoresolu}[Démontrer la formule du quotient]
\textbf{Énoncé.} Soient $a > 0$ et $b > 0$. En écrivant $\dfrac{a}{b} = a \times \dfrac{1}{b}$, démontrer que $\ln\left(\dfrac{a}{b}\right) = \ln a - \ln b$.
\tcblower
\textbf{Solution.} On écrit le quotient comme un produit :
\[ \ln\left(\frac{a}{b}\right) = \ln\left(a \times \frac{1}{b}\right). \]
On applique la relation fonctionnelle (licite car $a > 0$ et $\dfrac{1}{b} > 0$) :
\[ \ln\left(a \times \frac{1}{b}\right) = \ln a + \ln\left(\frac{1}{b}\right). \]
Or $\ln\left(\dfrac{1}{b}\right) = -\ln b$, d'où :
\[ \ln\left(\frac{a}{b}\right) = \ln a - \ln b. \qquad \blacksquare \]
\end{exoresolu}

\begin{exemplebox}[Application numérique]
Calculons $\ln\left(\dfrac{1}{2}\right)$ et $\ln\left(\dfrac{3}{2}\right)$ sachant que $\ln 2 \approx 0{,}693$ et $\ln 3 \approx 1{,}099$ :
\[ \ln\left(\frac{1}{2}\right) = -\ln 2 \approx -0{,}693 \quad ; \quad \ln\left(\frac{3}{2}\right) = \ln 3 - \ln 2 \approx 1{,}099 - 0{,}693 = 0{,}406. \]
\end{exemplebox}

\courssub{Logarithme d'une puissance et d'une racine carrée}

\begin{propriete}[Logarithme d'une puissance]
Pour tout réel $a > 0$ et tout entier $n \in \mathbb{Z}$ :
\[ \ln\left(a^n\right) = n \ln a. \]
Cette formule s'étend à tout exposant rationnel : pour $r \in \mathbb{Q}$, $\ln(a^r) = r \ln a$.
\end{propriete}

\begin{experience}[Démonstration guidée pour un entier naturel]
\textbf{Cas $n \in \mathbb{N}$.} On raisonne par récurrence sur $n$.
\begin{enumerate}
\item \textbf{Initialisation.} Pour $n = 0$ : $\ln(a^0) = \ln 1 = 0 = 0 \times \ln a$. La propriété est vraie au rang $0$.
\item \textbf{Hérédité.} Supposons $\ln(a^n) = n\ln a$ pour un entier naturel $n$ fixé (hypothèse de récurrence). Alors :
\[ \ln\left(a^{n+1}\right) = \ln\left(a^n \times a\right) = \ln\left(a^n\right) + \ln a = n\ln a + \ln a = (n+1)\ln a. \]
La propriété est donc vraie au rang $n+1$.
\item \textbf{Conclusion.} Par récurrence, $\ln(a^n) = n\ln a$ pour tout $n \in \mathbb{N}$.
\item \textbf{Cas des entiers négatifs.} Si $n < 0$, alors $-n > 0$ et $a^n = \dfrac{1}{a^{-n}}$, donc en appliquant la formule de l'inverse puis le cas positif :
\[ \ln(a^n) = -\ln\left(a^{-n}\right) = -(-n)\ln a = n\ln a. \qquad \blacksquare \]
\end{enumerate}
\textbf{Extension aux rationnels.} Si $r = \dfrac{p}{q}$ avec $p \in \mathbb{Z}$ et $q \in \mathbb{N}^*$, alors $\left(a^{p/q}\right)^q = a^p$, donc en prenant le logarithme : $q\ln\left(a^{p/q}\right) = p\ln a$, c'est-à-dire $\ln(a^{p/q}) = \dfrac{p}{q}\ln a$.
\end{experience}

\begin{propriete}[Logarithme d'une racine carrée]
Pour tout réel $a > 0$ :
\[ \ln\left(\sqrt{a}\right) = \frac{1}{2}\ln a. \]
\end{propriete}

En effet, $\sqrt{a} = a^{1/2}$, et la formule des puissances avec $r = \dfrac{1}{2}$ donne directement le résultat. On peut aussi le démontrer sans les exposants fractionnaires : comme $\sqrt{a} \times \sqrt{a} = a$, la relation fonctionnelle donne $\ln(\sqrt{a}) + \ln(\sqrt{a}) = \ln a$, donc $2\ln(\sqrt{a}) = \ln a$, d'où le résultat en divisant par $2$.

\begin{exemplebox}[Applications numériques]
\begin{itemize}
\item $\ln 8 = \ln(2^3) = 3\ln 2 \approx 3 \times 0{,}693 = 2{,}079$ ;
\item $\ln\left(\dfrac{1}{16}\right) = \ln(2^{-4}) = -4\ln 2 \approx -2{,}773$ ;
\item $\ln\sqrt{5} = \dfrac{1}{2}\ln 5$ ; en particulier $\ln\sqrt{e} = \dfrac{1}{2}\ln e = \dfrac{1}{2}$.
\end{itemize}
\end{exemplebox}

\courssub{Tableau récapitulatif des formules}

\begin{aretenir}[Formulaire algébrique du logarithme népérien]
\begin{center}
\begin{tabularx}{\linewidth}{|l|X|l|}
\hline
\rowcolor{popblueL} \textbf{Formule} & \textbf{Énoncé} & \textbf{Conditions} \\
\hline
Produit & $\ln(ab) = \ln a + \ln b$ & $a > 0$, $b > 0$ \\
\hline
Inverse & $\ln\left(\dfrac{1}{b}\right) = -\ln b$ & $b > 0$ \\
\hline
Quotient & $\ln\left(\dfrac{a}{b}\right) = \ln a - \ln b$ & $a > 0$, $b > 0$ \\
\hline
Puissance entière & $\ln(a^n) = n\ln a$ & $a > 0$, $n \in \mathbb{Z}$ \\
\hline
Puissance rationnelle & $\ln(a^r) = r\ln a$ & $a > 0$, $r \in \mathbb{Q}$ \\
\hline
Racine carrée & $\ln(\sqrt{a}) = \dfrac{1}{2}\ln a$ & $a > 0$ \\
\hline
Valeurs remarquables & $\ln 1 = 0$ ; $\ln e = 1$ & --- \\
\hline
\end{tabularx}
\end{center}
\end{aretenir}

\begin{attention}[Le piège classique : il n'existe AUCUNE formule pour $\ln(a+b)$]
Les formules concernent le produit, le quotient et les puissances — \textbf{jamais la somme}. En général :
\[ \ln(a+b) \neq \ln a + \ln b \quad \text{et} \quad \ln(a-b) \neq \ln a - \ln b. \]
Contre-exemple : $\ln(1+1) = \ln 2 \approx 0{,}693$, alors que $\ln 1 + \ln 1 = 0 + 0 = 0$. De même, attention à ne pas confondre les opérations : $\ln a \times \ln b \neq \ln(ab)$ et $\dfrac{\ln a}{\ln b} \neq \ln\left(\dfrac{a}{b}\right)$. Enfin, n'oublie jamais la condition d'existence : $\ln x$ n'est défini que pour $x > 0$, donc avant d'appliquer une formule, vérifie que toutes les quantités sous les logarithmes sont strictement positives.
\end{attention}

\courssub{Simplifier et regrouper des expressions logarithmiques}

\begin{methode}[Regrouper plusieurs logarithmes en un seul]
Pour simplifier une somme ou une différence de logarithmes :
\begin{enumerate}
\item \textbf{Traiter d'abord les coefficients} : utiliser $k\ln a = \ln(a^k)$ pour faire « rentrer » chaque coefficient dans le logarithme comme exposant.
\item \textbf{Regrouper} : les sommes deviennent des produits ($\ln a + \ln b = \ln(ab)$), les différences deviennent des quotients ($\ln a - \ln b = \ln(a/b)$).
\item \textbf{Simplifier l'argument} obtenu (fraction à réduire, puissance à factoriser).
\item \textbf{Ré-exprimer} éventuellement le résultat en fonction d'un logarithme de référence (souvent $\ln 2$ ou $\ln 3$).
\end{enumerate}
\end{methode}

\begin{exoresolu}[Deux simplifications types du Baccalauréat]
\textbf{Énoncé.} 1) Montrer que $\ln 8 + \ln 4 - \ln 2 = \ln 16 = 4\ln 2$. \quad 2) Simplifier $A = 2\ln 3 + \ln 5 - \ln 15$.
\tcblower
\textbf{Solution.}
\begin{enumerate}
\item On regroupe : $\ln 8 + \ln 4 = \ln(8 \times 4) = \ln 32$, puis $\ln 32 - \ln 2 = \ln\left(\dfrac{32}{2}\right) = \ln 16$. Or $16 = 2^4$, donc $\ln 16 = \ln(2^4) = 4\ln 2$. Ainsi :
\[ \ln 8 + \ln 4 - \ln 2 = \ln 16 = 4\ln 2. \]
\item On fait d'abord rentrer le coefficient : $2\ln 3 = \ln(3^2) = \ln 9$. Puis :
\begin{align*}
A &= \ln 9 + \ln 5 - \ln 15 \\
&= \ln(9 \times 5) - \ln 15 \\
&= \ln 45 - \ln 15 = \ln\left(\frac{45}{15}\right) = \ln 3.
\end{align*}
Donc $A = \ln 3$. Remarque : on pouvait aussi remarquer dès le départ que $15 = 3 \times 5$, d'où $\ln 15 = \ln 3 + \ln 5$, et $A = 2\ln 3 + \ln 5 - \ln 3 - \ln 5 = \ln 3$.
\end{enumerate}
\end{exoresolu}

\begin{exoresolu}[Écrire en fonction de $\ln 2$ et $\ln 3$]
\textbf{Énoncé.} Exprimer $B = \ln 72 - \ln\sqrt{3} + \ln\left(\dfrac{1}{4}\right)$ en fonction de $\ln 2$ et $\ln 3$.
\tcblower
\textbf{Solution.} On décompose chaque terme. Comme $72 = 8 \times 9 = 2^3 \times 3^2$ :
\[ \ln 72 = \ln(2^3 \times 3^2) = 3\ln 2 + 2\ln 3. \]
Ensuite, $\ln\sqrt{3} = \dfrac{1}{2}\ln 3$ et $\ln\left(\dfrac{1}{4}\right) = -\ln 4 = -\ln(2^2) = -2\ln 2$. Donc :
\begin{align*}
B &= 3\ln 2 + 2\ln 3 - \frac{1}{2}\ln 3 - 2\ln 2 \\
&= (3-2)\ln 2 + \left(2 - \frac{1}{2}\right)\ln 3 \\
&= \ln 2 + \frac{3}{2}\ln 3.
\end{align*}
\end{exoresolu}

\courssub{Lien avec les suites géométriques}

Voici une application élégante et très appréciée au Baccalauréat. Considérons une suite géométrique $(u_n)$ de premier terme $u_0 > 0$ et de raison $q > 0$ : tous ses termes sont strictement positifs, donc on peut prendre leurs logarithmes. Posons $v_n = \ln(u_n)$. Alors, pour tout entier naturel $n$ :
\[ v_{n+1} = \ln(u_{n+1}) = \ln(q \times u_n) = \ln q + \ln(u_n) = v_n + \ln q. \]

\begin{propriete}[Logarithme d'une suite géométrique]
Si $(u_n)$ est une suite géométrique de raison $q > 0$ et de premier terme $u_0 > 0$, alors la suite $(v_n)$ définie par $v_n = \ln(u_n)$ est une suite \cle{arithmétique} de raison $r = \ln q$ et de premier terme $v_0 = \ln(u_0)$. En particulier :
\[ \ln(u_n) = \ln(u_0) + n\ln q. \]
\end{propriete}

Autrement dit, le logarithme transforme les progressions géométriques (croissance multiplicative) en progressions arithmétiques (croissance additive). C'est exactement le même phénomène que la relation fonctionnelle, vu cette fois à travers les suites.

\begin{exemplebox}[Application : épargne à intérêts composés]
Mama Fotso place un capital qui augmente de $5\,\%$ par an. Le capital (en milliers de FCFA) suit une suite géométrique de raison $q = 1{,}05$. La suite $v_n = \ln(u_n)$ est alors arithmétique de raison $\ln(1{,}05) \approx 0{,}0488$ : chaque année, le \emph{logarithme} du capital augmente de la même quantité constante. C'est pourquoi, pour savoir en combien d'années le capital double, on résout $1{,}05^n = 2$, c'est-à-dire $n\ln(1{,}05) = \ln 2$, soit $n = \dfrac{\ln 2}{\ln(1{,}05)} \approx \dfrac{0{,}693}{0{,}0488} \approx 14{,}2$ : il faut donc $15$ ans (on arrondit à l'année supérieure, car le doublement n'est atteint qu'en fin d'année).
\end{exemplebox}

\begin{exercice}[Pour t'entraîner]
\begin{enumerate}
\item Simplifier $C = 3\ln 2 + 2\ln 3 - \ln 12$ et $D = \ln\left(\dfrac{1}{9}\right) + \ln\sqrt{27}$.
\item Soit $(u_n)$ la suite géométrique de premier terme $u_0 = 3$ et de raison $q = 2$. Montrer que la suite $(v_n)$ définie par $v_n = \ln(u_n)$ est arithmétique, préciser sa raison, puis déterminer le plus petit entier $n$ tel que $u_n > 1000$.
\end{enumerate}
\end{exercice}

\begin{corrige}[Exercice]
\begin{enumerate}
\item $C = \ln(2^3) + \ln(3^2) - \ln 12 = \ln 8 + \ln 9 - \ln 12 = \ln\left(\dfrac{8 \times 9}{12}\right) = \ln\left(\dfrac{72}{12}\right) = \ln 6$. Pour $D$ : $\ln\left(\dfrac{1}{9}\right) = -\ln 9 = -2\ln 3$ et $\ln\sqrt{27} = \ln(27^{1/2}) = \ln(3^{3/2}) = \dfrac{3}{2}\ln 3$, donc $D = -2\ln 3 + \dfrac{3}{2}\ln 3 = -\dfrac{1}{2}\ln 3$.
\item Pour tout entier naturel $n$ : $v_{n+1} = \ln(u_{n+1}) = \ln(2u_n) = \ln 2 + \ln(u_n) = v_n + \ln 2$. Donc $(v_n)$ est arithmétique de raison $r = \ln 2$ et de premier terme $v_0 = \ln 3$. Ensuite, $u_n = 3 \times 2^n$, et comme la fonction $\ln$ est strictement croissante :
\[ u_n > 1000 \iff 3 \times 2^n > 1000 \iff n\ln 2 > \ln\left(\frac{1000}{3}\right) \iff n > \frac{\ln(1000/3)}{\ln 2} \approx \frac{5{,}808}{0{,}693} \approx 8{,}38. \]
Le plus petit entier convenable est donc $n = 9$. Vérification : $u_9 = 3 \times 2^9 = 3 \times 512 = 1536 > 1000$, tandis que $u_8 = 3 \times 256 = 768 < 1000$.
\end{enumerate}
\end{corrige}

\begin{savaistu}[Les tables de logarithmes : la calculatrice d'avant la calculatrice]
Avant l'arrivée des calculatrices électroniques dans les années 1970, comment les ingénieurs, les astronomes ou les comptables effectuaient-ils de longues multiplications ? Grâce aux \emph{tables de logarithmes}, publiées dès le XVII\textsuperscript{e} siècle (Neper, Briggs). Pour calculer un produit $a \times b$, on lisait $\ln a$ et $\ln b$ dans la table, on les \emph{additionnait} — opération facile —, puis on cherchait dans la table le nombre dont le logarithme était cette somme : c'était le produit cherché. Tout le travail reposait sur la seule formule $\ln(ab) = \ln a + \ln b$. La règle à calcul, utilisée par les générations d'ingénieurs qui ont construit ponts et barrages (y compris au Cameroun), est une application mécanique directe de cette même propriété. Laplace disait que les logarithmes, en abrégeant les calculs, « doublaient la vie des astronomes ».
\end{savaistu}

En résumé, tu disposes maintenant de tout l'arsenal algébrique du logarithme : produit, inverse, quotient, puissances et racines. Ces formules vont être tes outils de travail dans la section suivante, où nous les mettrons en action pour résoudre des équations, des inéquations et des systèmes faisant intervenir $\ln$.

\coursec{Équations, inéquations et inconnue auxiliaire}

Dans la section précédente, nous avons établi les propriétés algébriques du logarithme népérien : $\ln(ab) = \ln a + \ln b$, $\ln\!\left(\dfrac{a}{b}\right) = \ln a - \ln b$, $\ln(a^n) = n\ln a$. Ces formules ne sont pas de simples curiosités : ce sont de véritables \cle{outils de résolution}. Leur grand pouvoir est de transformer des produits en sommes et des puissances en produits, c'est-à-dire de \emph{simplifier la structure} des équations. Cette section t'apprend à exploiter ces propriétés, combinées à deux faits essentiels — la fonction $\ln$ est \cle{définie uniquement} sur $]0\,;\,+\infty[$ et elle est \cle{strictement croissante} — pour résoudre équations, inéquations et systèmes. C'est l'un des savoir-faire les plus fréquemment évalués au Baccalauréat.

\courssub{Le principe fondamental : injectivité et croissance de ln}

Avant toute formule, comprenons l'intuition. La fonction $\ln$ est strictement croissante sur $]0\,;\,+\infty[$ : elle ne « repasse jamais deux fois par la même valeur ». Deux nombres \emph{distincts} ont donc toujours des logarithmes \emph{distincts}. Autrement dit, si deux logarithmes sont égaux, c'est que les nombres de départ étaient égaux. On dit que $\ln$ est \cle{injective}. De même, la croissance conserve l'ordre : comparer les logarithmes revient à comparer les nombres eux-mêmes.

\begin{propriete}[Équivalences fondamentales]
Pour tous réels $u > 0$ et $v > 0$, et pour tout réel $k$ :
\[ \ln u = \ln v \iff u = v \qquad ; \qquad \ln u = k \iff u = e^{k} \]
\[ \ln u < \ln v \iff 0 < u < v \qquad ; \qquad \ln u \geq k \iff u \geq e^{k} \quad (u > 0). \]
\end{propriete}

\begin{experience}[Justification de l'équivalence $\ln u < \ln v \iff 0 < u < v$]
Justifions proprement cette équivalence, qui est la clé de toutes les inéquations de cette section.
\begin{enumerate}
\item \textbf{Sens direct.} Supposons $\ln u < \ln v$. Comme $\ln$ n'est définie que sur $]0\,;\,+\infty[$, on a déjà $u > 0$ et $v > 0$. Supposons par l'absurde que $u \geq v$. Alors, par croissance de $\ln$, on aurait $\ln u \geq \ln v$, ce qui contredit l'hypothèse. Donc $u < v$.
\item \textbf{Sens réciproque.} Supposons $0 < u < v$. La stricte croissance de $\ln$ sur $]0\,;\,+\infty[$ donne immédiatement $\ln u < \ln v$.
\item \textbf{Conclusion.} Les deux sens étant établis, l'équivalence est démontrée : $\ln u < \ln v \iff 0 < u < v$. Le même raisonnement vaut pour $\leq$, $>$ et $\geq$.
\end{enumerate}
\end{experience}

\begin{attention}[Le réflexe n°1 : l'ensemble de validité]
L'erreur la plus fréquente — et la plus sanctionnée au Bac — est de résoudre mécaniquement sans vérifier que les quantités sous les logarithmes sont \textbf{strictement positives}. Écrire $\ln x = \ln(-2)$ n'a \emph{aucun sens} : $-2$ n'est pas dans l'ensemble de définition de $\ln$. De même, résoudre $\ln(x-1) = \ln(2x-5)$ en écrivant simplement $x - 1 = 2x - 5$, donc $x = 4$, sans vérifier que $x - 1 > 0$ et $2x - 5 > 0$, peut faire accepter des \cle{solutions parasites}. Ici $x = 4$ convient ($3 > 0$ et $3 > 0$), mais si l'on avait trouvé $x = 1$, il aurait fallu le \textbf{rejeter}. Toujours : conditions d'abord, résolution ensuite, vérification à la fin.
\end{attention}

\courssub{Équations du type $\ln u(x) = \ln v(x)$ ou $\ln u(x) = k$}

\begin{methode}[Résoudre une équation faisant intervenir ln]
\begin{enumerate}
\item \textbf{Chercher l'ensemble de validité} $\mathcal{E}$ : on résout les conditions de stricte positivité de \emph{tous} les arguments des logarithmes présents.
\item \textbf{Se ramener}, à l'aide des propriétés algébriques, à une forme canonique : $\ln u = \ln v$ ou $\ln u = k$.
\item \textbf{Résoudre} l'équation simplifiée : $u = v$, ou $u = e^{k}$.
\item \textbf{Vérifier} que chaque solution trouvée appartient à $\mathcal{E}$, et ne conserver que celles-là.
\end{enumerate}
\end{methode}

\begin{exemplebox}[Équation $\ln u = \ln v$]
Résolvons dans $\mathbb{R}$ l'équation $\ln(2x+1) = \ln(x+4)$.
\begin{enumerate}
\item \textbf{Validité :} $2x+1 > 0 \iff x > -\dfrac{1}{2}$ et $x+4 > 0 \iff x > -4$. Donc $\mathcal{E} = \left]-\dfrac{1}{2}\,;\,+\infty\right[$.
\item \textbf{Équivalence :} sur $\mathcal{E}$, $\ln(2x+1) = \ln(x+4) \iff 2x+1 = x+4$.
\item \textbf{Résolution :} $x = 3$.
\item \textbf{Vérification :} $3 \in \mathcal{E}$. Donc $\mathcal{S} = \{3\}$.
\end{enumerate}
\end{exemplebox}

\begin{exemplebox}[Équation $\ln u = k$]
Résolvons $\ln(3x - 2) = 2$.
\begin{enumerate}
\item \textbf{Validité :} $3x - 2 > 0 \iff x > \dfrac{2}{3}$, donc $\mathcal{E} = \left]\dfrac{2}{3}\,;\,+\infty\right[$.
\item \textbf{Équivalence :} $3x - 2 = e^{2}$.
\item \textbf{Résolution :} $x = \dfrac{e^{2} + 2}{3}$.
\item \textbf{Vérification :} $e^{2} > 0$ donc $\dfrac{e^{2}+2}{3} > \dfrac{2}{3}$ : la solution est valide. $\mathcal{S} = \left\{\dfrac{e^{2}+2}{3}\right\}$.
\end{enumerate}
\end{exemplebox}

\courssub{Inéquations avec ln : la croissance fait tout le travail}

Pour les inéquations, le principe est identique, avec un bonus : comme $\ln$ est \emph{strictement croissante}, on peut « faire passer » l'inégalité à travers $\ln$ \textbf{sans changer le sens}. C'est une différence majeure avec les fonctions décroissantes (comme $x \mapsto \dfrac{1}{x}$), qui renversent les inégalités.

\begin{aretenir}[Inéquations types]
Sous réserve de validité (arguments strictement positifs) :
\[ \ln u(x) \leq \ln v(x) \iff 0 < u(x) \leq v(x) \qquad ; \qquad \ln u(x) \geq k \iff u(x) \geq e^{k}. \]
La condition $u(x) > 0$ fait \emph{partie intégrante} de l'équivalence : elle ne s'oublie jamais.
\end{aretenir}

\begin{exoresolu}[Inéquation $\ln u \leq k$]
Résoudre dans $\mathbb{R}$ l'inéquation $\ln(2x - 3) \leq 1$.
\tcblower
\textbf{Étape 1 — Ensemble de validité.} Il faut $2x - 3 > 0$, c'est-à-dire $x > \dfrac{3}{2}$. On travaille donc sur $\mathcal{E} = \left]\dfrac{3}{2}\,;\,+\infty\right[$.

\textbf{Étape 2 — Équivalence.} Écrivons $1 = \ln e$. L'inéquation devient $\ln(2x-3) \leq \ln e$. Par stricte croissance de $\ln$ :
\[ \ln(2x-3) \leq \ln e \iff 0 < 2x - 3 \leq e. \]

\textbf{Étape 3 — Résolution.} $2x - 3 \leq e \iff x \leq \dfrac{e + 3}{2}$.

\textbf{Étape 4 — Conclusion.} On intersecte avec $\mathcal{E}$ :
\[ \mathcal{S} = \left]\dfrac{3}{2}\,;\,\dfrac{e+3}{2}\right]. \]
Remarque : $\dfrac{e+3}{2} \approx \num{2,86}$, donc l'intervalle solution est bien non vide.
\end{exoresolu}

\begin{exoresolu}[Inéquation $\ln u \leq \ln v$]
Résoudre $\ln(x+1) \leq \ln(5 - x)$.
\tcblower
\textbf{Validité :} $x + 1 > 0$ et $5 - x > 0$, soit $\mathcal{E} = ]-1\,;\,5[$.

\textbf{Équivalence :} par croissance de $\ln$, $\ln(x+1) \leq \ln(5-x) \iff x + 1 \leq 5 - x$, c'est-à-dire $2x \leq 4$, soit $x \leq 2$.

\textbf{Conclusion :} en intersectant avec $\mathcal{E}$ : $\mathcal{S} = ]-1\,;\,2]$.
\end{exoresolu}

\courssub{Inconnue auxiliaire $X = \ln x$ : équations du second degré}

Certaines équations ne ressemblent à rien de connu... jusqu'à ce qu'on change de regard. Considérons $(\ln x)^2 - \ln x - 2 = 0$. Si l'on pose $X = \ln x$, l'équation devient $X^2 - X - 2 = 0$ : un simple \cle{trinôme du second degré} ! Cette technique d'\cle{inconnue auxiliaire} (ou changement de variable) est un grand classique du Baccalauréat.

\begin{methode}[Résoudre $a(\ln x)^2 + b\ln x + c = 0$]
\begin{enumerate}
\item \textbf{Condition :} $x > 0$ (pour que $\ln x$ existe).
\item \textbf{Poser} $X = \ln x$. L'équation devient $aX^2 + bX + c = 0$.
\item \textbf{Résoudre} le trinôme en $X$ (discriminant $\Delta = b^2 - 4ac$).
\item \textbf{Revenir à $x$} : chaque solution $X_0$ donne $x = e^{X_0}$. Comme $e^{X_0} > 0$ toujours, toutes les solutions en $X$ fournissent des solutions en $x$.
\end{enumerate}
\end{methode}

\begin{exoresolu}[Équation trinôme en $\ln x$]
Résoudre dans $\mathbb{R}$ l'équation $(\ln x)^2 - \ln x - 2 = 0$.
\tcblower
\textbf{Condition :} $x > 0$.

\textbf{Changement d'inconnue :} posons $X = \ln x$. L'équation s'écrit $X^2 - X - 2 = 0$.

\textbf{Résolution en $X$ :} $\Delta = (-1)^2 - 4 \times 1 \times (-2) = 1 + 8 = 9 > 0$. Deux solutions :
\[ X_1 = \dfrac{1 - 3}{2} = -1 \qquad \text{et} \qquad X_2 = \dfrac{1 + 3}{2} = 2. \]

\textbf{Retour à $x$ :} $\ln x = -1 \iff x = e^{-1} = \dfrac{1}{e}$ ; $\ln x = 2 \iff x = e^{2}$.

\textbf{Conclusion :} $\mathcal{S} = \left\{e^{-1}\,;\,e^{2}\right\}$, soit approximativement $\mathcal{S} \approx \{0,37\,;\,7,39\}$.
\end{exoresolu}

Pour les \textbf{inéquations du second degré en $\ln x$}, on ajoute une étape : l'étude du \cle{signe du trinôme} en $X$, puis le retour à $x$ en utilisant la croissance de la fonction exponentielle ($X \leq X_0 \iff x \leq e^{X_0}$).

\begin{exoresolu}[Inéquation trinôme en $\ln x$]
Résoudre $(\ln x)^2 - \ln x - 2 \leq 0$.
\tcblower
\textbf{Condition :} $x > 0$. Posons $X = \ln x$ : on étudie $X^2 - X - 2 \leq 0$.

\textbf{Signe du trinôme :} les racines sont $X_1 = -1$ et $X_2 = 2$ (calculées ci-dessus). Un trinôme est du signe de $a$ (ici $a = 1 > 0$) \emph{à l'extérieur} des racines, et du signe opposé \emph{entre} les racines. Donc :
\[ X^2 - X - 2 \leq 0 \iff -1 \leq X \leq 2. \]

\textbf{Retour à $x$ :} $-1 \leq \ln x \leq 2$. Par stricte croissance de $\exp$ :
\[ e^{-1} \leq x \leq e^{2}. \]

\textbf{Conclusion :} $\mathcal{S} = \left[e^{-1}\,;\,e^{2}\right]$. La condition $x > 0$ est automatiquement satisfaite puisque $e^{-1} > 0$.
\end{exoresolu}

\begin{attention}[Pièges du changement d'inconnue]
\begin{itemize}
\item Ne jamais conclure « $X = -1$ et $X = 2$ » : la question porte sur $x$, pas sur $X$ ! Le retour $x = e^{X}$ est \textbf{obligatoire}.
\item Ne pas écrire « $\ln x = -1$ donc pas de solution car un logarithme est positif » : c'est faux, $\ln x$ peut être négatif (c'est $x$ qui doit être positif, et $e^{-1} > 0$ convient parfaitement).
\item Pour les inéquations, attention au sens : $\ln x \leq 2 \iff x \leq e^{2}$ car $\exp$ est croissante. Le sens de l'inégalité est \emph{conservé}.
\end{itemize}
\end{attention}

\courssub{Systèmes à inconnues auxiliaires}

La même idée s'étend aux systèmes : lorsque les inconnues $x$ et $y$ n'apparaissent qu'à travers $\ln x$ et $\ln y$, on pose $X = \ln x$ et $Y = \ln y$, on résout un système linéaire classique en $(X, Y)$, puis on revient par $x = e^{X}$, $y = e^{Y}$. Souvent, les propriétés algébriques permettent même de se passer du changement de variable en faisant apparaître $xy$ et $\dfrac{x}{y}$.

\begin{exoresolu}[Système en logarithmes]
Résoudre le système $\begin{cases} \ln x + \ln y = \ln 6 \\ \ln x - \ln y = \ln 2 \end{cases}$
\tcblower
\textbf{Conditions :} $x > 0$ et $y > 0$.

\textbf{Transformation :} par les propriétés algébriques, le système équivaut à
\[ \begin{cases} \ln(xy) = \ln 6 \\ \ln\!\left(\dfrac{x}{y}\right) = \ln 2 \end{cases} \iff \begin{cases} xy = 6 \\ \dfrac{x}{y} = 2 \end{cases} \]

\textbf{Résolution :} de la deuxième équation, $x = 2y$. En reportant dans la première : $2y \times y = 6$, soit $y^2 = 3$. Comme $y > 0$, $y = \sqrt{3}$, puis $x = 2\sqrt{3}$.

\textbf{Vérification :} $x = 2\sqrt{3} > 0$, $y = \sqrt{3} > 0$ ; $xy = 2 \times 3 = 6$ et $\dfrac{x}{y} = 2$. Tout est conforme.

\textbf{Conclusion :} $\mathcal{S} = \left\{\left(2\sqrt{3}\,;\,\sqrt{3}\right)\right\}$.

\emph{Variante par inconnues auxiliaires :} en posant $X = \ln x$, $Y = \ln y$, on obtient $X + Y = \ln 6$ et $X - Y = \ln 2$, d'où par addition $2X = \ln 12$, $X = \ln\sqrt{12} = \ln(2\sqrt{3})$, et $Y = \ln\sqrt{3}$. On retrouve $x = 2\sqrt{3}$, $y = \sqrt{3}$.
\end{exoresolu}

\courssub{Organigramme de décision : quelle méthode pour quelle équation ?}

Face à une équation ou inéquation avec $\ln$, le bon réflexe est de \emph{classer} la situation avant de se lancer. L'arbre de décision suivant résume la stratégie complète de cette section.

\begin{popfigure}
\begin{center}
\begin{tikzpicture}[
  font=\small,
  decision/.style={diamond, draw=poporange, fill=poporangeL, text=popdark, aspect=2.2, inner sep=1pt, align=center},
  bloc/.style={rectangle, rounded corners=3pt, draw=popblue, fill=popblueL, text=popdark, align=center, inner sep=4pt},
  final/.style={rectangle, rounded corners=3pt, draw=popgreen, fill=popgreenL, text=popdark, align=center, inner sep=4pt},
  fleche/.style={-{Stealth[length=2.5mm]}, thick, popdark}
]
\node[decision] (q1) {Forme de l'équation ?};
\node[bloc, below left=1.1cm and 2.6cm of q1, align=center] (a) {$\ln u = \ln v$\\[2pt] Conditions $u>0$, $v>0$\\ puis $u = v$};
\node[bloc, below=1.1cm of q1, align=center] (b) {$\ln u = k$\\[2pt] Condition $u>0$\\ puis $u = e^{k}$};
\node[bloc, below right=1.1cm and 2.6cm of q1, align=center] (c) {$a(\ln x)^2+b\ln x+c=0$\\[2pt] Poser $X=\ln x$,\\ trinôme en $X$};
\node[final, below=1.2cm of b, align=center] (fin) {Retour à $x$ ($x=e^{X}$ si besoin)\\ puis \textbf{vérifier} les conditions de validité};
\draw[fleche] (q1) -- (a);
\draw[fleche] (q1) -- (b);
\draw[fleche] (q1) -- (c);
\draw[fleche] (a) |- (fin);
\draw[fleche] (b) -- (fin);
\draw[fleche] (c) |- (fin);
\end{tikzpicture}
\end{center}
\legende{Arbre de décision : identifier la forme, appliquer l'équivalence adaptée, puis toujours vérifier les conditions de validité.}
\end{popfigure}

\begin{savaistu}[Pourquoi tant d'équations se résolvent avec ln ?]
Le logarithme a été inventé par John Napier en 1614 précisément pour \emph{simplifier les calculs} : transformer les multiplications en additions. Quatre siècles plus tard, cette idée reste vivante : en sismologie (échelle de Richter), en acoustique (décibels), en finance (taux de croissance continus) et même dans l'étude de la propagation des épidémies, on « passe au logarithme » pour rendre les équations résolubles. Au Cameroun, les ingénieurs télécoms qui dimensionnent les réseaux MTN ou Orange utilisent des formules en logarithmes pour exprimer les puissances de signal en dBm. Maîtriser les techniques de cette section, c'est donc acquérir un outil professionnel autant qu'un savoir-faire d'examen.
\end{savaistu}

\begin{exercice}[Pour s'entraîner]
\begin{enumerate}
\item Résoudre $\ln(x^2 - 1) = \ln(2x + 1)$ (attention à l'ensemble de validité !).
\item Résoudre $(\ln x)^2 - 3\ln x + 2 = 0$.
\item Résoudre le système $\begin{cases} \ln x + \ln y = \ln 8 \\ \ln x - \ln y = 0 \end{cases}$
\end{enumerate}
\end{exercice}

\begin{corrige}[Exercice]
\textbf{1.} Validité : $x^2 - 1 > 0$ et $2x+1 > 0$, soit $\mathcal{E} = ]1\,;\,+\infty[$. Alors $x^2 - 1 = 2x + 1 \iff x^2 - 2x - 2 = 0$, $\Delta = 12$, $x = 1 \pm \sqrt{3}$. Seul $1 + \sqrt{3} \approx 2,73$ appartient à $\mathcal{E}$ ; $1 - \sqrt{3} < 0$ est rejeté. $\mathcal{S} = \{1 + \sqrt{3}\}$.

\textbf{2.} $X = \ln x$ : $X^2 - 3X + 2 = 0$, $\Delta = 1$, $X = 1$ ou $X = 2$. Donc $\mathcal{S} = \{e\,;\,e^2\}$.

\textbf{3.} $\ln x = \ln y$ donc $x = y$, puis $\ln(x^2) = \ln 8$ donc $x^2 = 8$, $x = 2\sqrt{2}$ (car $x > 0$). $\mathcal{S} = \{(2\sqrt{2}\,;\,2\sqrt{2})\}$.
\end{corrige}

\coursec{Limites usuelles et croissances comparées}

Dans la section précédente, tu as appris à résoudre des équations et des inéquations grâce aux propriétés algébriques du logarithme népérien. Mais pour \emph{étudier} une fonction faisant intervenir $\ln$ — dresser son tableau de variations, tracer sa courbe — il faut connaître son comportement aux bornes de son ensemble de définition. Que vaut $\ln x$ lorsque $x$ devient très grand ? Lorsque $x$ s'approche de $0$ ? Et comment comparer la croissance de $\ln x$ à celle de $x$ ? C'est tout l'objet de cette section, qui te fournira les outils indispensables pour lever les formes indéterminées au Baccalauréat.

\courssub{Limite de $\ln x$ en $+\infty$}

\textbf{Intuition.}\par
Prenons des valeurs de plus en plus grandes : $\ln 10 \approx 2{,}3$, $\ln 100 \approx 4{,}6$, $\ln 1000 \approx 6{,}9$, $\ln 10^6 \approx 13{,}8$. La fonction $\ln$ croît \emph{lentement} — il faut multiplier $x$ par $10$ pour ajouter seulement $\ln 10 \approx 2{,}3$ à son logarithme — mais elle ne s'arrête jamais : $\ln(10^n) = n\ln 10$ devient aussi grand qu'on veut. C'est un peu comme un marcheur très lent mais infatigable : il finit toujours par dépasser n'importe quel point du chemin.

\begin{propriete}[Limite de $\ln$ en $+\infty$ (démonstration exigible)]
\[
\lim_{x \to +\infty} \ln x = +\infty.
\]
\end{propriete}

\begin{experience}[Démonstration : $\ln x$ dépasse toute borne]
Rappelons la définition : dire que $\lim\limits_{x \to +\infty} \ln x = +\infty$ signifie que pour tout réel $A$ (aussi grand soit-il), il existe un rang au-delà duquel $\ln x > A$.

Soit donc $A$ un réel quelconque. La fonction $\ln$ est strictement croissante sur $]0\,;+\infty[$ et réalise une bijection de $]0\,;+\infty[$ sur $\mathbb{R}$. On a les équivalences :
\[
\ln x > A \iff x > e^{A}.
\]
Posons $B = e^{A}$. Alors, pour tout $x > B$, on a $\ln x > A$.

Conclusion : quel que soit le réel $A$, dès que $x > e^{A}$, on a $\ln x > A$. C'est exactement la définition de $\lim\limits_{x \to +\infty} \ln x = +\infty$. \hfill$\blacksquare$
\end{experience}

\begin{exemplebox}[Un seuil concret]
On cherche à partir de quelle valeur de $x$ on est sûr que $\ln x > 100$. D'après la démonstration, il suffit que $x > e^{100}$. Or $e^{100} \approx 2{,}7 \times 10^{43}$ : un nombre astronomique ! Cela confirme l'intuition : $\ln$ atteint $+\infty$, mais \emph{très} lentement.
\end{exemplebox}

\begin{attention}[Piège classique : « $\ln x$ finit par se stabiliser »]
Beaucoup d'élèves, voyant la courbe de $\ln$ s'aplatir, concluent que $\ln x$ tend vers une limite finie. C'est \textbf{faux} : la courbe monte de moins en moins vite, mais elle monte \emph{sans cesse} et dépasse toute borne. Une pente qui diminue n'empêche pas l'altitude d'augmenter indéfiniment. Retiens : \cle{croissance lente} ne signifie pas \cle{croissance bornée}.
\end{attention}

\courssub{Limite de $\ln x$ en $0^+$ et asymptote verticale}

\textbf{Intuition.}\par
Puisque $\ln 1 = 0$ et que $\ln x < 0$ pour $0 < x < 1$, que se passe-t-il lorsque $x$ se rapproche de $0$ ? Calculons : $\ln 0{,}1 \approx -2{,}3$, $\ln 0{,}01 \approx -4{,}6$, $\ln 10^{-6} \approx -13{,}8$. Les valeurs plongent vers le bas, de plus en plus profondément.

\begin{propriete}[Limite de $\ln$ en $0^+$ (démonstration exigible)]
\[
\lim_{\substack{x \to 0 \\ x > 0}} \ln x = -\infty.
\]
\end{propriete}

\begin{experience}[Démonstration par le changement de variable $X = \dfrac{1}{x}$]
Posons $X = \dfrac{1}{x}$. Lorsque $x$ tend vers $0$ par valeurs positives, $X$ tend vers $+\infty$. De plus, par la propriété algébrique du logarithme :
\[
\ln x = \ln\left(\frac{1}{X}\right) = -\ln X.
\]
Or $\lim\limits_{X \to +\infty} \ln X = +\infty$ (théorème précédent), donc par le théorème sur les limites et le changement de variable :
\[
\lim_{\substack{x \to 0 \\ x > 0}} \ln x = \lim_{X \to +\infty} (-\ln X) = -\infty. \qquad \blacksquare
\]
\end{experience}

\begin{aretenir}[Interprétation graphique : asymptote verticale]
Puisque $\lim\limits_{x \to 0^+} \ln x = -\infty$, la droite d'équation $x = 0$ (l'\cle{axe des ordonnées}) est \textbf{asymptote verticale} à la courbe représentative de la fonction $\ln$. La courbe « colle » à l'axe des ordonnées en plongeant vers $-\infty$, sans jamais le toucher.
\end{aretenir}

\begin{popfigure}
\begin{center}
\begin{tikzpicture}
\begin{axis}[
    width=0.85\linewidth, height=7cm,
    axis lines=middle,
    xlabel={$x$}, ylabel={$y$},
    xmin=-0.5, xmax=8, ymin=-4, ymax=3,
    xtick={1,2,3,4,5,6,7}, ytick={-3,-2,-1,1,2},
    legend style={at={(0.98,0.35)}, anchor=east, font=\small}
]
\addplot[popcurve, domain=0.05:8, samples=300] {ln(x)};
\addlegendentry{$y = \ln x$}
\addplot[popink, dashed, thick] coordinates {(0,-4) (0,3)};
\addlegendentry{asymptote $x = 0$}
\addplot[poporange, only marks, mark=*] coordinates {(1,0)};
\node[poporange, anchor=south west] at (axis cs:1,0) {\small $(1\,;0)$};
\end{axis}
\end{tikzpicture}
\end{center}
\legende{La courbe de $\ln$ : asymptote verticale $x = 0$ en rouge, passage par $(1\,;0)$, et croissance lente mais infinie vers $+\infty$.}
\end{popfigure}

\courssub{Croissances comparées : la puissance l'emporte sur le logarithme}

\textbf{Intuition.}\par
Nous avons vu que $\ln x$ et $x$ tendent tous les deux vers $+\infty$. Mais lequel « gagne » la course ? Comparons : pour $x = 10^6$, on a $\ln x \approx 13{,}8$ alors que $x = 1\,000\,000$. Le rapport $\dfrac{\ln x}{x} \approx 1{,}4 \times 10^{-5}$ est minuscule. La fonction puissance écrase littéralement le logarithme.

\begin{propriete}[Croissance comparée (admise)]
\[
\lim_{x \to +\infty} \frac{\ln x}{x} = 0.
\]
Plus généralement, pour tout entier $n \geq 1$ : $\lim\limits_{x \to +\infty} \dfrac{\ln x}{x^n} = 0$. On dit que \cle{toute puissance de $x$ l'emporte sur $\ln x$}.
\end{propriete}

\begin{experience}[Pourquoi c'est plausible : une démonstration guidée par encadrement]
Cette démonstration n'est pas exigible, mais elle est formatrice. Suivons les étapes.
\begin{enumerate}
\item On admet l'inégalité suivante, valable pour tout $t > 0$ : $\ln t \leq t - 1$ (la courbe de $\ln$ est sous sa tangente en $1$).
\item Appliquons-la à $t = \sqrt{x}$ : $\ln \sqrt{x} \leq \sqrt{x} - 1$, c'est-à-dire $\dfrac{1}{2}\ln x \leq \sqrt{x}$, donc $\ln x \leq 2\sqrt{x}$.
\item Pour $x \geq 1$, on a $\ln x \geq 0$, d'où l'encadrement :
\[
0 \leq \frac{\ln x}{x} \leq \frac{2\sqrt{x}}{x} = \frac{2}{\sqrt{x}}.
\]
\item Or $\lim\limits_{x \to +\infty} \dfrac{2}{\sqrt{x}} = 0$. Par le \textbf{théorème des gendarmes}, $\lim\limits_{x \to +\infty} \dfrac{\ln x}{x} = 0$. \hfill$\blacksquare$
\end{enumerate}
\end{experience}

\begin{propriete}[Conséquence : limite de $x\ln x$ en $0^+$]
\[
\lim_{\substack{x \to 0 \\ x > 0}} x \ln x = 0.
\]
\end{propriete}

\begin{experience}[Démonstration par le changement de variable $X = \dfrac{1}{x}$]
Posons $X = \dfrac{1}{x}$. Quand $x \to 0^+$, $X \to +\infty$, et :
\[
x \ln x = \frac{1}{X} \ln\left(\frac{1}{X}\right) = -\frac{\ln X}{X}.
\]
Par croissance comparée, $\lim\limits_{X \to +\infty} \dfrac{\ln X}{X} = 0$, donc $\lim\limits_{x \to 0^+} x \ln x = 0$. \hfill$\blacksquare$

Remarque : c'est une forme indéterminée du type « $0 \times (-\infty)$ » que le changement de variable a transformée en croissance comparée connue. Retiens cette \cle{technique du changement} $X = \dfrac{1}{x}$ : elle échange les problèmes en $0^+$ et en $+\infty$.
\end{experience}

\begin{popfigure}
\begin{center}
\begin{tikzpicture}
\begin{axis}[
    width=0.85\linewidth, height=7cm,
    axis lines=middle,
    xlabel={$x$}, ylabel={$y$},
    xmin=0, xmax=5, ymin=0, ymax=6,
    xtick={1,2,3,4,5}, ytick={1,2,3,4,5},
    legend style={at={(0.03,0.97)}, anchor=north west, font=\small}
]
\addplot[popcurve, domain=0.05:5, samples=200] {ln(x)};
\addlegendentry{$y = \ln x$}
\addplot[poporange, thick, domain=0:2.4, samples=50] {x};
\addlegendentry{$y = x$}
\addplot[popgreen, thick, domain=0:1.8, samples=100] {exp(x)};
\addlegendentry{$y = e^x$}
\end{axis}
\end{tikzpicture}
\end{center}
\legende{Hiérarchie des croissances : $e^x$ explose, $x$ croît modérément, $\ln x$ croît très lentement. En $+\infty$ : $\ln x \ll x \ll e^x$.}
\end{popfigure}

\courssub{Limite de référence : $\lim\limits_{h \to 0} \dfrac{\ln(1+h)}{h} = 1$}

\textbf{Intuition.}\par
Près de $x = 1$, la courbe de $\ln$ ressemble à sa tangente, de pente $\ln'(1) = 1$. Autrement dit, pour $h$ proche de $0$, $\ln(1+h) \approx h$ : le logarithme de $(1+h)$ se comporte comme $h$ lui-même.

\begin{propriete}[Limite de référence (démonstration exigible)]
\[
\lim_{h \to 0} \frac{\ln(1+h)}{h} = 1.
\]
\end{propriete}

\begin{experience}[Démonstration guidée : reconnaître un taux d'accroissement]
\begin{enumerate}
\item La fonction $\ln$ est dérivable en $1$, de nombre dérivé $\ln'(1) = \dfrac{1}{1} = 1$.
\item Par définition du nombre dérivé :
\[
\ln'(1) = \lim_{h \to 0} \frac{\ln(1+h) - \ln 1}{h}.
\]
\item Or $\ln 1 = 0$, donc :
\[
\lim_{h \to 0} \frac{\ln(1+h)}{h} = \ln'(1) = 1. \qquad \blacksquare
\]
\end{enumerate}
\textbf{Moralité} : cette limite n'est rien d'autre que le \cle{taux d'accroissement} de $\ln$ en $1$. Savoir la reconnaître sous cette forme est un réflexe de Bac.
\end{experience}

\begin{exemplebox}[Application directe]
Calculer $\lim\limits_{x \to 0} \dfrac{\ln(1+3x)}{x}$.

On fait apparaître la limite de référence en écrivant :
\[
\frac{\ln(1+3x)}{x} = 3 \times \frac{\ln(1+3x)}{3x}.
\]
Posons $h = 3x$ : quand $x \to 0$, $h \to 0$, et $\dfrac{\ln(1+h)}{h} \to 1$. Donc :
\[
\lim_{x \to 0} \frac{\ln(1+3x)}{x} = 3 \times 1 = 3.
\]
\end{exemplebox}

\courssub{Technique : lever les indéterminations par factorisation}

Face à une forme indéterminée mêlant $\ln x$ et des puissances de $x$ (types « $\infty - \infty$ » ou « $\dfrac{\infty}{\infty}$ »), la stratégie reine consiste à \cle{factoriser par le terme dominant}, c'est-à-dire la puissance de $x$, puis à faire apparaître une croissance comparée.

\begin{methode}[Lever une indétermination avec $\ln$]
\begin{enumerate}
\item \textbf{Identifier} la forme indéterminée (« $\infty - \infty$ », « $\dfrac{\infty}{\infty}$ », « $0 \times \infty$ »).
\item \textbf{Factoriser} par la puissance de $x$ dominante (c'est elle qui l'emporte sur $\ln x$).
\item \textbf{Faire apparaître} un quotient du type $\dfrac{\ln x}{x}$ (ou $\dfrac{\ln x}{x^n}$), qui tend vers $0$.
\item \textbf{Conclure} par les théorèmes sur les opérations sur les limites.
\item Si l'indétermination est en $0^+$, penser au changement de variable $X = \dfrac{1}{x}$.
\end{enumerate}
\end{methode}

\begin{exoresolu}[Deux limites classiques au Baccalauréat]
\textbf{1.} Déterminer $\lim\limits_{x \to +\infty} (x - \ln x)$.

\textbf{2.} Déterminer $\lim\limits_{x \to +\infty} \dfrac{\ln x}{\sqrt{x}}$.
\tcblower
\textbf{1.} Il s'agit d'une forme indéterminée « $+\infty - \infty$ ». Factorisons par le terme dominant $x$ :
\[
x - \ln x = x\left(1 - \frac{\ln x}{x}\right).
\]
Par croissance comparée, $\lim\limits_{x \to +\infty} \dfrac{\ln x}{x} = 0$, donc $\lim\limits_{x \to +\infty} \left(1 - \dfrac{\ln x}{x}\right) = 1$. Par produit :
\[
\lim_{x \to +\infty} x\left(1 - \frac{\ln x}{x}\right) = +\infty \times 1 = +\infty.
\]
Conclusion : $\lim\limits_{x \to +\infty} (x - \ln x) = +\infty$. La puissance $x$ « avale » le logarithme.

\textbf{2.} Forme indéterminée « $\dfrac{\infty}{\infty}$ ». Astuce : écrire $\sqrt{x} = x^{1/2}$ et faire apparaître la croissance comparée. Posons $X = \sqrt{x}$, de sorte que $x = X^2$ et $\ln x = 2\ln X$. Quand $x \to +\infty$, $X \to +\infty$ :
\[
\frac{\ln x}{\sqrt{x}} = \frac{2\ln X}{X} \xrightarrow[X \to +\infty]{} 2 \times 0 = 0.
\]
Conclusion : $\lim\limits_{x \to +\infty} \dfrac{\ln x}{\sqrt{x}} = 0$. Même la modeste puissance $\sqrt{x}$ l'emporte sur $\ln x$ !
\end{exoresolu}

\begin{attention}[Erreurs fréquentes à bannir]
\begin{itemize}
\item Écrire « $\lim\limits_{x \to +\infty} \ln x = +\infty$ donc $\lim\limits_{x \to +\infty}(x - \ln x) = +\infty - \infty = 0$ » : les opérations sur les infinis n'existent pas ! Il faut \textbf{factoriser} avant de conclure.
\item Croire que $\dfrac{\ln x}{x}$ tend vers $1$ : c'est $\dfrac{\ln(1+h)}{h}$ qui tend vers $1$ (en $0$), pas $\dfrac{\ln x}{x}$ (qui tend vers $0$ en $+\infty$). Ne confonds pas les deux limites de référence.
\item Oublier que la limite en $0$ de $\ln$ n'a de sens que \textbf{par valeurs supérieures} ($x > 0$), car $\ln$ n'est pas définie sur $\mathbb{R}_-^*$.
\end{itemize}
\end{attention}

\begin{aretenir}[Les quatre limites à connaître par cœur]
\begin{center}
\begin{tabularx}{\linewidth}{|X|X|}
\hline
\rowcolor{popblueL}
\textbf{Limite} & \textbf{Interprétation} \\
\hline
$\lim\limits_{x \to +\infty} \ln x = +\infty$ & croissance lente mais infinie \\
\hline
$\lim\limits_{x \to 0^+} \ln x = -\infty$ & asymptote verticale $x = 0$ \\
\hline
$\lim\limits_{x \to +\infty} \dfrac{\ln x}{x} = 0$ \quad et \quad $\lim\limits_{x \to 0^+} x\ln x = 0$ & toute puissance l'emporte sur $\ln$ \\
\hline
$\lim\limits_{h \to 0} \dfrac{\ln(1+h)}{h} = 1$ & nombre dérivé de $\ln$ en $1$ \\
\hline
\end{tabularx}
\end{center}
\end{aretenir}

\begin{savaistu}[Le logarithme, champion de la lenteur]
La fonction $\ln$ croît si lentement qu'elle est utilisée pour mesurer des phénomènes gigantesques : l'échelle de Richter (séismes), les décibels (son) et le pH (chimie) sont des échelles logarithmiques ! Un séisme de magnitude $7$ libère une énergie $10$ fois plus grande qu'un séisme de magnitude $6$, mais l'échelle n'augmente que d'une unité. Au Cameroun, lorsque les ingénieurs de CAMTEL ou des opérateurs télécoms mesurent l'affaiblissement d'un signal sur une ligne, ils raisonnent en décibels — donc en logarithmes. La lenteur de $\ln$ est précisément ce qui permet de compresser des ordres de grandeur immenses en nombres maniables.
\end{savaistu}

\begin{exercice}[À toi de jouer]
Déterminer les limites suivantes, en justifiant chaque étape :
\begin{enumerate}
\item $\lim\limits_{x \to +\infty} \dfrac{\ln x}{x^2}$ ;
\item $\lim\limits_{x \to +\infty} (2x - 3\ln x)$ ;
\item $\lim\limits_{x \to 0^+} x^2 \ln x$ ;
\item $\lim\limits_{x \to 0} \dfrac{\ln(1 - 2x)}{x}$.
\end{enumerate}
\end{exercice}

\begin{corrige}[Exercice : À toi de jouer]
\textbf{1.} $\dfrac{\ln x}{x^2} = \dfrac{1}{x} \times \dfrac{\ln x}{x}$. Or $\dfrac{\ln x}{x} \to 0$ et $\dfrac{1}{x} \to 0$, donc par produit la limite vaut $0$.

\textbf{2.} Forme « $\infty - \infty$ ». On factorise par $x$ : $2x - 3\ln x = x\left(2 - 3\dfrac{\ln x}{x}\right)$. Comme $\dfrac{\ln x}{x} \to 0$, la parenthèse tend vers $2$, et par produit la limite est $+\infty$.

\textbf{3.} Posons $X = \dfrac{1}{x}$ : $x^2 \ln x = \dfrac{1}{X^2} \times (-\ln X) = -\dfrac{\ln X}{X^2} \to 0$ quand $X \to +\infty$. Donc $\lim\limits_{x \to 0^+} x^2 \ln x = 0$.

\textbf{4.} On écrit $\dfrac{\ln(1-2x)}{x} = -2 \times \dfrac{\ln(1-2x)}{-2x}$. Avec $h = -2x \to 0$, le quotient tend vers $1$, donc la limite vaut $-2$.
\end{corrige}

Tu maîtrises désormais le comportement asymptotique du logarithme népérien et la technique de factorisation par le terme dominant. Ces résultats seront les piliers de la section suivante, où nous étudierons la dérivation et les primitives liées à $\ln$, avant de mener l'étude complète des fonctions du type $\ln(ax+b)$.

\coursec{Dérivation et primitives liées à ln}

Dans la section précédente, nous avons étudié le comportement de la fonction $\ln$ aux bornes de son ensemble de définition : $\ln x \to -\infty$ quand $x \to 0^+$ et $\ln x \to +\infty$ quand $x \to +\infty$, avec une croissance « lente » face aux puissances de $x$. Une question naturelle se pose alors : \emph{comment varie la fonction $\ln$ entre ces bornes, et à quelle vitesse ?} Pour répondre à cette question, l'outil roi est la dérivation. C'est toute l'ambition de cette section : calculer la dérivée de $\ln$, en déduire celle des fonctions composées du type $\ln(ax+b)$, puis « remonter » le calcul pour trouver des primitives des fractions rationnelles simples du type $\dfrac{a}{bx+c}$. Ces deux savoir-faire — dériver et primitiver — sont au cœur de l'épreuve de mathématiques du Baccalauréat.

\courssub{Dérivabilité et dérivée de la fonction ln}

\textbf{Intuition.} La courbe de la fonction $\ln$ est obtenue à partir de celle de l'exponentielle par symétrie par rapport à la droite d'équation $y = x$. La courbe de l'exponentielle est « lisse » (elle admet une tangente en chacun de ses points, sans angle ni rupture) ; sa symétrique l'est donc aussi. On admet alors, grâce au théorème de dérivabilité des fonctions réciproques, que $\ln$ est dérivable sur tout son ensemble de définition $]0\,;+\infty[$. Reste à calculer sa dérivée : c'est l'objet du théorème suivant, dont la démonstration est \cle{exigible} au Baccalauréat.

\begin{propriete}[Théorème — Dérivée de la fonction ln (démonstration exigible)]
La fonction $\ln$ est dérivable sur $]0\,;+\infty[$ et, pour tout réel $x > 0$ :
\[
(\ln x)' = \frac{1}{x}.
\]
\end{propriete}

\begin{experience}[Démonstration du théorème $(\ln x)' = \dfrac{1}{x}$]
L'idée est lumineuse : on ne connaît pas encore la dérivée de $\ln$, mais on connaît celle de l'exponentielle, et les deux fonctions sont liées par la relation fondamentale $\mathrm{e}^{\ln x} = x$. Dérivons les deux membres !
\begin{enumerate}
\item \textbf{Point de départ.} Pour tout $x > 0$, par définition de la fonction réciproque de l'exponentielle :
\[
\mathrm{e}^{\ln x} = x.
\]
\item \textbf{Dérivation membre à membre.} Le membre de gauche est une fonction composée : c'est $\exp \circ\, \ln$. Par la formule de dérivation des fonctions composées $(v \circ u)' = (v' \circ u) \times u'$, avec $u = \ln$ et $v = \exp$ (dont la dérivée est elle-même), on obtient :
\[
\left(\mathrm{e}^{\ln x}\right)' = \mathrm{e}^{\ln x} \times (\ln x)'.
\]
Le membre de droite a pour dérivée $(x)' = 1$.
\item \textbf{Égalité des dérivées.} Deux fonctions égales ont des dérivées égales, donc pour tout $x > 0$ :
\[
\mathrm{e}^{\ln x} \times (\ln x)' = 1.
\]
\item \textbf{Conclusion.} Comme $\mathrm{e}^{\ln x} = x$ et que $x > 0$ (donc $x \neq 0$), on peut diviser par $\mathrm{e}^{\ln x}$ :
\[
(\ln x)' = \frac{1}{\mathrm{e}^{\ln x}} = \frac{1}{x}. \qquad \blacksquare
\]
\end{enumerate}
\end{experience}

\begin{aretenir}[Ce qu'il faut retenir]
Pour tout $x > 0$ : $(\ln x)' = \dfrac{1}{x}$. La dérivée de $\ln$ est donc une fonction \emph{strictement positive} sur $]0\,;+\infty[$, ce qui redémontre au passage que $\ln$ est \cle{strictement croissante} sur $]0\,;+\infty[$ — résultat que nous avions admis dans la première section. De plus, $\dfrac{1}{x}$ décroît vers $0$ quand $x$ grandit : la courbe de $\ln$ monte de plus en plus « doucement », ce qui confirme la croissance lente observée dans la section précédente.
\end{aretenir}

\begin{popfigure}
\begin{center}
\begin{tikzpicture}
\begin{axis}[
  width=0.85\linewidth, height=6.5cm,
  axis lines=middle, xlabel={$x$}, ylabel={$y$},
  xmin=-0.5, xmax=6.5, ymin=-3, ymax=3,
  xtick={1,2,3,4,5,6}, ytick={-2,-1,1,2},
  legend style={at={(0.98,0.55)},anchor=east,font=\small},
  grid=both, grid style={popline!40},
]
\addplot[popcurve,domain=0.05:6.3,samples=200]{ln(x)};
\addlegendentry{$y=\ln x$}
\addplot[popcurve,poporange,domain=0.34:6.3,samples=200]{1/x};
\addlegendentry{$y=\dfrac{1}{x}$ (sa dérivée)}
\end{axis}
\end{tikzpicture}
\end{center}
\legende{La courbe de $\ln$ (bleu) et celle de sa dérivée $x \mapsto \dfrac{1}{x}$ (orange). La dérivée est toujours positive : $\ln$ est strictement croissante. Quand $x$ grandit, $\dfrac{1}{x}$ se rapproche de $0$ : la courbe de $\ln$ s'aplatit.}
\end{popfigure}

\courssub{Dérivée de \texorpdfstring{$\ln u$}{ln u} et cas particulier \texorpdfstring{$\ln(ax+b)$}{ln(ax+b)}}

\textbf{Intuition.} Au Baccalauréat, on dérive rarement $\ln x$ seul : on dérive des fonctions \emph{composées} comme $\ln(3x-2)$ ou $\ln(x^2+1)$. La règle de dérivation des fonctions composées fait tout le travail.

\begin{propriete}[Dérivée de $\ln u$]
Soit $u$ une fonction dérivable et \cle{strictement positive} sur un intervalle $I$. Alors la fonction $\ln u : x \mapsto \ln\big(u(x)\big)$ est dérivable sur $I$ et :
\[
(\ln u)' = \frac{u'}{u}.
\]
En particulier, pour $u(x) = ax + b$ (avec $a \neq 0$), sur tout intervalle où $ax + b > 0$ :
\[
\big(\ln(ax+b)\big)' = \frac{a}{ax+b}.
\]
\end{propriete}

\begin{experience}[Démonstration par la dérivation composée]
La fonction $\ln u$ est la composée $\ln \circ\, u$. La formule de dérivation des fonctions composées s'écrit :
\[
(\ln \circ\, u)'(x) = \ln'\big(u(x)\big) \times u'(x).
\]
Or $\ln'(t) = \dfrac{1}{t}$ pour tout $t > 0$. Comme $u(x) > 0$ sur $I$, on peut appliquer cette formule en $t = u(x)$ :
\[
(\ln u)'(x) = \frac{1}{u(x)} \times u'(x) = \frac{u'(x)}{u(x)}. \qquad \blacksquare
\]
Pour le cas particulier $u(x) = ax + b$, on a $u'(x) = a$, d'où $\big(\ln(ax+b)\big)' = \dfrac{a}{ax+b}$, valable sur tout intervalle contenu dans l'ensemble où $ax+b > 0$.
\end{experience}

\begin{exemplebox}[Dériver $f(x) = \ln(3x-2)$]
\textbf{Étape 1 — Ensemble de dérivabilité.} Il faut $3x - 2 > 0$, c'est-à-dire $x > \dfrac{2}{3}$. La fonction $f$ est donc définie et dérivable sur $I = \left]\dfrac{2}{3}\,;+\infty\right[$.

\textbf{Étape 2 — Calcul de la dérivée.} On pose $u(x) = 3x - 2$, donc $u'(x) = 3$. Par la formule $(\ln u)' = \dfrac{u'}{u}$ :
\[
f'(x) = \frac{3}{3x-2} \quad \text{pour tout } x > \frac{2}{3}.
\]
\textbf{Étape 3 — Conséquence.} Sur $I$, $3x - 2 > 0$ et $3 > 0$, donc $f'(x) > 0$ : $f$ est strictement croissante sur $\left]\dfrac{2}{3}\,;+\infty\right[$.
\end{exemplebox}

\begin{attention}[Les deux erreurs classiques au Bac]
\begin{itemize}
\item \textbf{Oublier le facteur $a$.} La dérivée de $\ln(3x-2)$ n'est \emph{pas} $\dfrac{1}{3x-2}$ : c'est $\dfrac{3}{3x-2}$. Le numérateur est la \emph{dérivée de ce qui est dans le logarithme}. Réflexe de survie : « dérivée de $\ln(\text{truc})$ égale $\dfrac{\text{dérivée du truc}}{\text{truc}}$ ».
\item \textbf{Oublier la condition d'existence.} Écrire $f'(x) = \dfrac{3}{3x-2}$ sans préciser que $x > \dfrac{2}{3}$ est une rédaction incomplète : la fonction $f$ n'existe tout simplement pas pour $x \leqslant \dfrac{2}{3}$. Toujours commencer par déterminer l'ensemble de définition.
\end{itemize}
\end{attention}

\courssub{Primitives liées au logarithme}

\textbf{Intuition.} Primitiver, c'est « remonter » la dérivation : chercher une fonction $F$ dont la dérivée est la fonction donnée $f$. Puisque $(\ln x)' = \dfrac{1}{x}$, la réponse à « quelle fonction a pour dérivée $\dfrac{1}{x}$ ? » est immédiate : c'est $\ln x$ ! Toute la richesse vient ensuite de la formule composée $(\ln u)' = \dfrac{u'}{u}$, lue \emph{à l'envers}.

\begin{propriete}[Primitives de $\dfrac{1}{x}$ et de $\dfrac{u'}{u}$]
\begin{itemize}
\item Une primitive de $x \mapsto \dfrac{1}{x}$ sur $]0\,;+\infty[$ est $x \mapsto \ln x$. L'ensemble des primitives est $x \mapsto \ln x + k$, $k \in \mathbb{R}$.
\item Plus généralement, si $u$ est dérivable et strictement positive sur un intervalle $I$, alors une primitive de $\dfrac{u'}{u}$ sur $I$ est $\ln u$.
\end{itemize}
\end{propriete}

\begin{propriete}[Primitive de $\dfrac{a}{bx+c}$ — cas du programme]
Soit $b \neq 0$. Sur tout intervalle où $bx + c > 0$, une primitive de
\[
f(x) = \frac{a}{bx+c}
\quad \text{est} \quad
F(x) = \frac{a}{b}\,\ln(bx+c).
\]
\end{propriete}

\begin{experience}[Pourquoi ça marche ?]
Vérifions en dérivant $F(x) = \dfrac{a}{b}\ln(bx+c)$. Par la formule de dérivation de $\ln u$ avec $u(x) = bx + c$, donc $u'(x) = b$ :
\[
F'(x) = \frac{a}{b} \times \frac{b}{bx+c} = \frac{a}{bx+c} = f(x). \qquad \blacksquare
\]
L'idée profonde : la dérivée du dénominateur $bx + c$ est la constante $b$. On « fabrique » donc cette dérivée au numérateur en multipliant et divisant par $b$, ce qui fait apparaître le facteur correctif $\dfrac{a}{b}$.
\end{experience}

\begin{methode}[Primitiver une fraction rationnelle simple]
Pour trouver une primitive de $f(x) = \dfrac{\alpha}{bx+c}$ sur un intervalle où $bx + c > 0$ :
\begin{enumerate}
\item \textbf{Identifier} la dérivée du dénominateur : $(bx+c)' = b$.
\item \textbf{Écrire le numérateur} sous la forme $k \times b$, c'est-à-dire $\alpha = \dfrac{\alpha}{b} \times b$. On a donc $f(x) = \dfrac{\alpha}{b} \times \dfrac{b}{bx+c} = \dfrac{\alpha}{b} \times \dfrac{u'(x)}{u(x)}$.
\item \textbf{Reconnaître} la forme $k \times \dfrac{u'}{u}$, dont une primitive est $k \ln u$ : $F(x) = \dfrac{\alpha}{b}\ln(bx+c)$.
\item \textbf{Vérifier} en redérivant $F$ : on doit retrouver $f$. Cette vérification prend dix secondes et élimine toute erreur de coefficient.
\end{enumerate}
\end{methode}

\begin{exoresolu}[Primitive de $f(x) = \dfrac{4}{2x+1}$]
Déterminer une primitive de $f(x) = \dfrac{4}{2x+1}$ sur l'intervalle $I = \left]-\dfrac{1}{2}\,;+\infty\right[$.
\tcblower
\textbf{Solution détaillée.}

Sur $I$, on a $2x + 1 > 0$, donc le logarithme de $2x+1$ est bien défini.

\textbf{Étape 1.} La dérivée du dénominateur est $(2x+1)' = 2$.

\textbf{Étape 2.} On fait apparaître cette dérivée au numérateur : $4 = 2 \times 2$, donc
\[
f(x) = \frac{4}{2x+1} = 2 \times \frac{2}{2x+1} = 2 \times \frac{u'(x)}{u(x)} \quad \text{avec } u(x) = 2x+1.
\]

\textbf{Étape 3.} Une primitive de $\dfrac{u'}{u}$ est $\ln u$, donc une primitive de $f$ est :
\[
F(x) = 2\ln(2x+1).
\]

\textbf{Étape 4 — Vérification.} $F'(x) = 2 \times \dfrac{2}{2x+1} = \dfrac{4}{2x+1} = f(x)$. \checkmark

L'ensemble des primitives de $f$ sur $I$ est donc $x \mapsto 2\ln(2x+1) + k$, $k \in \mathbb{R}$.
\end{exoresolu}

\begin{attention}[Conditions d'existence : le piège silencieux]
La primitive $F(x) = \dfrac{a}{b}\ln(bx+c)$ n'existe que là où $bx + c > 0$. Par exemple, $2\ln(2x+1)$ n'est une primitive de $\dfrac{4}{2x+1}$ que sur $\left]-\dfrac{1}{2}\,;+\infty\right[$, jamais sur $\mathbb{R}$ tout entier. Au Baccalauréat, l'énoncé précise presque toujours l'intervalle de travail : lis-le attentivement et justifie la positivité de la quantité sous le logarithme.
\end{attention}

\begin{savaistu}[Complément utile : la notation $\ln|x|$ (hors programme strict)]
Que se passe-t-il pour $x < 0$ ? La fonction $x \mapsto \ln(-x)$ est définie sur $]-\infty\,;0[$ et, par dérivation composée, sa dérivée est $\dfrac{-1}{-x} = \dfrac{1}{x}$. Ainsi, $\dfrac{1}{x}$ admet aussi des primitives sur $]-\infty\,;0[$ : les fonctions $\ln(-x) + k$. On regroupe les deux situations dans la notation condensée : une primitive de $\dfrac{1}{x}$ sur tout intervalle ne contenant pas $0$ est $\ln|x|$. Cette notation est très utilisée en classe de Première année universitaire et dans les concours ; retiens simplement que la valeur absolue garantit que la quantité sous le logarithme est toujours strictement positive.
\end{savaistu}

\courssub{Tableaux récapitulatifs : dérivées et primitives}

Voici l'outil de révision idéal : deux tableaux à double entrée résumant toutes les formules de cette section. Recopie-les dans ton cahier de formules et apprends-les dans les \emph{deux sens} de lecture.

\begin{center}
\renewcommand{\arraystretch}{2.2}
\begin{tabularx}{\linewidth}{|l|X|X|}
\hline
\rowcolor{popblueL} \textbf{Fonction $f(x)$} & \textbf{Dérivée $f'(x)$} & \textbf{Conditions} \\
\hline
$\ln x$ & $\dfrac{1}{x}$ & $x > 0$ \\
\hline
$\ln(ax+b)$ & $\dfrac{a}{ax+b}$ & $ax+b > 0$ \\
\hline
$\ln\big(u(x)\big)$ & $\dfrac{u'(x)}{u(x)}$ & $u$ dérivable, $u(x) > 0$ \\
\hline
\end{tabularx}
\end{center}

\begin{center}
\renewcommand{\arraystretch}{2.2}
\begin{tabularx}{\linewidth}{|l|X|X|}
\hline
\rowcolor{popgreenL} \textbf{Fonction $f(x)$} & \textbf{Une primitive $F(x)$} & \textbf{Conditions} \\
\hline
$\dfrac{1}{x}$ & $\ln x$ & $x > 0$ \\
\hline
$\dfrac{a}{bx+c}$ \; ($b \neq 0$) & $\dfrac{a}{b}\ln(bx+c)$ & $bx+c > 0$ \\
\hline
$\dfrac{u'(x)}{u(x)}$ & $\ln\big(u(x)\big)$ & $u$ dérivable, $u(x) > 0$ \\
\hline
\end{tabularx}
\end{center}

\begin{exoresolu}[Dérivation et primitive dans un même exercice]
On considère la fonction $g$ définie par $g(x) = \ln(2x+6)$.
\begin{enumerate}
\item Déterminer l'ensemble de définition de $g$, puis calculer $g'(x)$.
\item En déduire une primitive de la fonction $h(x) = \dfrac{1}{x+3}$ sur $]-3\,;+\infty[$.
\end{enumerate}
\tcblower
\textbf{Solution détaillée.}

\textbf{1.} Il faut $2x + 6 > 0$, c'est-à-dire $x > -3$. Donc $D_g = \left]-3\,;+\infty\right[$. Avec $u(x) = 2x+6$ et $u'(x) = 2$ :
\[
g'(x) = \frac{2}{2x+6} = \frac{2}{2(x+3)} = \frac{1}{x+3}.
\]

\textbf{2.} Le calcul précédent montre exactement que $g'(x) = h(x)$ sur $]-3\,;+\infty[$. Par définition d'une primitive, $g$ est donc une primitive de $h$ :
\[
G(x) = \ln(2x+6).
\]
\emph{Autre présentation équivalente :} comme $2x+6 = 2(x+3)$, on a $\ln(2x+6) = \ln 2 + \ln(x+3)$ ; la constante $\ln 2$ disparaît par dérivation, donc $x \mapsto \ln(x+3)$ est aussi une primitive de $h$. Les deux réponses sont correctes : deux primitives d'une même fonction diffèrent d'une constante.
\end{exoresolu}

\begin{exercice}[À toi de jouer]
\begin{enumerate}
\item Dériver les fonctions suivantes en précisant leur ensemble de dérivabilité :
\[
f(x) = \ln(5x-1) \qquad ; \qquad g(x) = \ln(x^2+1).
\]
\item Déterminer une primitive de chacune des fonctions suivantes sur l'intervalle indiqué :
\[
h(x) = \frac{6}{3x+2} \ \text{sur } \left]-\tfrac{2}{3}\,;+\infty\right[ \qquad ; \qquad k(x) = \frac{1}{4x-3} \ \text{sur } \left]\tfrac{3}{4}\,;+\infty\right[.
\]
\end{enumerate}
\end{exercice}

\begin{corrige}[Exercice « À toi de jouer »]
\textbf{1.} Pour $f$ : il faut $5x - 1 > 0$, soit $x > \dfrac{1}{5}$. Alors $f'(x) = \dfrac{5}{5x-1}$ sur $\left]\dfrac{1}{5}\,;+\infty\right[$.

Pour $g$ : $x^2 + 1 > 0$ pour tout réel $x$, donc $g$ est dérivable sur $\mathbb{R}$ et $g'(x) = \dfrac{2x}{x^2+1}$.

\textbf{2.} Pour $h$ : la dérivée du dénominateur est $3$, et $6 = 2 \times 3$, donc $h(x) = 2 \times \dfrac{3}{3x+2}$. Une primitive est $H(x) = 2\ln(3x+2)$.

Pour $k$ : la dérivée du dénominateur est $4$, et $1 = \dfrac{1}{4} \times 4$, donc $k(x) = \dfrac{1}{4} \times \dfrac{4}{4x-3}$. Une primitive est $K(x) = \dfrac{1}{4}\ln(4x-3)$. Vérifie en redérivant : $K'(x) = \dfrac{1}{4} \times \dfrac{4}{4x-3} = \dfrac{1}{4x-3}$. \checkmark
\end{corrige}

\begin{aretenir}[L'essentiel de la section]
\begin{itemize}
\item $(\ln x)' = \dfrac{1}{x}$ pour $x > 0$ — démonstration exigible : dériver $\mathrm{e}^{\ln x} = x$.
\item $(\ln u)' = \dfrac{u'}{u}$ pour $u > 0$ ; en particulier $\big(\ln(ax+b)\big)' = \dfrac{a}{ax+b}$.
\item Une primitive de $\dfrac{u'}{u}$ est $\ln u$ ; une primitive de $\dfrac{a}{bx+c}$ est $\dfrac{a}{b}\ln(bx+c)$ là où $bx+c > 0$.
\item Réflexe permanent : \textbf{toujours} vérifier la condition de positivité de la quantité sous le logarithme, et \textbf{toujours} vérifier une primitive en la redérivant.
\end{itemize}
Ces outils en main, nous sommes prêts pour la dernière section : l'étude complète de fonctions du type $\ln(ax+b)$, de l'ensemble de définition jusqu'au tracé de la courbe.
\end{aretenir}

\coursec{Étude complète de fonctions du type ln(ax + b)}

Dans la section précédente, tu as appris à dériver une fonction du type $x \mapsto \ln(ax+b)$ et à trouver des primitives de $x \mapsto \dfrac{a}{bx+c}$. Ces outils de calcul ne sont pas une fin en soi : ils servent précisément à \emph{étudier} les fonctions logarithmes, c'est-à-dire à déterminer leur comportement complet (domaine, limites, variations, courbe). C'est exactement ce que le Baccalauréat te demandera : « Étudier la fonction $f$ définie par $f(x) = \ln(2x+4)$ et tracer sa courbe représentative ». Cette dernière section te donne la méthode complète, étape par étape, puis deux exemples de synthèse de niveau Bac.

\courssub{Le plan d'étude général}

Étudier une fonction, c'est répondre dans l'ordre à six questions précises. L'ordre compte : chaque étape s'appuie sur les précédentes.

\begin{methode}[Plan d'étude d'une fonction du type $f(x) = \ln(ax+b)$]
\begin{enumerate}
\item \textbf{Ensemble de définition} : résoudre l'inéquation $ax + b > 0$ (inégalité \emph{stricte}, car $\ln$ n'est définie que sur $]0\,;\,+\infty[$).
\item \textbf{Limites aux bornes} : calculer la limite de $f$ à chaque borne ouverte du domaine, en utilisant les limites usuelles de $\ln$.
\item \textbf{Dérivée} : calculer $f'(x) = \dfrac{a}{ax+b}$ et étudier son signe.
\item \textbf{Tableau de variation} : synthétiser limites et variations dans un tableau complet.
\item \textbf{Points remarquables} : intersection avec l'axe des abscisses ($f(x)=0$), ordonnée à l'origine si elle existe, équation d'une tangente demandée.
\item \textbf{Tracé de la courbe} : placer l'asymptote verticale, les points remarquables, la tangente, puis dessiner la courbe en respectant le tableau de variation.
\end{enumerate}
\end{methode}

Détaillons maintenant chacune de ces étapes.

\courssub{Étape 1 : l'ensemble de définition}

\begin{propriete}[Ensemble de définition de $\ln(ax+b)$]
La fonction $f : x \mapsto \ln(ax+b)$ (avec $a \neq 0$) est définie sur l'ensemble des réels $x$ tels que $ax + b > 0$.
\begin{itemize}
\item Si $a > 0$ : $D_f = \left]-\dfrac{b}{a}\,;\,+\infty\right[$.
\item Si $a < 0$ : $D_f = \left]-\infty\,;\,-\dfrac{b}{a}\right[$.
\end{itemize}
\end{propriete}

L'idée est simple : le logarithme népérien n'accepte que des nombres \cle{strictement positifs}. La quantité $ax+b$ joue le rôle de « l'intérieur » du logarithme : c'est elle qui doit être strictement positive. On résout donc une inéquation du premier degré, avec le piège classique du changement de sens de l'inégalité.

\begin{attention}[Le piège du signe de $a$]
Pour déterminer le domaine de $f(x) = \ln(-2x+6)$, on résout $-2x + 6 > 0$, c'est-à-dire $-2x > -6$. En divisant par $-2$ (\textbf{nombre négatif}), on doit \cle{inverser le sens de l'inégalité} :
\[ -2x + 6 > 0 \iff x < 3 \quad \text{donc} \quad D_f = \left]-\infty\,;\,3\right[. \]
L'erreur fréquente consiste à écrire $x > 3$ en oubliant d'inverser l'inégalité. Vérifie toujours ta réponse avec une valeur test : pour $x = 0$, on a $-2\times 0 + 6 = 6 > 0$, donc $0$ doit appartenir au domaine — ce qui confirme $D_f = \left]-\infty\,;\,3\right[$.
\end{attention}

\courssub{Étape 2 : limites aux bornes et asymptote verticale}

Le domaine de $f(x) = \ln(ax+b)$ possède une borne finie $-\dfrac{b}{a}$ (ouverte) et une borne infinie. Aux deux bornes, on utilise les limites usuelles par composition.

\textbf{À la borne finie.} Quand $x$ tend vers $-\dfrac{b}{a}$ du côté du domaine, la quantité $ax+b$ tend vers $0$ \emph{par valeurs positives}. Or $\lim_{X \to 0^+} \ln X = -\infty$. Donc :
\[ \lim_{x \to -\frac{b}{a}} \ln(ax+b) = -\infty. \]
Graphiquement, la courbe « plonge » vers le bas en s'approchant de la droite verticale d'équation $x = -\dfrac{b}{a}$ sans jamais la toucher : c'est une \cle{asymptote verticale}.

\begin{propriete}[Asymptote verticale]
Pour tous réels $a \neq 0$ et $b$, la courbe représentative de $f(x) = \ln(ax+b)$ admet toujours une asymptote verticale d'équation
\[ x = -\frac{b}{a}. \]
C'est la droite verticale passant par la valeur qui annule $ax+b$.
\end{propriete}

\textbf{À la borne infinie.} Quand $x$ tend vers l'infini du côté du domaine, $ax+b$ tend vers $+\infty$ (car $ax+b > 0$ sur tout le domaine et $|ax+b|$ devient arbitrairement grand). Or $\lim_{X \to +\infty} \ln X = +\infty$. Donc :
\[ \lim \ln(ax+b) = +\infty \quad \text{à la borne infinie du domaine.} \]
La courbe présente une \cle{branche infinie} montante : elle monte indéfiniment, mais de plus en plus lentement (c'est la lenteur caractéristique du logarithme).

\courssub{Étape 3 : dérivée et sens de variation}

On a vu dans la section précédente que si $f(x) = \ln(ax+b)$, alors :
\[ f'(x) = \frac{a}{ax+b}. \]
Le point capital est le suivant : sur l'ensemble de définition, le dénominateur $ax+b$ est \textbf{strictement positif} (c'est précisément la condition d'existence de $f$). Le signe de $f'(x)$ est donc le signe de $a$, qui est \emph{constant} :

\begin{propriete}[Monotonie de $\ln(ax+b)$]
La fonction $f(x) = \ln(ax+b)$ est :
\begin{itemize}
\item \textbf{strictement croissante} sur son ensemble de définition si $a > 0$ ;
\item \textbf{strictement décroissante} sur son ensemble de définition si $a < 0$.
\end{itemize}
\end{propriete}

C'est une grande simplification par rapport aux fonctions polynômes : il n'y a jamais de changement de variation, jamais d'extremum. Une fonction $\ln(ax+b)$ est \emph{toujours monotone}. Cela se comprend intuitivement : la fonction $\ln$ elle-même est croissante ; multiplier l'intérieur par $a > 0$ conserve l'ordre, multiplier par $a < 0$ le renverse.

\courssub{Étape 4 : le tableau de variation}

On rassemble tout dans un tableau. Pour $a > 0$, il a l'allure suivante (6 colonnes : borne finie, intervalle, borne infinie) :

\begin{center}
\begin{tabular}{|c|ccccc|}
\hline
$x$ & $-\dfrac{b}{a}$ & & & & $+\infty$ \\
\hline
$f'(x)$ & & $+$ & & $+$ & \\
\hline
$f(x)$ & $-\infty$ & $\nearrow$ & & $\nearrow$ & $+\infty$ \\
\hline
\end{tabular}
\end{center}

La double barre verticale sous $-\dfrac{b}{a}$ (ici représentée par le bord du tableau) matérialise l'asymptote : la valeur $-\dfrac{b}{a}$ est \emph{interdite}. Pour $a < 0$, le domaine devient $\left]-\infty\,;\,-\dfrac{b}{a}\right[$, la dérivée est négative et la fonction décroît de $+\infty$ vers $-\infty$ (tableau symétrique en sens inverse).

\courssub{Étape 5 : points remarquables et tangentes}

\textbf{Intersection avec l'axe des abscisses.} On résout $f(x) = 0$ :
\[ \ln(ax+b) = 0 \iff ax + b = \mathrm{e}^0 = 1 \iff x = \frac{1-b}{a}. \]
La courbe coupe donc l'axe des abscisses au point $\left(\dfrac{1-b}{a}\,;\,0\right)$. Retiens l'idée clé : $\ln(\text{quelque chose}) = 0$ exactement quand ce « quelque chose » vaut $1$.

\textbf{Ordonnée à l'origine.} Elle n'existe que si $0 \in D_f$, c'est-à-dire si $b > 0$. Dans ce cas : $f(0) = \ln b$.

\textbf{Tangente en un point.} L'équation de la tangente à la courbe de $f$ au point d'abscisse $x_0$ est :
\[ y = f'(x_0)\,(x - x_0) + f(x_0). \]
Un cas fondamental à connaître par cœur : la tangente à la courbe de $y = \ln x$ au point $(1\,;\,0)$. Comme $f(1) = \ln 1 = 0$ et $f'(1) = \dfrac{1}{1} = 1$, on obtient :
\[ y = 1 \times (x - 1) + 0 = x - 1. \]

\begin{aretenir}[Tangente de référence]
La tangente à la courbe de $\ln$ au point $(1\,;\,0)$ a pour équation $y = x - 1$. Elle revient très souvent dans les sujets de Bac.
\end{aretenir}

\courssub{Exemple de synthèse 1 : étude complète de $f(x) = \ln(2x+4)$}

Appliquons la méthode en suivant scrupuleusement les six étapes.

\textbf{Étape 1 — Ensemble de définition.} On résout $2x + 4 > 0 \iff x > -2$. Donc $D_f = \left]-2\,;\,+\infty\right[$.

\textbf{Étape 2 — Limites aux bornes.} Quand $x \to -2^+$, $2x+4 \to 0^+$, donc $\lim_{x \to -2^+} f(x) = -\infty$ : la droite d'équation $x = -2$ est asymptote verticale. Quand $x \to +\infty$, $2x+4 \to +\infty$, donc $\lim_{x \to +\infty} f(x) = +\infty$.

\textbf{Étape 3 — Dérivée.} $f'(x) = \dfrac{2}{2x+4}$. Sur $D_f$, $2x+4 > 0$, donc $f'(x) > 0$ : $f$ est strictement croissante.

\textbf{Étape 4 — Tableau de variation.} On ajoute la valeur $x = -\frac{3}{2}$ (zéro de la fonction) pour précision.

\begin{center}
\begin{tabular}{|c|ccccc|}
\hline
$x$ & $-2$ & & $-\dfrac{3}{2}$ & & $+\infty$ \\
\hline
$f'(x)$ & & $+$ & $2$ & $+$ & \\
\hline
$f(x)$ & $-\infty$ & $\nearrow$ & $0$ & $\nearrow$ & $+\infty$ \\
\hline
\end{tabular}
\end{center}

\textbf{Étape 5 — Points remarquables.} $f(x) = 0 \iff 2x+4 = 1 \iff x = -\dfrac{3}{2}$. La courbe coupe l'axe des abscisses en $\left(-\dfrac{3}{2}\,;\,0\right)$. Comme $b = 4 > 0$, l'ordonnée à l'origine existe : $f(0) = \ln 4 \approx 1,39$. Tangente en $x_0 = -\dfrac{3}{2}$ : on a $f\left(-\dfrac{3}{2}\right) = 0$ et $f'\left(-\dfrac{3}{2}\right) = \dfrac{2}{2 \times (-\frac{3}{2}) + 4} = \dfrac{2}{1} = 2$, d'où :
\[ y = 2\left(x + \frac{3}{2}\right) = 2x + 3. \]

\textbf{Étape 6 — Tracé.} On place l'asymptote $x = -2$, le point $\left(-\dfrac{3}{2}\,;\,0\right)$, le point $(0\,;\,\ln 4)$, la tangente $y = 2x+3$, puis on trace la courbe croissante qui plonge vers $-\infty$ près de l'asymptote et monte lentement vers $+\infty$.

\begin{popfigure}
\begin{center}
\begin{tikzpicture}
\begin{axis}[
  width=0.95\linewidth, height=7.5cm,
  axis lines=middle,
  xlabel={$x$}, ylabel={$y$},
  xmin=-4.5, xmax=4.5, ymin=-4, ymax=4.5,
  xtick={-2,-1.5,0,2}, ytick={-2,2},
  grid=both, grid style={popline!40},
  legend style={at={(0.02,0.98)}, anchor=north west, font=\small}
]
\addplot[popcurve,domain=-1.93:4.3,samples=200]{ln(2*x+4)};
\addlegendentry{$y=\ln(2x+4)$}
\addplot[poporange, dashed, thick] coordinates {(-2,-4) (-2,4.5)};
\addlegendentry{Asymptote $x=-2$}
\addplot[poppurple, thick, domain=-2.1:0.6]{2*x+3};
\addlegendentry{Tangente $y=2x+3$}
\addplot[only marks, mark=*, popink, mark size=2pt] coordinates {(-1.5,0) (0,1.386)};
\node[popink, anchor=south west, font=\small] at (axis cs:-1.5,0.15){$\left(-\frac{3}{2}\,;\,0\right)$};
\node[popink, anchor=west, font=\small] at (axis cs:0.1,1.5){$(0\,;\,\ln 4)$};
\end{axis}
\end{tikzpicture}
\end{center}
\legende{Courbe de $f(x)=\ln(2x+4)$ : asymptote verticale $x=-2$, passage par $\left(-\frac{3}{2}\,;\,0\right)$, tangente $y=2x+3$ en ce point.}
\end{popfigure}

\courssub{Complément type Bac : position de la courbe de ln par rapport à sa tangente}

Voici un grand classique du Baccalauréat, qui illustre la puissance de l'étude de fonctions pour démontrer des inégalités.

\begin{propriete}[Inégalité fondamentale du logarithme — démonstration exigible]
Pour tout réel $x > 0$ : $\ln x \leq x - 1$, avec égalité si et seulement si $x = 1$.
\end{propriete}

\begin{experience}[Démonstration guidée de l'inégalité $\ln x \leq x - 1$]
Posons $g(x) = x - 1 - \ln x$ pour $x > 0$. Nous allons montrer que $g(x) \geq 0$.
\begin{enumerate}
\item \textbf{Dérivée} : $g'(x) = 1 - \dfrac{1}{x} = \dfrac{x-1}{x}$.
\item \textbf{Signe de $g'$} : comme $x > 0$, le signe de $g'(x)$ est celui de $x-1$. Donc $g'(x) < 0$ sur $]0\,;\,1[$ et $g'(x) > 0$ sur $]1\,;\,+\infty[$.
\item \textbf{Variations} : $g$ est strictement décroissante sur $]0\,;\,1]$, strictement croissante sur $[1\,;\,+\infty[$. Elle admet donc un \cle{minimum} en $x = 1$.
\item \textbf{Conclusion} : ce minimum vaut $g(1) = 1 - 1 - \ln 1 = 0$. Ainsi, pour tout $x > 0$, $g(x) \geq 0$, c'est-à-dire $\ln x \leq x - 1$. L'égalité n'a lieu qu'au minimum, c'est-à-dire en $x = 1$.
\end{enumerate}
\end{experience}

L'interprétation géométrique est remarquable : la courbe de $\ln$ est \emph{toujours située en dessous} de sa tangente au point $(1\,;\,0)$, et ne la touche qu'en ce point. On dit que la fonction $\ln$ est \cle{concave}.

\begin{popfigure}
\begin{center}
\begin{tikzpicture}
\begin{axis}[
  width=0.95\linewidth, height=7.5cm,
  axis lines=middle,
  xlabel={$x$}, ylabel={$y$},
  xmin=-1, xmax=5, ymin=-3.5, ymax=4,
  xtick={1,2,3,4}, ytick={-2,-1,1,2,3},
  grid=both, grid style={popline!40},
  legend style={at={(0.02,0.98)}, anchor=north west, font=\small}
]
\addplot[popcurve,domain=0.06:4.8,samples=200]{ln(x)};
\addlegendentry{$y=\ln x$}
\addplot[poporange, thick, domain=-0.5:4.5]{x-1};
\addlegendentry{Tangente $y=x-1$}
\addplot[only marks, mark=*, popink, mark size=2pt] coordinates {(1,0)};
\node[popink, anchor=south, font=\small] at (axis cs:1,0.15){$(1\,;\,0)$};
\end{axis}
\end{tikzpicture}
\end{center}
\legende{La courbe de $\ln$ est entièrement située sous sa tangente $y = x-1$ au point $(1\,;\,0)$ : pour tout $x>0$, $\ln x \leq x-1$.}
\end{popfigure}

\courssub{Exemple de synthèse 2 : étude de $g(x) = x - \ln x$}

Les sujets de Bac aiment mélanger le logarithme avec d'autres fonctions. Voici un exercice résolu complet, dans l'esprit de l'épreuve.

\begin{exoresolu}[Étude de $g(x) = x - \ln x$]
Soit $g$ la fonction définie par $g(x) = x - \ln x$.
\begin{enumerate}
\item Déterminer l'ensemble de définition de $g$ et les limites aux bornes.
\item Calculer $g'(x)$, dresser le tableau de variation de $g$.
\item En déduire le signe de $g(x)$ sur son domaine.
\end{enumerate}
\tcblower
\textbf{Solution détaillée.}

\textbf{1. Ensemble de définition et limites.} Le terme $\ln x$ impose $x > 0$, donc $D_g = \left]0\,;\,+\infty\right[$.

En $0^+$ : $\ln x \to -\infty$, donc $-\ln x \to +\infty$ et $\lim_{x \to 0^+} g(x) = +\infty$. L'axe des ordonnées ($x = 0$) est asymptote verticale.

En $+\infty$ : attention à la forme indéterminée « $+\infty - \infty$ ». On factorise par $x$ :
\[ g(x) = x\left(1 - \frac{\ln x}{x}\right). \]
Par croissances comparées, $\lim_{x \to +\infty} \dfrac{\ln x}{x} = 0$, donc la parenthèse tend vers $1$ et $\lim_{x \to +\infty} g(x) = +\infty$.

\textbf{2. Dérivée et variations.} $g'(x) = 1 - \dfrac{1}{x} = \dfrac{x-1}{x}$. Comme $x > 0$, $g'(x)$ a le signe de $x-1$ : négatif sur $]0\,;\,1[$, nul en $1$, positif sur $]1\,;\,+\infty[$. De plus $g(1) = 1 - \ln 1 = 1$.

\begin{center}
\begin{tabular}{|c|ccccccc|}
\hline
$x$ & $0$ & & & $1$ & & & $+\infty$ \\
\hline
$g'(x)$ & & $-$ & & $0$ & & $+$ & \\
\hline
$g(x)$ & $+\infty$ & $\searrow$ & & $1$ & & $\nearrow$ & $+\infty$ \\
\hline
\end{tabular}
\end{center}

\textbf{3. Signe de $g$.} D'après le tableau, $g$ admet sur $]0\,;\,+\infty[$ un minimum égal à $1$. Or $1 > 0$, donc pour tout $x > 0$, $g(x) \geq 1 > 0$ : la fonction $g$ est \textbf{strictement positive} sur tout son domaine. En particulier, l'équation $x = \ln x$ n'a aucune solution — la courbe de $\ln$ ne rencontre jamais la droite $y = x$.
\end{exoresolu}

\begin{attention}[Formes indéterminées avec le logarithme]
Face à une limite du type $x - \ln x$ en $+\infty$, ne conclus jamais « $+\infty - \infty = 0$ » : c'est une forme indéterminée. La bonne réaction est de \cle{factoriser par le terme dominant} (ici $x$) pour faire apparaître le quotient $\dfrac{\ln x}{x}$, dont la limite $0$ est un résultat de cours (croissances comparées).
\end{attention}

\courssub{Pour aller plus loin : les échelles logarithmiques}

\begin{savaistu}[Le logarithme mesure le monde qui nous entoure]
Pourquoi le logarithme est-il si important en dehors de la classe ? Parce que de nombreuses grandeurs varient sur des plages gigantesques, et le logarithme les « compresse » sur une échelle lisible.
\begin{itemize}
\item \textbf{L'échelle de Richter} (séismes) : la magnitude est un logarithme de l'énergie libérée. Un séisme de magnitude $7$ est $10$ fois plus puissant qu'un séisme de magnitude $6$, et $100$ fois plus qu'un séisme de magnitude $5$.
\item \textbf{Le pH en chimie} : $\mathrm{pH} = -\log[\mathrm{H}^+]$. Un jus de citron (pH $2$) est $100\,000$ fois plus acide que l'eau pure (pH $7$).
\item \textbf{Les décibels en acoustique} : l'intensité sonore est mesurée par un logarithme. Quand le son double d'intensité, le niveau n'augmente que d'environ $3$ dB.
\end{itemize}
Même au Cameroun, les sismologues qui surveillent l'activité du mont Cameroun, volcan actif de la région du Sud-Ouest, utilisent ces échelles logarithmiques pour classer les secousses enregistrées.
\end{savaistu}

\begin{exercice}[Pour t'entraîner]
\begin{enumerate}
\item Étudier complètement la fonction $f(x) = \ln(-2x+6)$ : ensemble de définition, limites, dérivée, variations, intersection avec l'axe des abscisses, puis tracer sa courbe.
\item Soit $h(x) = \ln(3x - 6)$. Déterminer son domaine, son asymptote verticale, le point où sa courbe coupe l'axe des abscisses, et l'équation de la tangente en ce point.
\item Démontrer que pour tout $x > 0$, $\ln x \geq 1 - \dfrac{1}{x}$ (on étudiera la fonction $k(x) = \ln x - 1 + \dfrac{1}{x}$).
\end{enumerate}
\end{exercice}

\begin{corrige}[Exercices proposés]
\textbf{1.} Domaine : $-2x+6 > 0 \iff x < 3$, donc $D_f = \left]-\infty\,;\,3\right[$. Limites : $\lim_{x \to 3^-} f(x) = -\infty$ (asymptote $x = 3$) et $\lim_{x \to -\infty} f(x) = +\infty$. Dérivée : $f'(x) = \dfrac{-2}{-2x+6} < 0$ sur $D_f$, donc $f$ strictement décroissante. Intersection avec l'axe des abscisses : $-2x+6 = 1 \iff x = \dfrac{5}{2}$.

\textbf{2.} Domaine : $3x - 6 > 0 \iff x > 2$, donc $D_h = \left]2\,;\,+\infty\right[$ ; asymptote $x = 2$. Intersection : $3x - 6 = 1 \iff x = \dfrac{7}{3}$. Dérivée : $h'(x) = \dfrac{3}{3x-6}$, d'où $h'\left(\dfrac{7}{3}\right) = \dfrac{3}{1} = 3$ et $h\left(\dfrac{7}{3}\right) = 0$ ; tangente : $y = 3\left(x - \dfrac{7}{3}\right) = 3x - 7$.

\textbf{3.} $k'(x) = \dfrac{1}{x} - \dfrac{1}{x^2} = \dfrac{x-1}{x^2}$, du signe de $x-1$. Donc $k$ décroît sur $]0\,;\,1]$ puis croît sur $[1\,;\,+\infty[$, avec minimum $k(1) = 0$. Ainsi $k(x) \geq 0$ pour tout $x > 0$, ce qui prouve l'inégalité demandée.
\end{corrige}

\begin{aretenir}[L'essentiel de la section]
\begin{itemize}
\item L'étude de $f(x) = \ln(ax+b)$ suit un plan en six étapes : domaine, limites, dérivée, tableau, points remarquables, tracé.
\item Le domaine s'obtient en résolvant $ax+b > 0$ — attention à inverser l'inégalité si $a < 0$.
\item La courbe admet toujours une asymptote verticale $x = -\dfrac{b}{a}$, et $f$ est monotone : croissante si $a > 0$, décroissante si $a < 0$.
\item $f(x) = 0 \iff ax+b = 1$ : la courbe coupe l'axe des abscisses en $x = \dfrac{1-b}{a}$.
\item La tangente à $y = \ln x$ en $(1\,;\,0)$ est $y = x - 1$, et la courbe reste toujours en dessous : $\ln x \leq x - 1$ pour tout $x > 0$.
\end{itemize}
\end{aretenir}

\newpage

\coursec{Activité d'intégration}

\courssub{Objectifs de l'activité}

À l'issue de cette activité d'intégration, tu seras capable de :
\begin{itemize}
\item mobiliser les \cle{propriétés algébriques} de la fonction logarithme népérien ($\ln(a^n) = n\ln a$, $\ln(ab) = \ln a + \ln b$) dans un contexte financier concret ;
\item résoudre une \cle{équation} et une \cle{inéquation} à l'aide de la fonction $\ln$, en respectant le sens de variation (croissance stricte) de cette fonction ;
\item étudier et dériver une fonction du type $\ln(f(x))$ où $f$ est une exponentielle, et interpréter concrètement la dérivée obtenue ;
\item comparer deux situations d'évolution et rédiger un \cle{avis argumenté} destiné à un non-spécialiste, compétence essentielle de la série A.
\end{itemize}

\courssub{Situation}

Mme Etoundi, commerçante au marché B de Bafoussam, a vendu sa récolte de pommes de terre et dispose d'un capital de $500\,000$ FCFA. Une coopérative de microfinance agréée COBAC lui propose un placement à \cle{intérêts composés} au taux annuel de $8\,\%$ : chaque année, le capital est multiplié par $1,08$. Au bout de $n$ années (avec $n \in \mathbb{N}$), le capital disponible est donc :
\[ C(n) = 500\,000 \times (1,08)^n \quad \text{(en FCFA)}. \]

Mme Etoundi se pose plusieurs questions :
\begin{itemize}
\item au bout de combien d'années son capital aura-t-il \textbf{doublé}, c'est-à-dire atteint $1\,000\,000$ FCFA ?
\item elle rêve d'acheter une camionnette à $1\,500\,000$ FCFA pour transporter ses marchandises : au bout de combien d'années disposera-t-elle \textbf{d'au moins} cette somme ?
\item une coopérative concurrente propose un taux de $5\,\%$ l'an. Laquelle des deux offres permet de doubler le capital \textbf{en moins de 12 ans} ?
\end{itemize}

Pour affiner son étude, son neveu, étudiant en gestion, modélise le capital en \cle{continu} par la fonction définie sur $[0\,;\,20]$ par :
\[ f(x) = 500\,000\, e^{0,077x}, \]
où $x$ est le temps écoulé en années. Il s'intéresse alors à la fonction $g(x) = \ln\big(f(x)\big)$.

Tu es le(la) conseiller(ère) mathématique de Mme Etoundi. Ta mission : répondre à toutes ses questions et rédiger un court avis argumenté.

\courssub{Consignes}

\textbf{Partie A — Étude du placement à $8\,\%$}\par
\begin{enumerate}
\item \begin{enumerate}
\item Calculer $C(0)$, $C(1)$ et $C(2)$. Vérifier que la suite $(C(n))$ est géométrique et préciser sa raison.
\item Exprimer $\ln\big(C(n)\big)$ en fonction de $n$, en utilisant les propriétés algébriques de $\ln$.
\item En déduire que la suite $\big(\ln C(n)\big)$ est arithmétique ; préciser sa raison et son premier terme.
\end{enumerate}
\item \begin{enumerate}
\item Résoudre dans $\mathbb{R}$ l'équation $C(n) = 1\,000\,000$ en appliquant la fonction $\ln$ aux deux membres.
\item En déduire le plus petit entier $n$ tel que le capital ait doublé. Conclure par une phrase pour Mme Etoundi.
\end{enumerate}
\item Résoudre dans $\mathbb{R}$ l'inéquation $C(n) \geq 1\,500\,000$, puis déterminer le nombre minimal d'années (entier) permettant d'acheter la camionnette.
\end{enumerate}

\textbf{Partie B — Modélisation continue}\par
\begin{enumerate}
\setcounter{enumi}{3}
\item \begin{enumerate}
\item Simplifier l'expression de $g(x) = \ln\big(f(x)\big)$ pour $x \in [0\,;\,20]$.
\item Justifier que $g$ est dérivable sur $[0\,;\,20]$ et calculer $g'(x)$.
\item Étudier le signe de $g'(x)$ et dresser le tableau de variations de $g$ sur $[0\,;\,20]$.
\item Interpréter concrètement la valeur de $g'(x)$ pour le capital de Mme Etoundi.
\end{enumerate}
\end{enumerate}

\textbf{Partie C — Comparaison des offres et avis}\par
\begin{enumerate}
\setcounter{enumi}{4}
\item Avec un taux de $5\,\%$, le capital est $D(n) = 500\,000 \times (1,05)^n$. Résoudre $D(n) = 1\,000\,000$ et déterminer le nombre d'années entières nécessaires pour doubler le capital. Laquelle des deux coopératives permet de doubler le capital en moins de 12 ans ?
\item \textbf{Production écrite.} Rédige un avis d'une dizaine de lignes à remettre à Mme Etoundi : tu y présenteras les résultats chiffrés, la comparaison des deux offres et ta recommandation finale, dans un langage accessible à une non-spécialiste.
\end{enumerate}

\courssub{Ressources à mobiliser}

\begin{aretenir}[Boîte à outils de la leçon]
\begin{itemize}
\item \textbf{Définition} : pour tout $a > 0$, $\ln a$ est l'unique réel tel que $e^{\ln a} = a$ ; en particulier $\ln(e^x) = x$ pour tout réel $x$.
\item \textbf{Propriétés algébriques} : pour $a, b > 0$ et $n \in \mathbb{Z}$ :
\[ \ln(ab) = \ln a + \ln b, \qquad \ln\!\left(\frac{a}{b}\right) = \ln a - \ln b, \qquad \ln(a^n) = n \ln a. \]
\item \textbf{Équations et inéquations} : $\ln$ est \emph{strictement croissante} sur $]0\,;\,+\infty[$, donc pour $a, b > 0$ :
\[ \ln a = \ln b \iff a = b \qquad \text{et} \qquad \ln a \leq \ln b \iff a \leq b. \]
\item \textbf{Limites usuelles} (exigibles au Bac) :
\[ \lim_{x \to 0^+} \ln x = -\infty, \qquad \lim_{x \to +\infty} \ln x = +\infty. \]
\item \textbf{Dérivation} : si $u$ est dérivable et strictement positive sur $I$, alors $(\ln u)' = \dfrac{u'}{u}$. En particulier $\big(\ln(ax+b)\big)' = \dfrac{a}{ax+b}$ pour $ax+b>0$.
\item \textbf{Primitives} : pour $a,b,c \in \mathbb{R}$, $b \neq 0$, sur tout intervalle où $bx+c > 0$ :
\[ \int \frac{a}{bx+c} \, dx = \frac{a}{b} \ln(bx+c) + k \quad (k \in \mathbb{R}). \]
\item \textbf{Suites} : $(u_n)$ est géométrique de raison $q$ si $u_{n+1} = q\, u_n$ ; arithmétique de raison $r$ si $u_{n+1} = u_n + r$.
\end{itemize}
\end{aretenir}

\begin{savaistu}[Taux actuariel vs taux nominal]
Le coefficient $0,077$ utilisé dans la modélisation continue $f(x)=500\,000 e^{0,077x}$ est une valeur approchée du \cle{taux actuariel} (ou taux continu) $r$ tel que $e^r = 1,08$. On a exactement $r = \ln 1,08 \approx 0,07696$. C'est le taux qui, composé en continu, produit le même effet qu'un taux nominal de $8\,\%$ composé annuellement. Les banques et institutions de microfinance (COBAC, BEAC) utilisent ce formalisme pour comparer des produits aux fréquences de capitalisation différentes.
\end{savaistu}

\courssub{Démarche et corrigé commenté}

\begin{exoresolu}[Corrigé complet de l'activité]
\textbf{Partie A.}
\begin{enumerate}
\item \begin{enumerate}
\item $C(0) = 500\,000 \times (1,08)^0 = 500\,000$ FCFA.\\
$C(1) = 500\,000 \times 1,08 = 540\,000$ FCFA.\\
$C(2) = 500\,000 \times (1,08)^2 = 500\,000 \times 1,1664 = 583\,200$ FCFA.\\
Pour tout $n \in \mathbb{N}$, $C(n+1) = 500\,000 \times (1,08)^{n+1} = 1,08 \times \big(500\,000 \times (1,08)^n\big) = 1,08 \times C(n)$.\\
La suite $(C(n))$ est donc \textbf{géométrique} de raison $q = 1,08$ et de premier terme $C(0) = 500\,000$.
\item Pour tout $n \in \mathbb{N}$, $C(n) > 0$, donc $\ln C(n)$ est bien défini. En utilisant les propriétés du logarithme :
\[ \ln C(n) = \ln\big(500\,000 \times (1,08)^n\big) = \ln 500\,000 + \ln\big((1,08)^n\big) = \ln 500\,000 + n \ln 1,08. \]
\item Posons $v_n = \ln C(n)$. Alors $v_{n+1} - v_n = \big(\ln 500\,000 + (n+1)\ln 1,08\big) - \big(\ln 500\,000 + n\ln 1,08\big) = \ln 1,08$.\\
Cette différence est constante : la suite $(v_n)$ est \textbf{arithmétique} de raison $r = \ln 1,08 \approx 0,07696$ et de premier terme $v_0 = \ln 500\,000 \approx 13,1224$.
\end{enumerate}
\item \begin{enumerate}
\item L'équation $C(n) = 1\,000\,000$ s'écrit $500\,000 \times (1,08)^n = 1\,000\,000$, soit $(1,08)^n = 2$.\\
Les deux membres sont strictement positifs, on peut appliquer $\ln$ (qui est bijective de $]0;+\infty[$ sur $\mathbb{R}$) :
\[ \ln\big((1,08)^n\big) = \ln 2 \iff n \ln 1,08 = \ln 2 \iff n = \frac{\ln 2}{\ln 1,08}. \]
Numériquement : $\ln 2 \approx 0,69315$, $\ln 1,08 \approx 0,07696$, donc $n \approx \dfrac{0,69315}{0,07696} \approx 9,006$.
\item Comme $n$ représente un nombre d'années entières et que la suite $(C(n))$ est strictement croissante (raison $>1$), le capital a doublé dès que $n \ge 9,006\ldots$ Le plus petit entier convenable est \textbf{$n = 10$}.\\
Vérification : $C(9) = 500\,000 \times 1,08^9 \approx 999\,563$ FCFA $< 1\,000\,000$ ; $C(10) \approx 1\,079\,462$ FCFA $> 1\,000\,000$.\\
\textbf{Conclusion :} Mme Etoundi verra son capital doubler au bout de \textbf{10 ans}.
\end{enumerate}
\item L'inéquation $C(n) \geq 1\,500\,000$ s'écrit $500\,000 \times (1,08)^n \geq 1\,500\,000$, soit $(1,08)^n \geq 3$.\\
Comme $\ln$ est strictement croissante sur $]0;+\infty[$, on a l'équivalence :
\[ (1,08)^n \geq 3 \iff \ln\big((1,08)^n\big) \geq \ln 3 \iff n \ln 1,08 \geq \ln 3. \]
Comme $\ln 1,08 > 0$, la division ne change pas le sens de l'inégalité :
\[ n \geq \frac{\ln 3}{\ln 1,08} \approx \frac{1,09861}{0,07696} \approx 14,27. \]
Le plus petit entier $n$ satisfaisant cette condition est \textbf{$n = 15$}.\\
Vérification : $C(14) \approx 1\,468\,597$ FCFA $< 1\,500\,000$ ; $C(15) \approx 1\,586\,085$ FCFA.\\
Mme Etoundi pourra acheter sa camionnette au bout de \textbf{15 ans}.
\end{enumerate}

\textbf{Partie B.}
\begin{enumerate}
\setcounter{enumi}{3}
\item \begin{enumerate}
\item Pour $x \in [0\,;\,20]$, $f(x) = 500\,000\, e^{0,077x} > 0$, donc $g(x) = \ln f(x)$ est bien définie.
\[ g(x) = \ln\big(500\,000\, e^{0,077x}\big) = \ln 500\,000 + \ln\big(e^{0,077x}\big) = \ln 500\,000 + 0,077x. \]
\item La fonction $g$ est une fonction \textbf{affine} (somme d'une constante et d'une fonction linéaire), donc elle est dérivable sur $\mathbb{R}$, et a fortiori sur $[0\,;\,20]$. Sa dérivée est $g'(x) = 0,077$.\\
(Vérification par la formule de dérivation composée : $u(x) = 500\,000 e^{0,077x}$, $u'(x) = 500\,000 \times 0,077 e^{0,077x}$, donc $\frac{u'(x)}{u(x)} = 0,077 = g'(x)$.)
\item $g'(x) = 0,077 > 0$ pour tout $x \in [0\,;\,20]$. La fonction $g$ est donc \textbf{strictement croissante} sur $[0\,;\,20]$.\\
Tableau de variations de $g$ :
\begin{center}
\begin{tabular}{|c|ccc|}
\hline
$x$ & $0$ & & $20$ \\ \hline
$g'(x)$ & & $+$ & \\ \hline
$g(x)$ & $\ln 500\,000 \approx 13,12$ & $\nearrow$ & $\ln 500\,000 + 1,54 \approx 14,66$ \\ \hline
\end{tabular}
\end{center}
\item La dérivée $g'(x) = 0,077$ représente le \cle{taux de croissance instantané} (ou taux actuariel) du capital. Concrètement, cela signifie qu'à tout instant, le capital croît d'environ $7,7\,\%$ par an en continu. Ce taux continu $r = 0,077$ est l'équivalent exact du taux annuel nominal de $8\,\%$ composé une fois par an, car $e^{0,077} \approx 1,08$ (plus précisément $e^{\ln 1,08} = 1,08$).
\end{enumerate}
\end{enumerate}

\textbf{Partie C.}
\begin{enumerate}
\setcounter{enumi}{4}
\item Pour la coopérative à $5\,\%$ : $D(n) = 500\,000 \times (1,05)^n$.\\
L'équation $D(n) = 1\,000\,000$ donne $(1,05)^n = 2 \iff n \ln 1,05 = \ln 2 \iff n = \dfrac{\ln 2}{\ln 1,05}$.\\
Numériquement : $\ln 1,05 \approx 0,04879$, donc $n \approx \dfrac{0,69315}{0,04879} \approx 14,21$.\\
Le plus petit entier convenable est \textbf{$n = 15$} ans ($D(14) \approx 989\,966$ FCFA, $D(15) \approx 1\,039\,464$ FCFA).\\
\textbf{Comparaison :} À $8\,\%$, le capital double en \textbf{10 ans} ; à $5\,\%$, il faut \textbf{15 ans}.\\
Seule la coopérative proposant \textbf{$8\,\%$} permet de doubler le capital en moins de 12 ans.
\item \emph{Exemple d'avis argumenté :}\\
« Chère Mme Etoundi, \\
Après étude des deux offres, voici les éléments pour votre décision. Avec la première coopérative (taux $8\,\%$), vos 500 000 FCFA deviendront 1 000 000 FCFA au bout de \textbf{10 ans} (précisément 1 079 462 FCFA) et atteindront 1 500 000 FCFA au bout de \textbf{15 ans} (1 586 085 FCFA), ce qui vous permettra d'acheter votre camionnette. \\
La coopérative concurrente (taux $5\,\%$) ne fait doubler votre capital qu'au bout de \textbf{15 ans} : elle ne répond pas à votre objectif de doublement en moins de 12 ans. \\
Je vous recommande donc vivement la première offre à $8\,\%$. Veillez toutefois à vérifier la solidité financière de l'institution (agrément COBAC) et l'absence de frais de gestion cachés qui réduiraient le rendement réel. »
\end{enumerate}
\end{exoresolu}

\begin{attention}[Pièges fréquents dans ce type de problème]
\begin{itemize}
\item \textbf{Mauvaise répartition de la puissance :} Ne jamais écrire $\ln\big(500\,000 \times (1,08)^n\big) = n \ln\big(500\,000 \times 1,08\big)$. La puissance $n$ ne porte que sur $1,08$ ; on doit écrire $\ln 500\,000 + n \ln 1,08$.
\item \textbf{Arrondi à l'entier inférieur :} Pour une inéquation « au moins » ($C(n) \geq \dots$), la solution réelle $n \approx 9,01$ ou $14,27$ impose de prendre l'\textbf{entier supérieur} ($10$ et $15$). $9$ ans ne suffisent pas pour doubler, $14$ ans ne suffisent pas pour la camionnette.
\item \textbf{Sens de l'inégalité :} Diviser par $\ln 1,08$ ne change pas le sens car $\ln 1,08 > 0$. Si le taux était négatif (ex: $0,95$), $\ln 0,95 < 0$ et la division \textbf{renverserait} l'inégalité. Vérifie toujours le signe du diviseur !
\item \textbf{Unités et notation :} Écris « FCFA » (pas de $s$ à FCFA) et utilise la virgule pour les décimales ($1,08$ et non $1.08$).
\end{itemize}
\end{attention}

\courssub{Bilan de l'activité}

Cette activité a permis de mettre en évidence plusieurs idées centrales de la leçon MA06 :
\begin{itemize}
\item Le logarithme népérien \cle{transforme les multiplications en additions} et les puissances en produits : c'est pourquoi une suite géométrique (capital à intérêts composés) devient, après passage au $\ln$, une suite \textbf{arithmétique}, beaucoup plus simple à manipuler pour trouver $n$.
\item La fonction $\ln$ est l'outil naturel pour résoudre les équations du type $a^n = b$ (inconnue dans l'exposant) : on « fait descendre » l'exposant grâce à la formule $\ln(a^n) = n\ln a$.
\item La \cle{stricte croissance} de $\ln$ sur $]0;+\infty[$ garantit que les inéquations se résolvent sans changer le sens des inégalités (à condition de diviser par un nombre positif).
\item La dérivée de $\ln(f(x))$, c'est-à-dire $\frac{f'(x)}{f(x)}$, s'interprète comme un \cle{taux de croissance relatif} (ou taux instantané), notion fondamentale en économie, en gestion et en sciences de la vie.
\item Enfin, les mathématiques de Terminale A permettent de produire un \textbf{avis argumenté, chiffré et accessible} dans une situation financière authentique du quotidien camerounais : c'est exactement l'esprit de l'approche par les compétences (savoir, savoir-faire, savoir-être).
\end{itemize}

\newpage

\coursec{Méthodes et exercices résolus}

\coursec{Définition et premières propriétés de ln}

\begin{definition}[Logarithme népérien]
La fonction \cle{logarithme népérien}, notée $\ln$, est la fonction réciproque de la fonction exponentielle $x \mapsto e^{x}$. Elle est définie sur $\left]0\,;\,+\infty\right[$ par :
\[ \forall x > 0,\quad \ln x = y \iff e^{y} = x. \]
On peut aussi la définir par l'intégrale $\ln x = \int_{1}^{x} \frac{1}{t}\,\mathrm{d}t$ pour $x > 0$.
\end{definition}

\begin{propriete}[Ensemble de définition, ensemble image et valeurs remarquables]
\begin{itemize}
\item Ensemble de définition : $D_{\ln} = \left]0\,;\,+\infty\right[$.
\item Ensemble image : $\mathbb{R}$.
\item Valeurs remarquables : $\ln 1 = 0$ ; $\ln e = 1$.
\item Signe : $\ln x > 0 \iff x > 1$ ; $\ln x < 0 \iff 0 < x < 1$.
\item La fonction $\ln$ est \cle{strictement croissante} sur $\left]0\,;\,+\infty\right[$.
\end{itemize}
\end{propriete}

\begin{savaistu}[Le logarithme au Cameroun et dans le monde]
Le logarithme népérien est omniprésent dans les sciences appliquées. Au Cameroun, on l'utilise pour :
\begin{itemize}
\item Modéliser l'\textbf{atténuation du signal} dans les fibres optiques de CAMTEL ou les antennes MTN/Orange : la puissance reçue décroît exponentiellement avec la distance, le logarithme permet de linéariser la relation.
\item Calculer le \textbf{pH} des sols agricoles (acidité) : $\mathrm{pH} = -\log_{10}[\mathrm{H}^{+}]$, lié à $\ln$ par changement de base.
\item Estimer la \textbf{croissance démographique} (Douala, Yaoundé) ou la \textbf{radioactivité} (gestion des déchets médicaux) via des modèles exponentiels dont l'étude passe par $\ln$.
\item Mesurer l'intensité des séismes (échelle de Richter, logarithmique) ou du son (décibels).
\end{itemize}
\end{savaistu}

\coursec{Propriétés algébriques de ln}

\begin{propriete}[Propriétés algébriques fondamentales]
Pour tous réels strictement positifs $a$ et $b$, et pour tout réel $n$ :
\begin{align*}
\ln(ab) &= \ln a + \ln b \quad &\text{(logarithme d'un produit)} \\
\ln\!\left(\frac{a}{b}\right) &= \ln a - \ln b \quad &\text{(logarithme d'un quotient)} \\
\ln(a^{n}) &= n\ln a \quad &\text{(logarithme d'une puissance)} \\
\ln\!\left(\frac{1}{a}\right) &= -\ln a \quad &\text{(cas particulier $n=-1$)}
\end{align*}
Ces propriétés découlent directement des propriétés de l'exponentielle ($e^{x+y}=e^{x}e^{y}$, etc.).
\end{propriete}

\begin{attention}[Pièges classiques sur les propriétés]
\begin{enumerate}
\item \textbf{Jamais} $\ln(a+b) = \ln a + \ln b$ (sauf cas trivial $a=0$ ou $b=0$ exclus du domaine).
\item \textbf{Jamais} $\ln(a-b) = \ln a - \ln b$.
\item $(\ln x)^{2} \neq \ln(x^{2}) = 2\ln x$ (pour $x>0$). Le carré du logarithme n'est pas le logarithme du carré.
\item Les propriétés ne s'appliquent que si les arguments sont \textbf{strictement positifs}.
\end{enumerate}
\end{attention}

\coursec{Équations, inéquations et inconnue auxiliaire}

\begin{methode}[Résoudre une équation avec des logarithmes]
Pour résoudre une équation faisant intervenir $\ln$, on suit toujours le même plan de bataille :
\begin{enumerate}
\item \textbf{Poser les conditions d'existence} : chaque quantité sous un $\ln$ doit être \cle{strictement positive}. On résout ces conditions et on obtient l'ensemble de résolution $D$.
\item \textbf{Regrouper} les logarithmes à l'aide des propriétés algébriques :
\[ \ln a + \ln b = \ln(ab) \ ; \quad \ln a - \ln b = \ln\!\left(\frac{a}{b}\right) \ ; \quad \ln(a^n) = n\ln a \quad (a>0). \]
\item \textbf{Se ramener à une équation du type $\ln u = \ln v$} (ou $\ln u = k$), puis utiliser l'injectivité de $\ln$ :
\[ \ln u = \ln v \iff u = v \qquad \text{et} \qquad \ln u = k \iff u = e^{k}. \]
\item \textbf{Résoudre} l'équation obtenue (souvent polynomiale).
\item \textbf{Confronter chaque solution aux conditions} de l'étape 1 : une solution qui n'appartient pas à $D$ est rejetée. Conclure par une phrase.
\end{enumerate}
\textbf{Réflexe} : ne jamais écrire $\ln(a+b) = \ln a + \ln b$ : c'est faux ! Les propriétés portent sur le produit, le quotient et les puissances, jamais sur la somme.
\end{methode}

\begin{exoresolu}[Équation avec logarithmes (type Bac A)]
Résoudre dans $\mathbb{R}$ l'équation :
\[ \ln(x+2) + \ln(x-2) = \ln 5. \]
\tcblower
\textbf{Étape 1 : conditions d'existence.}
Il faut simultanément $x+2 > 0$ et $x-2 > 0$, c'est-à-dire $x > -2$ et $x > 2$. L'ensemble de résolution est donc
\[ D = \left]2\,;\,+\infty\right[. \]

\textbf{Étape 2 : regroupement des logarithmes.}
Pour $x \in D$, on a $x+2 > 0$ et $x-2 > 0$, donc la propriété $\ln a + \ln b = \ln(ab)$ s'applique :
\[ \ln(x+2) + \ln(x-2) = \ln\bigl((x+2)(x-2)\bigr) = \ln(x^2 - 4). \]

\textbf{Étape 3 : équation sans logarithme.}
L'équation devient $\ln(x^2 - 4) = \ln 5$. La fonction $\ln$ étant injective :
\[ x^2 - 4 = 5 \iff x^2 = 9 \iff x = 3 \ \text{ou}\ x = -3. \]

\textbf{Étape 4 : confrontation aux conditions.}
\begin{itemize}
\item $x = 3$ : on a $3 \in \left]2\,;\,+\infty\right[$, donc $3$ est \textbf{acceptée}.
\item $x = -3$ : on a $-3 \notin D$ (d'ailleurs $\ln(-3+2) = \ln(-1)$ n'existe pas), donc $-3$ est \textbf{rejetée}.
\end{itemize}

\textbf{Conclusion.} L'équation admet une unique solution :
\[ \boxed{S = \{3\}}. \]
\textbf{Vérification} : $\ln(3+2) + \ln(3-2) = \ln 5 + \ln 1 = \ln 5 + 0 = \ln 5$. \checkmark
\end{exoresolu}

\begin{methode}[Résoudre une inéquation avec des logarithmes]
\begin{enumerate}
\item \textbf{Poser les conditions d'existence} (quantités sous $\ln$ strictement positives).
\item \textbf{Regrouper} pour obtenir une inéquation du type $\ln u \leqslant \ln v$ ou $\ln u \leqslant k$.
\item \textbf{Utiliser la stricte croissance de $\ln$} sur $\left]0\,;\,+\infty\right[$ : la fonction $\ln$ \cle{conserve l'ordre} :
\[ \ln u \leqslant \ln v \iff u \leqslant v \qquad \text{et} \qquad \ln u \leqslant k \iff u \leqslant e^{k}. \]
\item \textbf{Résoudre} puis \textbf{intersecter} avec l'ensemble des conditions. Conclure par un intervalle.
\end{enumerate}
\textbf{Réflexes à retenir} :
\[ \ln x > 0 \iff x > 1 \ ; \qquad \ln x < 0 \iff 0 < x < 1. \]
\end{methode}

\begin{exoresolu}[Inéquation avec logarithmes (type Bac A)]
Résoudre dans $\mathbb{R}$ l'inéquation :
\[ \ln(3x - 1) \leqslant \ln(x + 3). \]
\tcblower
\textbf{Étape 1 : conditions d'existence.}
Il faut $3x - 1 > 0$ et $x + 3 > 0$, c'est-à-dire $x > \dfrac{1}{3}$ et $x > -3$. La condition la plus forte est $x > \dfrac{1}{3}$, donc
\[ D = \left]\frac{1}{3}\,;\,+\infty\right[. \]

\textbf{Étape 2 : utilisation de la croissance de $\ln$.}
La fonction $\ln$ est strictement croissante sur $\left]0\,;\,+\infty\right[$, donc pour $x \in D$ :
\[ \ln(3x-1) \leqslant \ln(x+3) \iff 3x - 1 \leqslant x + 3. \]

\textbf{Étape 3 : résolution.}
\[ 3x - 1 \leqslant x + 3 \iff 2x \leqslant 4 \iff x \leqslant 2. \]

\textbf{Étape 4 : intersection avec $D$.}
On doit avoir à la fois $x > \dfrac{1}{3}$ et $x \leqslant 2$.

\textbf{Conclusion.}
\[ \boxed{S = \left]\frac{1}{3}\,;\,2\right]}. \]
Remarque : la borne $\dfrac{1}{3}$ est exclue (condition d'existence stricte), la borne $2$ est incluse (inégalité large).
\end{exoresolu}

\begin{methode}[Poser l'inconnue auxiliaire $X = \ln x$]
Quand une équation ou inéquation contient $(\ln x)^2$ et $\ln x$, on procède par \cle{changement d'inconnue} :
\begin{enumerate}
\item \textbf{Condition} : $x > 0$ (pour que $\ln x$ existe).
\item \textbf{Poser $X = \ln x$}. L'équation devient une équation polynomiale en $X$ (souvent du second degré).
\item \textbf{Résoudre en $X$} : discriminant, factorisation, tableau de signes selon le cas.
\item \textbf{Revenir à $x$} grâce à l'équivalence : $\ln x = X \iff x = e^{X}$. Pour une inéquation, utiliser la croissance de $\ln$ : $\ln x > X_0 \iff x > e^{X_0}$.
\item \textbf{Conclure} sur $x$, en vérifiant la condition $x > 0$.
\end{enumerate}
\end{methode}

\begin{exoresolu}[Inéquation du second degré en $\ln x$]
Résoudre dans $\mathbb{R}$ l'inéquation :
\[ (\ln x)^2 - \ln x - 2 \geqslant 0. \]
\tcblower
\textbf{Étape 1 : condition d'existence.} Il faut $x > 0$, donc $D = \left]0\,;\,+\infty\right[$.

\textbf{Étape 2 : changement d'inconnue.} Posons $X = \ln x$. L'inéquation devient :
\[ X^2 - X - 2 \geqslant 0. \]

\textbf{Étape 3 : résolution en $X$.} Cherchons les racines du trinôme $X^2 - X - 2$ :
\[ \Delta = (-1)^2 - 4 \times 1 \times (-2) = 1 + 8 = 9 > 0, \qquad X_1 = \frac{1 - 3}{2} = -1, \quad X_2 = \frac{1+3}{2} = 2. \]
Donc $X^2 - X - 2 = (X+1)(X-2)$. Un trinôme est du signe de $a = 1 > 0$ à l'extérieur des racines :
\[ X^2 - X - 2 \geqslant 0 \iff X \leqslant -1 \ \text{ou}\ X \geqslant 2. \]

\textbf{Étape 4 : retour à $x$.}
\begin{itemize}
\item $\ln x \leqslant -1 \iff x \leqslant e^{-1}$ (croissance de $\ln$), soit $x \leqslant \dfrac{1}{e}$ ;
\item $\ln x \geqslant 2 \iff x \geqslant e^{2}$.
\end{itemize}

\textbf{Étape 5 : conclusion} (en n'oubliant pas $x > 0$) :
\[ \boxed{S = \left]0\,;\,\frac{1}{e}\right] \cup \left[e^{2}\,;\,+\infty\right[}. \]
\end{exoresolu}

\begin{exoresolu}[Modélisation de la croissance d'une population (Contexte Cameroun)]
La population d'une localité rurale électrifiée suit le modèle $P(t) = 5000\,e^{0,03t}$ où $t$ est le nombre d'années après 2020 et $P(t)$ le nombre d'habitants. On cherche la date à partir de laquelle la population dépassera 10\,000 habitants.
\tcblower
On résout l'inéquation $P(t) > 10\,000$ :
\[ 5000\,e^{0,03t} > 10\,000 \iff e^{0,03t} > 2. \]
On applique $\ln$ (strictement croissante) des deux côtés :
\[ 0,03t > \ln 2 \iff t > \frac{\ln 2}{0,03} \approx \frac{0,693}{0,03} \approx 23,1. \]
La population dépassera 10\,000 habitants \textbf{au cours de l'année 2043} (2020 + 23,1).
\end{exoresolu}

\coursec{Limites usuelles et croissances comparées}

\begin{propriete}[Limites usuelles du logarithme népérien]
\begin{align*}
\lim_{x \to 0^{+}} \ln x &= -\infty \quad &\text{(asymptote verticale $x=0$)} \\
\lim_{x \to +\infty} \ln x &= +\infty \\
\lim_{x \to +\infty} \frac{\ln x}{x} &= 0 \quad &\text{(croissance comparée : $x$ domine $\ln x$)} \\
\lim_{x \to 0^{+}} x\ln x &= 0 \\
\lim_{x \to 0} \frac{\ln(1+x)}{x} &= 1 \quad &\text{(équivalent : $\ln(1+x) \sim x$ en $0$)}
\end{align*}
\end{propriete}

\begin{propriete}[Croissances comparées : $\ln x$ face aux puissances de $x$]
Pour tout $\alpha > 0$ :
\[ \lim_{x \to +\infty} \frac{\ln x}{x^{\alpha}} = 0 \quad \text{et} \quad \lim_{x \to 0^{+}} x^{\alpha}\ln x = 0. \]
En $+\infty$, toute puissance positive de $x$ (même $x^{0,001}$) croît infiniment plus vite que $\ln x$.
En $0^{+}$, toute puissance positive de $x$ tend vers $0$ plus vite que $\ln x$ ne tend vers $-\infty$.
\end{propriete}

\begin{methode}[Lever une indétermination avec les limites usuelles de ln]
Démarche face à une forme indéterminée ($+\infty - \infty$, $0 \times \infty$, $\frac{\infty}{\infty}$, $\frac{0}{0}$) :
\begin{enumerate}
\item \textbf{Identifier la forme} indéterminée.
\item \textbf{Factoriser par le terme dominant} (souvent $x$ ou la plus grande puissance) pour faire apparaître $\frac{\ln x}{x}$ ou $x\ln x$.
\item \textbf{Appliquer les croissances comparées} : en $+\infty$, $x^{\alpha}$ l'emporte sur $\ln x$ ; en $0^{+}$, $x^{\alpha}$ l'emporte sur $\ln x$.
\item \textbf{Conclure} par les théorèmes sur les opérations et la composition des limites.
\end{enumerate}
\end{methode}

\begin{exoresolu}[Deux limites classiques]
Calculer les limites suivantes :
\[ \text{a) } \lim_{x \to +\infty} \bigl(x - \ln x\bigr) \ ; \qquad \text{b) } \lim_{x \to +\infty} \frac{\ln x}{x^2}. \]
\tcblower
\textbf{a) Limite de $x - \ln x$ en $+\infty$.}

On a $\displaystyle\lim_{x \to +\infty} x = +\infty$ et $\displaystyle\lim_{x \to +\infty} \ln x = +\infty$ : c'est une forme indéterminée « $+\infty - \infty$ ». On factorise par le terme dominant $x$ :
\[ x - \ln x = x\left(1 - \frac{\ln x}{x}\right). \]
Par croissances comparées, $\displaystyle\lim_{x \to +\infty} \frac{\ln x}{x} = 0$, donc $\displaystyle\lim_{x \to +\infty}\left(1 - \frac{\ln x}{x}\right) = 1$. Par produit :
\[ \boxed{\lim_{x \to +\infty}\bigl(x - \ln x\bigr) = +\infty}. \]

\textbf{b) Limite de $\dfrac{\ln x}{x^2}$ en $+\infty$.}

C'est une forme « $\dfrac{\infty}{\infty}$ ». On écrit :
\[ \frac{\ln x}{x^2} = \frac{\ln x}{x} \times \frac{1}{x}. \]
Or $\displaystyle\lim_{x \to +\infty} \frac{\ln x}{x} = 0$ (croissances comparées) et $\displaystyle\lim_{x \to +\infty} \frac{1}{x} = 0$. Par produit :
\[ \boxed{\lim_{x \to +\infty} \frac{\ln x}{x^2} = 0}. \]
Interprétation : en $+\infty$, $x^2$ croît infiniment plus vite que $\ln x$.
\end{exoresolu}

\begin{exoresolu}[Limite en $0^{+}$ : forme $0 \times (-\infty)$]
Calculer $\displaystyle\lim_{x \to 0^{+}} x^{2}\ln x$.
\tcblower
C'est une forme indéterminée $0 \times (-\infty)$. On utilise la croissance comparée en $0^{+}$ : pour $\alpha = 2 > 0$, $\lim_{x \to 0^{+}} x^{2}\ln x = 0$.
\[ \boxed{\lim_{x \to 0^{+}} x^{2}\ln x = 0}. \]
On peut aussi poser $X = \frac{1}{x} \to +\infty$ : $x^{2}\ln x = \frac{-\ln X}{X^{2}} \to 0$.
\end{exoresolu}

\coursec{Dérivation et primitives liées à ln}

\begin{propriete}[Dérivée de $\ln u$ et de $\ln(ax+b)$]
Soit $u$ une fonction dérivable et strictement positive sur un intervalle $I$.
\[ (\ln u)' = \frac{u'}{u} \quad \text{sur } I. \]
En particulier, pour $a \neq 0$ et $ax+b > 0$ :
\[ \bigl(\ln(ax+b)\bigr)' = \frac{a}{ax+b}. \]
\textbf{Démonstration exigible au Bac :} Par dérivation de la composée $f = \ln \circ u$, $(f \circ u)' = f'(u) \cdot u' = \frac{1}{u} \cdot u' = \frac{u'}{u}$.
\end{propriete}

\begin{methode}[Dériver les fonctions contenant un logarithme]
Formules à connaître par cœur :
\[ (\ln x)' = \frac{1}{x} \quad (x>0) \ ; \qquad \bigl(\ln u\bigr)' = \frac{u'}{u} \ ; \qquad \bigl(\ln(ax+b)\bigr)' = \frac{a}{ax+b} \quad \left(x > -\frac{b}{a}\right). \]
Démarche :
\begin{enumerate}
\item \textbf{Identifier la structure} : somme, produit $u \times v$, quotient $\dfrac{u}{v}$, ou composée $\ln u$.
\item \textbf{Appliquer la bonne formule} : $(uv)' = u'v + uv'$ ; $\left(\dfrac{u}{v}\right)' = \dfrac{u'v - uv'}{v^2}$ ; $(\ln u)' = \dfrac{u'}{u}$.
\item \textbf{Simplifier} le résultat (mettre au même dénominateur, factoriser) : une dérivée factorisée rend l'étude du signe immédiate.
\item \textbf{Préciser l'ensemble de dérivabilité} (là où la fonction sous $\ln$ est strictement positive).
\end{enumerate}
\end{methode}

\begin{exoresolu}[Dérivée d'un produit et d'une composée]
Soit $f$ la fonction définie par $f(x) = x\ln(2x+1) - x$.
\begin{enumerate}
\item Préciser l'ensemble de définition de $f$.
\item Calculer $f'(x)$ et l'écrire sous forme simplifiée.
\end{enumerate}
\tcblower
\textbf{1. Ensemble de définition.} Il faut $2x+1 > 0 \iff x > -\dfrac{1}{2}$. Donc
\[ D_f = \left]-\frac{1}{2}\,;\,+\infty\right[. \]

\textbf{2. Calcul de la dérivée.} $f$ est de la forme $f = uv - x$ avec $u(x) = x$ et $v(x) = \ln(2x+1)$.
\begin{itemize}
\item $u'(x) = 1$ ;
\item $v = \ln w$ avec $w(x) = 2x+1$, donc $v'(x) = \dfrac{w'(x)}{w(x)} = \dfrac{2}{2x+1}$.
\end{itemize}
Par la formule du produit $(uv)' = u'v + uv'$ :
\begin{align*}
f'(x) &= 1 \times \ln(2x+1) + x \times \frac{2}{2x+1} - 1 \\
&= \ln(2x+1) + \frac{2x}{2x+1} - 1.
\end{align*}
Réduisons au même dénominateur $2x+1$ :
\[ f'(x) = \ln(2x+1) + \frac{2x - (2x+1)}{2x+1} = \ln(2x+1) - \frac{1}{2x+1}. \]

\textbf{Conclusion.} Pour tout $x \in \left]-\dfrac{1}{2}\,;\,+\infty\right[$ :
\[ \boxed{f'(x) = \ln(2x+1) - \frac{1}{2x+1}}. \]
\end{exoresolu}

\begin{propriete}[Primitives de la forme $\frac{u'}{u}$ et $\frac{a}{bx+c}$]
Si $u$ est strictement positive et dérivable sur un intervalle $I$, alors une primitive de $\dfrac{u'}{u}$ sur $I$ est $\ln u$ :
\[ \int \frac{u'(x)}{u(x)}\,\mathrm{d}x = \ln\bigl(u(x)\bigr) + k, \quad k \in \mathbb{R}. \]
En particulier, pour $a,b \in \mathbb{R}$, $b \neq 0$ et $bx+c > 0$ :
\[ \int \frac{a}{bx+c}\,\mathrm{d}x = \frac{a}{b}\,\ln(bx+c) + k. \]
\end{propriete}

\begin{methode}[Primitives liées au logarithme]
Démarche :
\begin{enumerate}
\item \textbf{Faire apparaître au numérateur la dérivée du dénominateur} : la dérivée de $bx+c$ est $b$ ; on écrit $a = \dfrac{a}{b} \times b$.
\item \textbf{Sortir la constante} et appliquer la formule $\displaystyle\int \frac{u'}{u} = \ln u + k$.
\item \textbf{Vérifier par dérivation} (réflexe anti-erreur) : dériver la primitive trouvée doit redonner la fonction de départ.
\item Pour une \textbf{primitive particulière} vérifiant $F(x_0) = y_0$, déterminer $k$ avec cette condition.
\end{enumerate}
\end{methode}

\begin{exoresolu}[Primitive particulière de $\dfrac{a}{bx+c}$]
Soit $g$ la fonction définie sur $\left]\dfrac{1}{3}\,;\,+\infty\right[$ par $g(x) = \dfrac{4}{3x-1}$.
\begin{enumerate}
\item Déterminer les primitives de $g$ sur $\left]\dfrac{1}{3}\,;\,+\infty\right[$.
\item Déterminer la primitive $G$ de $g$ qui s'annule en $x = 1$.
\end{enumerate}
\tcblower
\textbf{1. Primitives de $g$.} La dérivée du dénominateur $3x-1$ est $3$. On fait apparaître $3$ au numérateur :
\[ g(x) = \frac{4}{3x-1} = \frac{4}{3} \times \frac{3}{3x-1}. \]
On reconnaît la forme $\dfrac{u'}{u}$ avec $u(x) = 3x - 1$ (qui est bien strictement positive sur l'intervalle). Donc :
\[ \boxed{G(x) = \frac{4}{3}\ln(3x-1) + k, \quad k \in \mathbb{R}}. \]
\textbf{Vérification} : $G'(x) = \dfrac{4}{3} \times \dfrac{3}{3x-1} = \dfrac{4}{3x-1} = g(x)$. \checkmark

\textbf{2. Primitive qui s'annule en 1.} On impose $G(1) = 0$ :
\[ G(1) = \frac{4}{3}\ln(3 \times 1 - 1) + k = \frac{4}{3}\ln 2 + k = 0 \iff k = -\frac{4}{3}\ln 2. \]
\textbf{Conclusion.} La primitive cherchée est :
\[ \boxed{G(x) = \frac{4}{3}\ln(3x-1) - \frac{4}{3}\ln 2 = \frac{4}{3}\ln\!\left(\frac{3x-1}{2}\right)}. \]
\end{exoresolu}

\begin{exoresolu}[Calcul d'aire : intégrale de $\ln x$]
Calculer l'aire du domaine plan borné par la courbe $y = \ln x$, l'axe des abscisses et les droites $x=1$ et $x=e$.
\tcblower
L'aire cherchée est $\mathcal{A} = \int_{1}^{e} \ln x\,\mathrm{d}x$.
On utilise une intégration par parties : $u = \ln x \Rightarrow u' = 1/x$ ; $v' = 1 \Rightarrow v = x$.
\[ \int \ln x\,\mathrm{d}x = x\ln x - \int x \cdot \frac{1}{x}\,\mathrm{d}x = x\ln x - x + k. \]
Donc :
\[ \mathcal{A} = \Bigl[ x\ln x - x \Bigr]_{1}^{e} = (e\ln e - e) - (1\ln 1 - 1) = (e - e) - (0 - 1) = 1. \]
L'aire vaut \textbf{1 unité d'aire}.
\end{exoresolu}

\coursec{Étude complète de fonctions du type $x \mapsto \ln(ax+b)$}

\begin{methode}[Plan d'étude de $f(x) = \ln(ax+b)$]
\begin{enumerate}
\item \textbf{Ensemble de définition} : résoudre $ax + b > 0$.
\begin{itemize}
\item Si $a > 0$ : $D_f = \left]-\dfrac{b}{a}\,;\,+\infty\right[$.
\item Si $a < 0$ : $D_f = \left]-\infty\,;\,-\dfrac{b}{a}\right[$.
\end{itemize}
\item \textbf{Limites aux bornes} : utiliser $\displaystyle\lim_{X \to 0^+} \ln X = -\infty$ et $\displaystyle\lim_{X \to +\infty} \ln X = +\infty$ (composition avec $X = ax+b$). En $x = -\dfrac{b}{a}$, la limite $-\infty$ donne une \cle{asymptote verticale} d'équation $x = -\dfrac{b}{a}$.
\item \textbf{Dérivée} : $f'(x) = \dfrac{a}{ax+b}$. Son signe est celui de $a$ (le dénominateur est strictement positif sur $D_f$) :
\begin{itemize}
\item Si $a > 0$, $f$ est strictement croissante.
\item Si $a < 0$, $f$ est strictement décroissante.
\end{itemize}
\item \textbf{Tableau de variations} complet (bornes, limites, flèches).
\item \textbf{Points remarquables} :
\begin{itemize}
\item Intersection avec l'axe des abscisses : $\ln(ax+b) = 0 \iff ax+b = 1 \iff x = \dfrac{1-b}{a}$.
\item Intersection avec l'axe des ordonnées si $0 \in D_f$ : $f(0) = \ln b$ (si $b > 0$).
\end{itemize}
\item \textbf{Tracé} : placer l'asymptote, les points d'intersection, quelques points auxiliaires, puis la courbe en respectant les variations.
\end{enumerate}
\end{methode}

\begin{exoresolu}[Étude complète de $f(x) = \ln(2x - 4)$]
On considère la fonction $f$ définie par $f(x) = \ln(2x-4)$ et on note $(\mathcal{C}_f)$ sa courbe dans un repère orthonormé.
\begin{enumerate}
\item Déterminer l'ensemble de définition de $f$.
\item Calculer les limites de $f$ aux bornes de son ensemble de définition et interpréter graphiquement.
\item Étudier les variations de $f$ et dresser son tableau de variations.
\item Déterminer les coordonnées du point d'intersection de $(\mathcal{C}_f)$ avec l'axe des abscisses.
\item Tracer $(\mathcal{C}_f)$.
\end{enumerate}
\tcblower
\textbf{1. Ensemble de définition.} $2x - 4 > 0 \iff x > 2$. Donc $D_f = \left]2\,;\,+\infty\right[$.

\textbf{2. Limites aux bornes.}
\begin{itemize}
\item En $2^+$ : quand $x \to 2^+$, $2x - 4 \to 0^+$, et $\displaystyle\lim_{X \to 0^+} \ln X = -\infty$, donc $\displaystyle\lim_{x \to 2^+} f(x) = -\infty$. La droite d'équation $\boxed{x = 2}$ est \textbf{asymptote verticale} à $(\mathcal{C}_f)$.
\item En $+\infty$ : $2x - 4 \to +\infty$ et $\displaystyle\lim_{X \to +\infty} \ln X = +\infty$, donc $\displaystyle\lim_{x \to +\infty} f(x) = +\infty$.
\end{itemize}

\textbf{3. Variations.} $f$ est dérivable sur $D_f$ et, par la formule $\bigl(\ln(ax+b)\bigr)' = \dfrac{a}{ax+b}$ :
\[ f'(x) = \frac{2}{2x-4}. \]
Sur $D_f$, $2x - 4 > 0$, donc $f'(x) > 0$ : $f$ est \textbf{strictement croissante} sur $\left]2\,;\,+\infty\right[$.

\begin{center}
\begin{tabular}{|c|ccc|}
\hline
\rowcolor{popblueL} $x$ & $2$ & \hspace{2.5cm} & $+\infty$ \\
\hline
$f'(x)$ & \multicolumn{3}{c|}{$+$} \\
\hline
$f(x)$ & $-\infty$ & $\nearrow$ & $+\infty$ \\
\hline
\end{tabular}
\end{center}

\textbf{4. Intersection avec l'axe des abscisses.}
\[ f(x) = 0 \iff \ln(2x-4) = 0 \iff 2x - 4 = e^{0} = 1 \iff x = \frac{5}{2}. \]
La courbe coupe l'axe des abscisses au point $\boxed{A\left(\dfrac{5}{2}\,;\,0\right)}$.

\textbf{5. Tracé.} Points auxiliaires : $f(4) = \ln 4 \approx 1{,}39$ ; $f(6) = \ln 8 \approx 2{,}08$.

\begin{popfigure}
\begin{center}
\begin{tikzpicture}[scale=0.85]
\draw[->,popdark] (-0.5,0) -- (7.5,0) node[below left]{$x$};
\draw[->,popdark] (0,-3) -- (0,3) node[below left]{$y$};
\draw[popmuted,dashed,thick] (2,-3) -- (2,3) node[above right,popmuted]{$x=2$};
\draw[popblue,very thick,domain=2.06:7.2,samples=200] plot (\x,{ln(2*\x-4)});
\fill[popink] (2.5,0) circle (2pt) node[below right,popink]{$A\left(\frac{5}{2}\,;\,0\right)$};
\fill[popink] (4,{ln 4}) circle (2pt);
\fill[popink] (6,{ln 8}) circle (2pt);
\foreach \i in {1,3,4,5,6,7} \draw (\i,0.06) -- (\i,-0.06) node[below,popdark]{\small $\i$};
\foreach \j in {-2,-1,1,2} \draw (0.06,\j) -- (-0.06,\j) node[left,popdark]{\small $\j$};
\end{tikzpicture}
\end{center}
\legende{Courbe de $f(x) = \ln(2x-4)$ : asymptote verticale $x = 2$, croissante, passant par $A\left(\frac{5}{2}\,;\,0\right)$.}
\end{popfigure}

\textbf{Conclusion.} $f$ est strictement croissante de $-\infty$ vers $+\infty$ sur $\left]2\,;\,+\infty\right[$, avec asymptote verticale $x = 2$ ; elle s'annule en $x = \dfrac{5}{2}$.
\end{exoresolu}

\begin{aretenir}[Cas $a < 0$ : symétrie et décroissance]
Si $a < 0$, par exemple $f(x) = \ln(4 - x)$ ($a=-1, b=4$) :
\begin{itemize}
\item $D_f = \left]-\infty\,;\,4\right[$.
\item Asymptote verticale $x = 4$ (limite $-\infty$ en $4^{-}$).
\item Limite $+\infty$ en $-\infty$.
\item $f'(x) = \frac{-1}{4-x} < 0$ : $f$ est \textbf{strictement décroissante}.
\item Tableau de variations : $-\infty \xrightarrow{\searrow} +\infty$ sur $]-\infty; 4[$.
\item La courbe est l'image de $x \mapsto \ln(x-4)$ (définie sur $]4,+\infty[$) par symétrie par rapport à la droite $x=4$.
\end{itemize}
\end{aretenir}

\begin{attention}[Les 5 erreurs qui coûtent des points au Bac]
\begin{enumerate}
\item \textbf{Oublier les conditions d'existence.} Avant toute résolution, écris que les quantités sous chaque $\ln$ sont strictement positives, et \emph{rejette} à la fin les solutions hors de l'ensemble de résolution. Une solution comme $x = -3$ pour $\ln(x+2) + \ln(x-2) = \ln 5$ doit être explicitement écartée.
\item \textbf{Écrire $\ln(a+b) = \ln a + \ln b$.} C'est faux ! Les propriétés algébriques portent sur $\ln(ab)$, $\ln\!\left(\dfrac{a}{b}\right)$ et $\ln(a^n)$. De même, $(\ln x)^2$ n'est pas $\ln(x^2)$ : on a $\ln(x^2) = 2\ln x$ seulement si $x > 0$.
\item \textbf{Changer le sens de l'inégalité.} La fonction $\ln$ est \emph{strictement croissante} : elle conserve l'ordre. On n'inverse jamais le sens en « passant le $\ln$ de l'autre côté ». (C'est avec des fonctions décroissantes comme $x \mapsto \left(\frac{1}{2}\right)^x$ qu'on inverse.)
\item \textbf{Oublier le coefficient dans la dérivation ou la primitivation.} $\bigl(\ln(ax+b)\bigr)' = \dfrac{a}{ax+b}$ et non $\dfrac{1}{ax+b}$ ; réciproquement $\displaystyle\int \frac{a}{bx+c}\,\mathrm{d}x = \dfrac{a}{b}\ln(bx+c) + k$ et non $a\ln(bx+c)$. Vérifie toujours ta primitive en la dérivant.
\item \textbf{Annoncer une limite fausse ou une forme indéterminée non traitée.} $\displaystyle\lim_{x \to 0^+} \ln x = -\infty$ (et non $0$ !) ; face à « $+\infty - \infty$ » ou « $\dfrac{\infty}{\infty}$ », il est interdit de conclure directement : factorise par le terme dominant et utilise $\displaystyle\lim_{x \to +\infty} \frac{\ln x}{x} = 0$.
\end{enumerate}
\end{attention}

\newpage

\coursec{Exercices}

\coursec{Fonction logarithme népérien}

\courssub{Je vérifie mes connaissances}

\begin{exercice}[QCM : une seule réponse exacte]
Pour chaque question, choisis la bonne réponse en justifiant brièvement.
\begin{enumerate}
\item L'ensemble de définition de la fonction $f : x \longmapsto \ln(2x-6)$ est :
\begin{enumerate}
\item[a)] $]0\,;\,+\infty[$ \quad b) $]3\,;\,+\infty[$ \quad c) $[3\,;\,+\infty[$ \quad d) $]-\infty\,;\,3[$
\end{enumerate}
\item Pour tout réel $a > 0$, $\ln(a^3) - \ln(a^2)$ est égal à :
\begin{enumerate}
\item[a)] $\ln a$ \quad b) $\ln(a^5)$ \quad c) $\dfrac{3}{2}\ln a$ \quad d) $\ln(a^6)$
\end{enumerate}
\item $\displaystyle\lim_{x \to +\infty} \dfrac{\ln x}{x}$ est égale à :
\begin{enumerate}
\item[a)] $+\infty$ \quad b) $1$ \quad c) $0$ \quad d) n'existe pas
\end{enumerate}
\item La dérivée de $g : x \longmapsto \ln(5x+1)$ sur $\left]-\dfrac{1}{5}\,;\,+\infty\right[$ est :
\begin{enumerate}
\item[a)] $g'(x) = \dfrac{1}{5x+1}$ \quad b) $g'(x) = \dfrac{5}{5x+1}$ \quad c) $g'(x) = \dfrac{5}{x}$ \quad d) $g'(x) = \ln 5$
\end{enumerate}
\end{enumerate}
\end{exercice}

\begin{exercice}[Vrai ou faux, en justifiant]
Précise si chacune des affirmations suivantes est vraie ou fausse, en justifiant ta réponse par une propriété du cours ou un contre-exemple.
\begin{enumerate}
\item Pour tous réels strictement positifs $a$ et $b$ : $\ln(a+b) = \ln a + \ln b$.
\item Pour tout réel $x > 1$ : $\ln(x^2 - 1) = \ln(x-1) + \ln(x+1)$.
\item $\ln 1 = 0$ et $\ln e = 1$.
\item Si $\ln x < \ln y$, alors $x < y$.
\item $\displaystyle\lim_{x \to 0^+} x \ln x = +\infty$.
\end{enumerate}
\end{exercice}

\begin{exercice}[Calculs directs]
Sans calculatrice, simplifie les expressions suivantes (les lettres désignent des réels strictement positifs) :
\begin{enumerate}
\item $A = \ln 12 + \ln 3 - \ln 9$ ;
\item $B = 2\ln 5 + \ln 4 - \ln 10$ ;
\item $C = \ln \sqrt{7} + \dfrac{1}{2}\ln 7$ ;
\item $D = \ln\left(\dfrac{a^3}{b^2}\right) - 3\ln a + \ln(b^2)$.
\end{enumerate}
\end{exercice}

\begin{exercice}[Définitions et réflexes de cours]
\begin{enumerate}
\item Donne la définition de la fonction logarithme népérien comme réciproque d'une fonction connue, et précise son ensemble de définition.
\item Complète : pour tous réels $a > 0$ et $b > 0$, $\ln(ab) = \ldots$ ; $\ln\left(\dfrac{a}{b}\right) = \ldots$ ; $\ln(a^n) = \ldots$ pour $n \in \mathbb{Z}$.
\item Donne une primitive de la fonction $x \longmapsto \dfrac{4}{3x-2}$ sur l'intervalle $\left]\dfrac{2}{3}\,;\,+\infty\right[$.
\item Résous dans $\mathbb{R}$ l'équation $\ln x = 3$, puis l'inéquation $\ln x \geq -1$.
\end{enumerate}
\end{exercice}

\courssub{Je m'entraîne}

\begin{methode}[Résolution d'une équation ou inéquation logarithmique]
\textbf{Étapes incontournables :}
\begin{enumerate}
\item Déterminer l'\textbf{ensemble de validité} (arguments des ln strictement positifs).
\item Utiliser les propriétés algébriques pour \textbf{réduire à une forme simple} : $\ln(\text{expr}_1) = \ln(\text{expr}_2)$ ou $\ln(\text{expr}) = k$.
\item Résoudre l'équation / inéquation \textbf{sans les ln} (produit, quotient, exponentielle).
\item \textbf{Sélectionner} les solutions appartenant à l'ensemble de validité.
\item Conclure par une phrase : « L'ensemble des solutions est... »
\end{enumerate}
\end{methode}

\begin{exercice}[Équations avec logarithmes]
Résous dans $\mathbb{R}$ chacune des équations suivantes, après avoir déterminé l'ensemble de validité :
\begin{enumerate}
\item $\ln(3x - 1) = \ln(x + 5)$ ;
\item $\ln(x+2) + \ln(x-1) = \ln(4x - 4)$ ;
\item $\ln(x^2) = 2\ln(3x - 2)$ ;
\item $\ln(x-1) + \ln(x+1) = \ln 3$.
\end{enumerate}
\end{exercice}

\begin{methode}[Inconnue auxiliaire $X = \ln x$]
Lorsque l'équation ou l'inéquation ne contient que $\ln x$ et $(\ln x)^2$ (ou puissances supérieures) :
\begin{enumerate}
\item Poser $X = \ln x$ (donc $x = e^X$, $x>0$).
\item Résoudre l'équation / inéquation en $X$ dans $\mathbb{R}$.
\item Revenir à $x$ par $x = e^X$ (toujours $>0$).
\end{enumerate}
\end{methode}

\begin{exercice}[Inconnue auxiliaire $X = \ln x$]
\begin{enumerate}
\item Résous dans $\mathbb{R}$ l'équation $(\ln x)^2 - 3\ln x + 2 = 0$. \emph{Indication : pose $X = \ln x$.}
\item Résous dans $\mathbb{R}$ l'inéquation $(\ln x)^2 - 3\ln x + 2 < 0$.
\item Résous le système d'inconnues $(x\,;\,y) \in \,]0\,;\,+\infty[^2$ :
\[
\begin{cases}
\ln x + \ln y = 3\ln 2 \\
(\ln x)(\ln y) = 2(\ln 2)^2
\end{cases}
\]
\emph{Indication : pose $X = \ln x$ et $Y = \ln y$, puis utilise la somme et le produit des racines d'un trinôme.}
\end{enumerate}
\end{exercice}

\begin{exercice}[Calcul de limites]
Calcule les limites suivantes en justifiant chaque étape (pense aux limites usuelles et aux croissances comparées) :
\begin{enumerate}
\item $\displaystyle\lim_{x \to +\infty} \dfrac{\ln x - x}{x + 1}$ ;
\item $\displaystyle\lim_{x \to 0^+} x\,\ln(x^2)$ ;
\item $\displaystyle\lim_{x \to +\infty} \dfrac{\ln(2x+3)}{x}$ ;
\item $\displaystyle\lim_{x \to +\infty} \left(\ln x - \sqrt{x}\right)$.
\end{enumerate}
\end{exercice}

\begin{exercice}[Dérivation et primitives]
\begin{enumerate}
\item Détermine la dérivée de chacune des fonctions suivantes, en précisant l'ensemble de dérivabilité :
\[
f(x) = \ln(4x - 8) \ ; \qquad g(x) = \ln(x^2 + 1) \ ; \qquad h(x) = x\ln x - x.
\]
\item Détermine les primitives de $f : x \longmapsto \dfrac{5}{2x-3}$ sur l'intervalle $\left]\dfrac{3}{2}\,;\,+\infty\right[$.
\item Détermine la primitive $F$ de $f$ sur cet intervalle qui s'annule en $x = 2$.
\end{enumerate}
\end{exercice}

\courssub{Je me prépare au Bac}

\begin{methode}[Étude complète d'une fonction $f(x) = \ln(ax+b)$ ($a \neq 0$)]
\begin{enumerate}
\item \textbf{Domaine :} $ax+b > 0 \iff x > -b/a$ (si $a>0$) ou $x < -b/a$ (si $a<0$).
\item \textbf{Limites :} Aux bornes du domaine, utiliser $\lim_{u\to 0^+} \ln u = -\infty$ et $\lim_{u\to +\infty} \ln u = +\infty$. Asymptote verticale à la borne finie.
\item \textbf{Dérivée :} $f'(x) = \dfrac{a}{ax+b}$. Signe constant (signe de $a$) sur le domaine.
\item \textbf{Variations :} Strictement croissante si $a>0$, décroissante si $a<0$.
\item \textbf{Courbe :} Passe par $(-b/a + 1/a, 0)$ car $\ln 1 = 0$. Tangente au point d'abscisse $x_0$ : $y = f'(x_0)(x-x_0) + f(x_0)$.
\end{enumerate}
\end{methode}

\begin{exercice}[Étude complète d'une fonction logarithme]
On considère la fonction $f$ définie par $f(x) = \ln(3x - 6)$, et on note $(\mathcal{C}_f)$ sa courbe représentative dans un repère orthonormé $(O\,;\,\vec{\imath}\,,\,\vec{\jmath})$ d'unité graphique \SI{1}{cm}.
\begin{enumerate}
\item Détermine l'ensemble de définition $D_f$ de $f$.
\item Calcule les limites de $f$ aux bornes de $D_f$. Interprète graphiquement le résultat obtenu en $2$.
\item Calcule la dérivée $f'$ de $f$ et étudie son signe sur $D_f$.
\item Dresse le tableau de variation complet de $f$.
\item Détermine une équation de la tangente $(T)$ à $(\mathcal{C}_f)$ au point d'abscisse $\dfrac{7}{3}$.
\item Détermine les coordonnées du point d'intersection de $(\mathcal{C}_f)$ avec l'axe des abscisses.
\item Trace soigneusement $(\mathcal{C}_f)$, son asymptote et la tangente $(T)$ dans le repère.
\item Résous graphiquement, puis algébriquement, l'inéquation $f(x) \geq 0$.
\end{enumerate}
\end{exercice}

\begin{savaistu}[Modélisation au Cameroun : Démographie et Bactériologie]
La croissance exponentielle modélise de nombreux phénomènes au Cameroun : démographie (population de Yaoundé ou Douala), épidémiologie (propagation d'un virus), finance (intérêts composés en microfinance), ou biologie (cultures bactériennes à l'Université de Yaoundé I ou à l'IRAD). Le logarithme népérien permet de « linéariser » ces croissances pour en extraire des paramètres (temps de doublement, taux continu) et faciliter les prévisions.
\end{savaistu}

\begin{exercice}[Croissance d'une population de bactéries]
Un laboratoire de l'Université de Yaoundé I étudie une culture de bactéries. On modélise le nombre de bactéries à l'instant $t$ (exprimé en heures) par :
\[
P(t) = 1000\,e^{0,5t}, \qquad t \geq 0.
\]
\textbf{Partie A : lecture directe du modèle}
\begin{enumerate}
\item Quel est le nombre initial de bactéries ?
\item Calcule le nombre de bactéries au bout de $4$ heures (on donnera une valeur approchée à l'unité près ; on rappelle que $e^2 \approx 7,39$).
\item Détermine, en résolvant une inéquation, au bout de combien de temps la population dépasse \num{100000} individus. Donne la valeur exacte, puis une valeur approchée à $0,1$ h près (on donne $\ln 100 \approx 4,6$).
\end{enumerate}
\textbf{Partie B : étude de la fonction auxiliaire}
On considère la fonction $h$ définie sur $[0\,;\,+\infty[$ par $h(t) = \ln\bigl(P(t)\bigr)$.
\begin{enumerate}
\item Montre que pour tout $t \geq 0$, $h(t) = 0,5t + \ln 1000$.
\item Quelle est la nature de la fonction $h$ ? Interprète ce résultat : pourquoi dit-on que le logarithme « linéarise » une croissance exponentielle ?
\item Résous l'équation $h(t) = \ln(100000)$ et retrouve le résultat de la question A.3.
\end{enumerate}
\end{exercice}

\begin{aretenir}[Inégalité fondamentale : $\ln x \leq x - 1$]
Pour tout réel $x > 0$ : $\ln x \leq x - 1$, avec égalité \textbf{si et seulement si} $x = 1$.
\textbf{Preuve standard :} Étudier $\varphi(x) = x - 1 - \ln x$ sur $]0, +\infty[$.
\textbf{Conséquence graphique :} La courbe $y = \ln x$ est \textbf{toujours en dessous} de sa tangente au point $(1,0)$ d'équation $y = x-1$, sauf au point de contact.
\textbf{Double inégalité utile :} Pour tout $x > 0$, $\quad 1 - \dfrac{1}{x} \leq \ln x \leq x - 1$.
\end{aretenir}

\begin{exercice}[Une inégalité classique et ses applications]
\textbf{Partie A : démonstration d'une inégalité de référence}

On considère la fonction $\varphi$ définie sur $]0\,;\,+\infty[$ par :
\[
\varphi(x) = x - 1 - \ln x.
\]
\begin{enumerate}
\item Calcule la dérivée $\varphi'$ de $\varphi$ et étudie son signe sur $]0\,;\,+\infty[$.
\item Dresse le tableau de variation de $\varphi$ (on calculera $\varphi(1)$).
\item En déduire que pour tout réel $x > 0$ : $\ln x \leq x - 1$. Dans quel cas a-t-on l'égalité ?
\item Interprète graphiquement cette inégalité en termes de positions relatives de la courbe de $\ln$ et d'une droite que l'on précisera.
\end{enumerate}
\textbf{Partie B : application}
\begin{enumerate}
\item En appliquant l'inégalité de la Partie A au réel $\dfrac{1}{x}$ (avec $x > 0$), démontre que pour tout $x > 0$ :
\[
1 - \dfrac{1}{x} \leq \ln x \leq x - 1.
\]
\item Encadrement d'un taux d'intérêt. Un commerçant de Douala place un capital dans une microfinance. Le capital est multiplié par $1{,}08$ en un an. Le taux continu équivalent $t$ vérifie $\ln(1{,}08) = t$. En utilisant la double inégalité de la question B.1 avec $x = 1{,}08$, donne un encadrement de $t$ par deux nombres décimaux, puis une valeur approchée de $t$ à $10^{-3}$ près.
\end{enumerate}
\end{exercice}

\newpage

\coursec{Corrigés des exercices}

\begin{corrige}[Exercice 1 — QCM : une seule réponse exacte]
\begin{enumerate}
\item \textbf{Réponse b) } $]3\,;\,+\infty[$. La fonction $\ln$ n'est définie que pour un argument strictement positif : $2x-6>0 \iff x>3$.
\item \textbf{Réponse a) } $\ln a$. En effet, pour $a>0$ :
\[ \ln(a^3)-\ln(a^2)=3\ln a-2\ln a=\ln a. \]
\item \textbf{Réponse c) } $0$. C'est une limite de croissance comparée : $\displaystyle\lim_{x\to+\infty}\dfrac{\ln x}{x}=0$.
\item \textbf{Réponse b) } $g'(x)=\dfrac{5}{5x+1}$. La dérivée de $x\mapsto\ln(ax+b)$ est $x\mapsto\dfrac{a}{ax+b}$ ; ici $a=5$, $b=1$.
\end{enumerate}
\end{corrige}

\begin{corrige}[Exercice 2 — Vrai ou faux, en justifiant]
\begin{enumerate}
\item \textbf{Faux.} Contre-exemple : $a=b=1$ donne $\ln(1+1)=\ln 2\approx 0{,}69$, alors que $\ln 1+\ln 1=0$. La propriété correcte est $\ln(ab)=\ln a+\ln b$ : le logarithme transforme un \emph{produit} en somme, pas une somme.
\item \textbf{Vrai.} Pour $x>1$, on a $x-1>0$ et $x+1>0$, donc les deux logarithmes du membre de droite existent, et :
\[ \ln(x-1)+\ln(x+1)=\ln\bigl((x-1)(x+1)\bigr)=\ln(x^2-1). \]
\item \textbf{Vrai.} Ce sont des valeurs de référence : $e^0=1$ donc $\ln 1=0$, et $e^1=e$ donc $\ln e=1$ (par définition de $\ln$ comme réciproque de l'exponentielle).
\item \textbf{Vrai.} La fonction $\ln$ est \textbf{strictement croissante} sur $]0\,;\,+\infty[$ ; elle conserve donc l'ordre : $\ln x<\ln y \iff x<y$ (pour $x,y>0$).
\item \textbf{Faux.} Par croissance comparée, $\displaystyle\lim_{x\to 0^+}x\ln x=0$ (et non $+\infty$) : $x$ « l'emporte » sur $\ln x$.
\end{enumerate}
\end{corrige}

\begin{corrige}[Exercice 3 — Calculs directs]
\begin{enumerate}
\item $A=\ln 12+\ln 3-\ln 9=\ln\left(\dfrac{12\times 3}{9}\right)=\ln 4=\ln(2^2)=2\ln 2$.
\item $B=2\ln 5+\ln 4-\ln 10=\ln 25+\ln 4-\ln 10=\ln\left(\dfrac{25\times 4}{10}\right)=\ln 10$.
\item $C=\ln\sqrt{7}+\dfrac{1}{2}\ln 7=\dfrac{1}{2}\ln 7+\dfrac{1}{2}\ln 7=\ln 7$.
\item $D=\ln\left(\dfrac{a^3}{b^2}\right)-3\ln a+\ln(b^2)=3\ln a-2\ln b-3\ln a+2\ln b=0$.
\end{enumerate}
\textbf{Vérification de D :} $\ln\left(\dfrac{a^3}{b^2}\right)+\ln(b^2)=\ln\left(\dfrac{a^3 b^2}{b^2}\right)=\ln(a^3)=3\ln a$, donc $D=3\ln a-3\ln a=0$. 
\end{corrige}

\begin{corrige}[Exercice 4 — Définitions et réflexes de cours]
\begin{enumerate}
\item La fonction \textbf{logarithme népérien}, notée $\ln$, est la \textbf{fonction réciproque} de la fonction exponentielle : pour tout $x>0$ et tout réel $y$,
\[ y=\ln x \iff x=e^y. \]
Son ensemble de définition est $]0\,;\,+\infty[$, et elle est à valeurs dans $\mathbb{R}$.
\item Pour tous $a>0$, $b>0$ et $n\in\mathbb{Z}$ :
\[ \ln(ab)=\ln a+\ln b \ ; \qquad \ln\left(\frac{a}{b}\right)=\ln a-\ln b \ ; \qquad \ln(a^n)=n\ln a. \]
\item On écrit $\dfrac{4}{3x-2}=\dfrac{4}{3}\times\dfrac{3}{3x-2}$ pour faire apparaître la forme $\dfrac{u'}{u}$ avec $u(x)=3x-2>0$ sur $\left]\dfrac{2}{3}\,;\,+\infty\right[$. Une primitive est donc :
\[ F(x)=\dfrac{4}{3}\ln(3x-2). \]
\item \textbf{Équation} $\ln x=3$ : condition $x>0$. Par définition, $x=e^3$. $\mathcal{S}=\{e^3\}$.

\textbf{Inéquation} $\ln x\geq -1$ : condition $x>0$. Par stricte croissance de l'exponentielle : $x\geq e^{-1}=\dfrac{1}{e}$. $\mathcal{S}=\left[\dfrac{1}{e}\,;\,+\infty\right[$.
\end{enumerate}
\end{corrige}

\begin{corrige}[Exercice 5 — Équations avec logarithmes]
\begin{enumerate}
\item \textbf{Validité :} $3x-1>0$ et $x+5>0$, soit $x>\dfrac{1}{3}$. Sur $\left]\dfrac{1}{3}\,;\,+\infty\right[$, par injectivité de $\ln$ :
\[ 3x-1=x+5 \iff 2x=6 \iff x=3. \]
$3>\dfrac{1}{3}$ : accepté. $\mathcal{S}=\{3\}$.
\item \textbf{Validité :} $x+2>0$, $x-1>0$ et $4x-4>0$, soit $x>1$. L'équation s'écrit :
\[ \ln\bigl((x+2)(x-1)\bigr)=\ln(4x-4) \iff x^2+x-2=4x-4 \iff x^2-3x+2=0. \]
$\Delta=9-8=1$ ; racines $x_1=\dfrac{3-1}{2}=1$ et $x_2=\dfrac{3+1}{2}=2$. Or $x=1$ est \textbf{exclu} (validité stricte $x>1$). $\mathcal{S}=\{2\}$.

\textbf{Vérification :} $\ln 4+\ln 1=\ln 4$ et $\ln(8-4)=\ln 4$. Correct.
\item \textbf{Validité :} $x^2>0$ et $3x-2>0$, soit $x>\dfrac{2}{3}$ (ce qui implique $x\neq 0$). Attention : $\ln(x^2)=2\ln|x|=2\ln x$ ici car $x>0$. L'équation devient :
\[ 2\ln x=2\ln(3x-2) \iff x=3x-2 \iff x=1. \]
$1>\dfrac{2}{3}$ : accepté. $\mathcal{S}=\{1\}$.
\item \textbf{Validité :} $x-1>0$ et $x+1>0$, soit $x>1$. L'équation s'écrit :
\[ \ln(x^2-1)=\ln 3 \iff x^2-1=3 \iff x^2=4 \iff x=\pm 2. \]
Seul $x=2$ appartient à $]1\,;\,+\infty[$. $\mathcal{S}=\{2\}$.
\end{enumerate}
\textbf{Point de vigilance :} la sélection finale des solutions dans l'ensemble de validité est \emph{indispensable} ; oublier cette étape fait perdre des points au Bac (cf. question 2 où $x=1$ est une solution parasite).
\end{corrige}

\begin{corrige}[Exercice 6 — Inconnue auxiliaire $X=\ln x$]
\begin{enumerate}
\item Condition : $x>0$. Posons $X=\ln x$. L'équation devient $X^2-3X+2=0$. $\Delta=9-8=1$ ; racines $X_1=1$ et $X_2=2$. Retour à $x$ : $x=e^1=e$ ou $x=e^2$.
\[ \mathcal{S}=\{e\,;\,e^2\}. \]
\item Le trinôme $X^2-3X+2=(X-1)(X-2)$ est strictement négatif \emph{entre} ses racines : $1<X<2$. Comme $x=e^X$ et que l'exponentielle est strictement croissante :
\[ e^1<x<e^2 \quad\Longrightarrow\quad \mathcal{S}=\left]e\,;\,e^2\right[. \]
\item Posons $X=\ln x$ et $Y=\ln y$. Le système devient :
\[ \begin{cases} X+Y=3\ln 2 \\ XY=2(\ln 2)^2 \end{cases} \]
$X$ et $Y$ sont les racines du trinôme $t^2-(X+Y)t+XY=0$, c'est-à-dire :
\[ t^2-3\ln 2\,t+2(\ln 2)^2=0. \]
$\Delta=9(\ln 2)^2-8(\ln 2)^2=(\ln 2)^2$, donc $\sqrt{\Delta}=\ln 2$ (car $\ln 2>0$). Racines :
\[ t_1=\dfrac{3\ln 2-\ln 2}{2}=\ln 2, \qquad t_2=\dfrac{3\ln 2+\ln 2}{2}=2\ln 2=\ln 4. \]
Donc $\{X,Y\}=\{\ln 2\,;\,\ln 4\}$, c'est-à-dire $\{x,y\}=\{2\,;\,4\}$. Le système étant symétrique :
\[ \mathcal{S}=\{(2\,;\,4)\,;\,(4\,;\,2)\}. \]
\textbf{Vérification :} $\ln 2+\ln 4=\ln 8=3\ln 2$ et $(\ln 2)(\ln 4)=2(\ln 2)^2$. Correct.
\end{enumerate}
\end{corrige}

\begin{corrige}[Exercice 7 — Calcul de limites]
\begin{enumerate}
\item Forme indéterminée $\dfrac{\infty}{\infty}$. Factorisons par $x$ au numérateur et au dénominateur :
\[ \dfrac{\ln x - x}{x+1}=\dfrac{x\left(\dfrac{\ln x}{x}-1\right)}{x\left(1+\dfrac{1}{x}\right)}=\dfrac{\dfrac{\ln x}{x}-1}{1+\dfrac{1}{x}}. \]
Or $\displaystyle\lim_{x\to+\infty}\dfrac{\ln x}{x}=0$ (croissance comparée) et $\displaystyle\lim_{x\to+\infty}\dfrac{1}{x}=0$. Donc la limite vaut $\dfrac{0-1}{1+0}=\boxed{-1}$.
\item $x\ln(x^2)=2x\ln x$ (car $x>0$). Par croissance comparée, $\displaystyle\lim_{x\to 0^+}x\ln x=0$, donc :
\[ \lim_{x\to 0^+}x\ln(x^2)=2\times 0=0. \]
\item Posons $u=2x+3$ : quand $x\to+\infty$, $u\to+\infty$, et $\dfrac{\ln(2x+3)}{x}=\dfrac{\ln u}{u}\times\dfrac{2x+3}{x}=\dfrac{\ln u}{u}\times\left(2+\dfrac{3}{x}\right)$. Comme $\displaystyle\lim_{u\to+\infty}\dfrac{\ln u}{u}=0$ et $2+\dfrac{3}{x}\to 2$, la limite est $0\times 2=\boxed{0}$.
\item Forme indéterminée $\infty-\infty$. Factorisons par $\sqrt{x}$ :
\[ \ln x-\sqrt{x}=\sqrt{x}\left(\dfrac{\ln x}{\sqrt{x}}-1\right). \]
Or $\dfrac{\ln x}{\sqrt{x}}=\dfrac{2\ln\sqrt{x}}{\sqrt{x}}=2\,\dfrac{\ln u}{u}$ avec $u=\sqrt{x}\to+\infty$, donc $\displaystyle\lim_{x\to+\infty}\dfrac{\ln x}{\sqrt{x}}=0$. Ainsi la parenthèse tend vers $-1$ et $\sqrt{x}\to+\infty$ :
\[ \lim_{x\to+\infty}\left(\ln x-\sqrt{x}\right)=-\infty. \]
\end{enumerate}
\end{corrige}

\begin{corrige}[Exercice 8 — Dérivation et primitives]
\begin{enumerate}
\item \textbf{Pour $f(x)=\ln(4x-8)$ :} définie et dérivable lorsque $4x-8>0$, c'est-à-dire sur $]2\,;\,+\infty[$. Forme $\ln u$ avec $u(x)=4x-8$, $u'(x)=4$ :
\[ f'(x)=\dfrac{4}{4x-8}=\dfrac{1}{x-2}. \]
\textbf{Pour $g(x)=\ln(x^2+1)$ :} $x^2+1>0$ pour tout réel $x$, donc $g$ est dérivable sur $\mathbb{R}$ :
\[ g'(x)=\dfrac{2x}{x^2+1}. \]
\textbf{Pour $h(x)=x\ln x - x$ :} dérivable sur $]0\,;\,+\infty[$. Dérivée du produit $x\ln x$ : $1\times\ln x+x\times\dfrac{1}{x}=\ln x+1$. Donc :
\[ h'(x)=\ln x+1-1=\ln x. \]
\emph{Remarque :} $h$ est donc \textbf{la} primitive de $\ln$ qui s'annule en $1$.
\item On fait apparaître la forme $\dfrac{u'}{u}$ avec $u(x)=2x-3>0$ sur $\left]\dfrac{3}{2}\,;\,+\infty\right[$ et $u'(x)=2$ :
\[ f(x)=\dfrac{5}{2x-3}=\dfrac{5}{2}\times\dfrac{2}{2x-3}. \]
Les primitives de $f$ sont donc les fonctions :
\[ F_k(x)=\dfrac{5}{2}\ln(2x-3)+k, \qquad k\in\mathbb{R}. \]
\item Condition $F_k(2)=0$ : $\dfrac{5}{2}\ln(4-3)+k=\dfrac{5}{2}\ln 1+k=k=0$. Donc $k=0$ et :
\[ F(x)=\dfrac{5}{2}\ln(2x-3). \]
\end{enumerate}
\end{corrige}

\begin{corrige}[Exercice 9 — Étude complète d'une fonction logarithme]
\emph{Barème indicatif (sur 10) : domaine 1 pt ; limites + interprétation 2 pts ; dérivée et signe 1,5 pt ; tableau 1 pt ; tangente 1,5 pt ; intersection axe 1 pt ; courbe 1 pt ; inéquation 1 pt.}

\begin{enumerate}
\item \textbf{Domaine :} $3x-6>0 \iff x>2$. Donc $D_f=\left]2\,;\,+\infty\right[$.
\item \textbf{Limites :} En $2^+$ : posons $u=3x-6$ ; quand $x\to 2^+$, $u\to 0^+$, et $\displaystyle\lim_{u\to 0^+}\ln u=-\infty$. Donc :
\[ \lim_{x\to 2^+}f(x)=-\infty. \]
\textbf{Interprétation :} la droite d'équation $x=2$ est \textbf{asymptote verticale} à $(\mathcal{C}_f)$.

En $+\infty$ : $3x-6\to+\infty$ et $\displaystyle\lim_{u\to+\infty}\ln u=+\infty$, donc $\displaystyle\lim_{x\to+\infty}f(x)=+\infty$.
\item \textbf{Dérivée :} $f$ est de la forme $\ln(ax+b)$ avec $a=3$, $b=-6$ :
\[ f'(x)=\dfrac{3}{3x-6}=\dfrac{1}{x-2}. \]
Sur $]2\,;\,+\infty[$, $x-2>0$, donc $f'(x)>0$ : $f$ est \textbf{strictement croissante} sur $D_f$.
\item \textbf{Tableau de variation :}
\begin{center}
\begin{tabular}{|c|c||cc|}\hline
$x$ & $2$ & & $+\infty$ \\ \hline
$f'(x)$ & & $+$ & \\ \hline
$f$ & $-\infty$ & $\nearrow$ & $+\infty$ \\ \hline
\end{tabular}
\end{center}
(La double barre en $2$ rappelle que $f$ n'y est pas définie.)
\item \textbf{Tangente en $x_0=\dfrac{7}{3}$ :} $f\left(\dfrac{7}{3}\right)=\ln(7-6)=\ln 1=0$ et $f'\left(\dfrac{7}{3}\right)=\dfrac{1}{\frac{7}{3}-2}=\dfrac{1}{\frac{1}{3}}=3$. Équation :
\[ y=f'(x_0)(x-x_0)+f(x_0)=3\left(x-\dfrac{7}{3}\right) \quad\text{soit}\quad (T):\ y=3x-7. \]
\item \textbf{Intersection avec l'axe des abscisses :} $f(x)=0 \iff \ln(3x-6)=0 \iff 3x-6=1 \iff x=\dfrac{7}{3}$. Le point est $\left(\dfrac{7}{3}\,;\,0\right)$ — c'est précisément le point de tangence de la question 5.
\item \textbf{Tracé :} on place l'asymptote $x=2$, le point $\left(\dfrac{7}{3}\,;\,0\right)\approx(2{,}33\,;\,0)$, la tangente $y=3x-7$ (qui coupe l'axe des ordonnées en $-7$ et passe par $(3\,;\,2)$), un point supplémentaire $f(4)=\ln 6\approx 1{,}8$, et la courbe strictement croissante, concave, s'approchant de l'asymptote en plongeant vers $-\infty$.
\item \textbf{Inéquation $f(x)\geq 0$ :}

\emph{Graphiquement :} $(\mathcal{C}_f)$ est au-dessus de l'axe des abscisses (ou le coupe) à partir de $x=\dfrac{7}{3}$.

\emph{Algébriquement :} $\ln(3x-6)\geq 0=\ln 1$. Par stricte croissance de $\ln$ : $3x-6\geq 1 \iff x\geq\dfrac{7}{3}$.
\[ \mathcal{S}=\left[\dfrac{7}{3}\,;\,+\infty\right[. \]
Les deux méthodes concordent.
\end{enumerate}
\textbf{Points de vigilance (Bac) :} citer la condition $3x-6>0$ avant tout calcul ; justifier la limite en $2^+$ par composition ($u=3x-6\to 0^+$) ; ne pas écrire $f(2)$ ; dans la question 8, transformer $0$ en $\ln 1$ pour utiliser la croissance de $\ln$.
\end{corrige}

\begin{corrige}[Exercice 10 — Croissance d'une population de bactéries]
\emph{Barème indicatif (sur 8) : Partie A 4 pts ; Partie B 4 pts.}

\textbf{Partie A}
\begin{enumerate}
\item $P(0)=1000\,e^0=1000$. La culture contient initialement \textbf{\num{1000} bactéries}.
\item $P(4)=1000\,e^{0,5\times 4}=1000\,e^2\approx 1000\times 7{,}39=\num{7390}$ bactéries (à l'unité près).
\item On résout $P(t)>\num{100000}$ :
\[ 1000\,e^{0,5t}>\num{100000} \iff e^{0,5t}>100 \iff 0,5t>\ln 100 \iff t>2\ln 100. \]
(La fonction $\ln$ étant strictement croissante, l'inégalité est conservée.) Valeur exacte : $t>2\ln 100=\ln(100^2)=\ln\num{10000}$. Valeur approchée : $t>2\times 4{,}6=9{,}2$ h. La population dépasse \num{100000} individus au bout d'environ \textbf{9,2 heures}.
\end{enumerate}
\textbf{Partie B}
\begin{enumerate}
\item Pour $t\geq 0$, $P(t)>0$, donc $h$ est bien définie et :
\[ h(t)=\ln\left(1000\,e^{0,5t}\right)=\ln 1000+\ln\left(e^{0,5t}\right)=\ln 1000+0,5t, \]
car $\ln(e^u)=u$ pour tout réel $u$ (réciprocité de $\ln$ et $\exp$).
\item $h(t)=0,5t+\ln 1000$ est une fonction \textbf{affine} de $t$ (droite de pente $0,5$ et d'ordonnée à l'origine $\ln 1000$). Le logarithme transforme la croissance exponentielle $t\mapsto 1000e^{0,5t}$ en une croissance \emph{linéaire} : on dit qu'il « linéarise » l'exponentielle. Conséquence pratique : sur un graphique en échelle logarithmique, les données s'alignent sur une droite, ce qui permet de lire directement le taux de croissance ($0,5$ par heure) comme pente.
\item $h(t)=\ln(100000) \iff 0,5t+\ln 1000=\ln 100000 \iff 0,5t=\ln\left(\dfrac{100000}{1000}\right)=\ln 100$, soit $t=2\ln 100\approx 9{,}2$ h. On retrouve exactement le résultat de la question A.3 (la frontière de l'inéquation).
\end{enumerate}
\textbf{Point de vigilance :} dans A.3, ne pas diviser par $\ln$ ; appliquer $\ln$ aux deux membres en justifiant par sa stricte croissance.
\end{corrige}

\begin{corrige}[Exercice 11 — Une inégalité classique et ses applications]
\emph{Barème indicatif (sur 10) : Partie A 5 pts ; Partie B 5 pts.}

\textbf{Partie A : démonstration}
\begin{enumerate}
\item $\varphi$ est dérivable sur $]0\,;\,+\infty[$ comme somme de fonctions dérivables, et :
\[ \varphi'(x)=1-\dfrac{1}{x}=\dfrac{x-1}{x}. \]
Comme $x>0$, le signe de $\varphi'(x)$ est celui de $x-1$ : $\varphi'(x)<0$ sur $]0\,;\,1[$, $\varphi'(1)=0$, $\varphi'(x)>0$ sur $]1\,;\,+\infty[$.
\item $\varphi(1)=1-1-\ln 1=0$. Limites : en $0^+$, $-\ln x\to+\infty$ donc $\varphi(x)\to+\infty$ ; en $+\infty$, $\varphi(x)=x\left(1-\dfrac{1}{x}-\dfrac{\ln x}{x}\right)\to+\infty$ par croissance comparée. Tableau :
\begin{center}
\begin{tabular}{|c|ccccc|}\hline
$x$ & $0$ & & $1$ & & $+\infty$ \\ \hline
$\varphi'(x)$ & & $-$ & $0$ & $+$ & \\ \hline
$\varphi$ & $+\infty$ & $\searrow$ & $0$ & $\nearrow$ & $+\infty$ \\ \hline
\end{tabular}
\end{center}
\item Le tableau montre que $\varphi$ admet sur $]0\,;\,+\infty[$ un \textbf{minimum} en $x=1$, valant $0$. Donc pour tout $x>0$ :
\[ \varphi(x)\geq 0 \iff x-1-\ln x\geq 0 \iff \ln x\leq x-1. \]
L'égalité a lieu \textbf{si et seulement si} $x=1$ (minimum atteint en un seul point, la décroissance puis la croissance étant strictes).
\item \textbf{Interprétation graphique :} la tangente à la courbe de $\ln$ au point d'abscisse $1$ a pour équation $y=\ln'(1)(x-1)+\ln 1=x-1$. L'inégalité $\ln x\leq x-1$ signifie que la courbe de $\ln$ est \textbf{toujours située en dessous de sa tangente} au point $(1\,;\,0)$, avec contact uniquement en ce point. (C'est une conséquence de la concavité de $\ln$.)
\end{enumerate}
\textbf{Partie B : application}
\begin{enumerate}
\item Appliquons l'inégalité de la Partie A au réel $\dfrac{1}{x}>0$ :
\[ \ln\left(\dfrac{1}{x}\right)\leq \dfrac{1}{x}-1 \iff -\ln x\leq \dfrac{1}{x}-1 \iff \ln x\geq 1-\dfrac{1}{x}. \]
En combinant avec $\ln x\leq x-1$, on obtient pour tout $x>0$ :
\[ 1-\dfrac{1}{x}\leq \ln x\leq x-1. \]
\item Avec $x=1{,}08$ : d'une part $x-1=0{,}08$ ; d'autre part $1-\dfrac{1}{1{,}08}=\dfrac{1{,}08-1}{1{,}08}=\dfrac{0{,}08}{1{,}08}\approx 0{,}0741$. Donc :
\[ 0{,}0741\lesssim t\lesssim 0{,}08. \]
La moyenne des bornes donne $t\approx\dfrac{0{,}0741+0{,}08}{2}\approx 0{,}077$. (Valeur de référence : $\ln(1{,}08)\approx 0{,}0770$ à $10^{-4}$ près, bien comprise dans l'encadrement.) Le taux continu équivalent est donc d'environ \textbf{7,7 \%}.
\end{enumerate}
\textbf{Points de vigilance (Bac) :} pour A.3, la conclusion doit être explicitement rattachée au \emph{minimum} de $\varphi$ ; pour B.1, justifier que $\dfrac{1}{x}>0$ avant d'appliquer l'inégalité ; dans B.2, l'encadrement doit être donné dans le bon sens ($1-\frac1x$ est la borne \emph{inférieure}).
\end{corrige}

\newpage

\coursec{Fiche bilan}

\coursec{Fonction logarithme népérien}

\begin{popfigure}
\begin{center}\tcbox[colback=popblue,colframe=popblue,coltext=white,arc=9pt,boxsep=4pt,left=6pt,right=6pt]{\bfseries\small Fonction logarithme népérien $\ln$}\end{center}
\vspace{-2pt}
\begin{tcbraster}[raster columns=3,raster equal height=rows,raster column skip=6pt,raster row skip=6pt,enhanced,blankest]
\begin{tcolorbox}[enhanced,colback=white,colframe=popblue,boxrule=0.9pt,arc=3pt,title={Définition et domaine},fonttitle=\bfseries\scriptsize,coltitle=white,colbacktitle=popblue,left=4pt,right=4pt,top=2pt,bottom=2pt,boxsep=1pt,toptitle=1.5pt,bottomtitle=1.5pt,halign=left]
\begin{itemize}[nosep,leftmargin=*,itemsep=1pt,font=\scriptsize]
\item $\ln x$ défini ssi $x>0$
\item $\ln 1=0$ ; $\ln e=1$
\item Stricte croissante sur $]0,+\infty[$
\end{itemize}
\end{tcolorbox}
\begin{tcolorbox}[enhanced,colback=white,colframe=poporange,boxrule=0.9pt,arc=3pt,title={Propriétés algébriques},fonttitle=\bfseries\scriptsize,coltitle=white,colbacktitle=poporange,left=4pt,right=4pt,top=2pt,bottom=2pt,boxsep=1pt,toptitle=1.5pt,bottomtitle=1.5pt,halign=left]
\begin{itemize}[nosep,leftmargin=*,itemsep=1pt,font=\scriptsize]
\item $\ln(ab)=\ln a+\ln b$
\item $\ln\frac{a}{b}=\ln a-\ln b$
\item $\ln(a^n)=n\ln a$
\item $\ln\sqrt{a}=\frac{1}{2}\ln a$
\end{itemize}
\end{tcolorbox}
\begin{tcolorbox}[enhanced,colback=white,colframe=popgreen,boxrule=0.9pt,arc=3pt,title={Équations et inéquations},fonttitle=\bfseries\scriptsize,coltitle=white,colbacktitle=popgreen,left=4pt,right=4pt,top=2pt,bottom=2pt,boxsep=1pt,toptitle=1.5pt,bottomtitle=1.5pt,halign=left]
\begin{itemize}[nosep,leftmargin=*,itemsep=1pt,font=\scriptsize]
\item $\ln A=\ln B \Leftrightarrow A=B$
\item $\ln A=k \Leftrightarrow A=e^k$
\item Inconnue auxiliaire $X=\ln x$
\end{itemize}
\end{tcolorbox}
\begin{tcolorbox}[enhanced,colback=white,colframe=poppurple,boxrule=0.9pt,arc=3pt,title={Limites usuelles},fonttitle=\bfseries\scriptsize,coltitle=white,colbacktitle=poppurple,left=4pt,right=4pt,top=2pt,bottom=2pt,boxsep=1pt,toptitle=1.5pt,bottomtitle=1.5pt,halign=left]
\begin{itemize}[nosep,leftmargin=*,itemsep=1pt,font=\scriptsize]
\item $\ln x \to +\infty$ en $+\infty$
\item $\ln x \to -\infty$ en $0^+$
\item $\frac{\ln x}{x}\to 0$ en $+\infty$
\item $\frac{\ln(1+x)}{x}\to 1$ en $0$
\end{itemize}
\end{tcolorbox}
\begin{tcolorbox}[enhanced,colback=white,colframe=popgold,boxrule=0.9pt,arc=3pt,title={Dérivation},fonttitle=\bfseries\scriptsize,coltitle=white,colbacktitle=popgold,left=4pt,right=4pt,top=2pt,bottom=2pt,boxsep=1pt,toptitle=1.5pt,bottomtitle=1.5pt,halign=left]
\begin{itemize}[nosep,leftmargin=*,itemsep=1pt,font=\scriptsize]
\item $(\ln u)'=\frac{u'}{u}$ ($u>0$)
\item $(\ln(ax+b))'=\frac{a}{ax+b}$
\end{itemize}
\end{tcolorbox}
\begin{tcolorbox}[enhanced,colback=white,colframe=poppink,boxrule=0.9pt,arc=3pt,title={Primitives},fonttitle=\bfseries\scriptsize,coltitle=white,colbacktitle=poppink,left=4pt,right=4pt,top=2pt,bottom=2pt,boxsep=1pt,toptitle=1.5pt,bottomtitle=1.5pt,halign=left]
\begin{itemize}[nosep,leftmargin=*,itemsep=1pt,font=\scriptsize]
\item $\int \frac{1}{x}=\ln x$ ($x>0$)
\item $\int \frac{u'}{u}=\ln u$ ($u>0$)
\item $\int \frac{a}{bx+c}=\frac{a}{b}\ln|bx+c|$
\end{itemize}
\end{tcolorbox}
\begin{tcolorbox}[enhanced,colback=white,colframe=popblue,boxrule=0.9pt,arc=3pt,title={Étude de $\ln(ax+b)$},fonttitle=\bfseries\scriptsize,coltitle=white,colbacktitle=popblue,left=4pt,right=4pt,top=2pt,bottom=2pt,boxsep=1pt,toptitle=1.5pt,bottomtitle=1.5pt,halign=left]
\begin{itemize}[nosep,leftmargin=*,itemsep=1pt,font=\scriptsize]
\item Domaine : $ax+b>0$
\item Asymptote verticale $x=-\frac{b}{a}$
\item Variations selon signe de $a$
\item Courbe et position relative
\end{itemize}
\end{tcolorbox}
\end{tcbraster}

\legende{Carte mentale de la leçon : la fonction logarithme népérien.}
\end{popfigure}

\courssub{Définition et propriétés fondamentales}

\begin{definition}[Fonction logarithme népérien]
La fonction \cle{logarithme népérien}, notée $\ln$, est l'unique application définie sur $]0\,;+\infty[$, dérivable, telle que :
\[
\forall x>0,\quad (\ln)'(x)=\frac{1}{x} \qquad\text{et}\qquad \ln 1 = 0.
\]
Le nombre \cle{$e$} (constante d'Euler) est l'unique réel vérifiant $\ln e = 1$ ; sa valeur approchée est $e \approx 2,718\,281\,828$.
\end{definition}

\begin{propriete}[Monotonie et signe]
La fonction $\ln$ est \cle{strictement croissante} sur $]0\,;+\infty[$.
Conséquences immédiates (pour $a,b > 0$) :
\begin{itemize}
  \item $\ln a = \ln b \;\Leftrightarrow\; a = b$ \quad (injectivité).
  \item $\ln a < \ln b \;\Leftrightarrow\; a < b$ \quad (ordre conservé).
  \item $\ln x > 0 \;\Leftrightarrow\; x > 1$ \quad ; \quad $\ln x = 0 \;\Leftrightarrow\; x = 1$ \quad ; \quad $\ln x < 0 \;\Leftrightarrow\; 0 < x < 1$.
\end{itemize}
\end{propriete}

\begin{propriete}[Propriétés algébriques]
Pour tous réels strictement positifs $a$, $b$ et tout entier relatif $n$ :
\begin{align*}
\ln(ab) &= \ln a + \ln b \\
\ln\left(\frac{a}{b}\right) &= \ln a - \ln b \quad;\quad \ln\left(\frac{1}{a}\right) = -\ln a \\
\ln(a^n) &= n\ln a \quad;\quad \ln\sqrt{a} = \frac{1}{2}\ln a
\end{align*}
\emph{Attention :} $\ln(a+b) \neq \ln a + \ln b$ et $\ln(a-b) \neq \ln a - \ln b$ (pièges classiques).
\end{propriete}

\begin{propriete}[Limites usuelles et croissances comparées]
\begin{align*}
\lim_{x\to +\infty} \ln x &= +\infty \\
\lim_{x\to 0^+} \ln x &= -\infty \quad\text{(asymptote verticale $x=0$)} \\
\lim_{x\to +\infty} \frac{\ln x}{x} &= 0 \quad\text{($\ln x$ est négligeable devant $x$)} \\
\lim_{x\to 0} \frac{\ln(1+x)}{x} &= 1 \quad\text{(limite remarquable, base de la dérivée)}
\end{align*}
\end{propriete}

\courssub{Équations, inéquations et inconnue auxiliaire}

\begin{methode}[Résolution d'une équation $\ln(\dots)=\dots$]
\begin{enumerate}
  \item \textbf{Domaine :} Écrire la ou les conditions $\text{argument}>0$ et déterminer $D$.
  \item \textbf{Transformation :} Utiliser les propriétés algébriques pour ramener l'équation à l'une des formes canoniques :
        \begin{itemize}
          \item $\ln A = \ln B$ (avec $A,B>0$ sur $D$) $\Rightarrow$ $A = B$.
          \item $\ln A = k$ ($k\in\mathbb{R}$, $A>0$ sur $D$) $\Rightarrow$ $A = e^k$.
        \end{itemize}
  \item \textbf{Résolution :} Résoudre l'équation algébrique obtenue.
  \item \textbf{Vérification :} Ne conserver que les solutions appartenant au domaine $D$.
\end{enumerate}
\end{methode}

\begin{methode}[Résolution d'une inéquation $\ln(\dots)\;<\; \dots$]
\begin{enumerate}
  \item \textbf{Domaine :} Déterminer $D$ (arguments $>0$).
  \item \textbf{Transformation :} Ramener à $\ln A < \ln B$ ou $\ln A < k$ sur $D$.
  \item \textbf{Ordre :} Comme $\ln$ est \textbf{strictement croissante}, le sens de l'inégalité est \textbf{conservé} :
        $\ln A < \ln B \;\Leftrightarrow\; A < B$ \quad et \quad $\ln A < k \;\Leftrightarrow\; A < e^k$.
  \item \textbf{Résolution :} Résoudre l'inéquation algébrique.
  \item \textbf{Conclusion :} Intersecter l'ensemble des solutions avec le domaine $D$.
\end{enumerate}
\end{methode}

\begin{methode}[Équation / inéquation à inconnue auxiliaire]
Lorsque l'inconnue $x$ n'apparaît que sous la forme $\ln x$ (et puissances) :
\begin{enumerate}
  \item Poser $X = \ln x$. Alors $X \in \mathbb{R}$ (pas de restriction sur $X$) et $x = e^X > 0$.
  \item Résoudre l'équation ou l'inéquation en $X$ (souvent polynomiale du second degré).
  \item Revenir à $x$ par $x = e^X$ (toujours $>0$, donc automatiquement dans le domaine).
\end{enumerate}
\end{methode}

\begin{attention}[Pièges à éviter absolument]
\begin{itemize}
  \item \textbf{Oublier le domaine :} Toute résolution doit commencer par « \emph{Domaine : $...>0$ donc $D=...$} ».
  \item \textbf{Appliquer $\ln$ des deux côtés sans vérifier la positivité :} $\ln A = \ln B \Rightarrow A=B$ \textbf{nécessite} $A>0$ et $B>0$.
  \item \textbf{Confondre $\ln(a+b)$ et $\ln a + \ln b$ :} Aucune formule ne simplifie $\ln(a+b)$.
  \item \textbf{Inconnue auxiliaire :} Ne pas oublier le retour $x=e^X$ ; $X$ peut être négatif, $x$ reste positif.
\end{itemize}
\end{attention}

\begin{exoresolu}[Équation et inéquation logarithmiques]
\textbf{Énoncé :} Résoudre dans $\mathbb{R}$ :
\begin{enumerate}
  \item $E_1 : \ln(2x-3) = \ln(x+1)$
  \item $E_2 : \ln(x^2-4) \ge \ln(2x)$
  \item $E_3 : (\ln x)^2 - 3\ln x + 2 = 0$
\end{enumerate}
\textbf{Solution détaillée :}
\tcblower
\textbf{1.} $E_1$ : Domaine : $\begin{cases} 2x-3>0 \\ x+1>0 \end{cases} \Leftrightarrow x>\frac{3}{2}$.
Sur $D$, $\ln(2x-3)=\ln(x+1) \Leftrightarrow 2x-3 = x+1 \Leftrightarrow x=4$.
$4 \in D$, donc $S = \{4\}$.

\textbf{2.} $E_2$ : Domaine : $\begin{cases} x^2-4>0 \\ 2x>0 \end{cases} \Leftrightarrow \begin{cases} x<-2 \text{ ou } x>2 \\ x>0 \end{cases} \Leftrightarrow x>2$.
Sur $D$, $\ln$ croissante $\Rightarrow$ $x^2-4 \ge 2x \Leftrightarrow x^2-2x-4 \ge 0$.
$\Delta = 20$, racines $1\pm\sqrt{5}$. Polynôme $\ge 0$ pour $x \le 1-\sqrt{5}$ ou $x \ge 1+\sqrt{5}$.
Intersection avec $D = ]2,+\infty[$ : $S = [1+\sqrt{5}, +\infty[$.

\textbf{3.} $E_3$ : Inconnue auxiliaire. Poser $X = \ln x$ ($X\in\mathbb{R}$).
$X^2 - 3X + 2 = 0 \Leftrightarrow (X-1)(X-2)=0 \Leftrightarrow X=1$ ou $X=2$.
Retour : $x = e^1 = e$ ou $x = e^2$.
$S = \{e\,;\, e^2\}$.
\end{exoresolu}

\courssub{Calcul différentiel et intégral}

\begin{propriete}[Dérivée de la fonction logarithme]
Soit $u$ une fonction dérivable et strictement positive sur un intervalle $I$.
La fonction $\ln \circ u$ est dérivable sur $I$ et :
\[
(\ln u)'(x) = \frac{u'(x)}{u(x)}.
\]
Cas particulier fondamental : pour $a,b\in\mathbb{R}$, $a\neq 0$, sur le domaine $ax+b>0$ :
\[
\big(\ln(ax+b)\big)' = \frac{a}{ax+b}.
\]
\end{propriete}

\begin{propriete}[Primitives liées au logarithme]
\begin{itemize}
  \item Sur $]0\,;+\infty[$, une primitive de $\dfrac{1}{x}$ est $\ln x$.
  \item Si $u$ est dérivable et $>0$ sur $I$, une primitive de $\dfrac{u'}{u}$ sur $I$ est $\ln u$.
  \item Pour $a,b,c\in\mathbb{R}$, $b\neq 0$, sur tout intervalle ne contenant pas $-\frac{c}{b}$ :
        \[
        \int \frac{a}{bx+c}\,\mathrm{d}x = \frac{a}{b}\ln|bx+c| + k \quad (k\in\mathbb{R}).
        \]
        \emph{Note :} La valeur absolue permet d'étendre la primitive aux intervalles où $bx+c<0$.
\end{itemize}
\end{propriete}

\begin{exoresolu}[Dérivation et recherche de primitives]
\textbf{Énoncé :}
\begin{enumerate}
  \item Dériver $f(x) = \ln(3x^2+2)$ sur $\mathbb{R}$.
  \item Trouver une primitive de $g(x) = \dfrac{4}{5-2x}$ sur $]-\infty, \frac{5}{2}[$.
\end{enumerate}
\textbf{Solution détaillée :}
\tcblower
\textbf{1.} $u(x)=3x^2+2 > 0$ sur $\mathbb{R}$. $u'(x)=6x$.
$f'(x) = \dfrac{6x}{3x^2+2}$.

\textbf{2.} $g(x) = \dfrac{4}{5-2x} = \dfrac{4}{-2x+5}$. Identifier $a=4$, $b=-2$, $c=5$.
Primitive : $\dfrac{4}{-2}\ln|-2x+5| = -2\ln|5-2x|$.
Sur $]-\infty, \frac{5}{2}[$, $5-2x > 0$, on peut écrire $-2\ln(5-2x)$.
\end{exoresolu}

\courssub{Étude complète de $f(x)=\ln(ax+b)$}

\begin{methode}[Étude systématique de $f(x)=\ln(ax+b)$ ($a\neq 0$)]
\begin{enumerate}
  \item \textbf{Domaine $D_f$ :} Résoudre $ax+b > 0$.
        \begin{itemize}
          \item Si $a>0$ : $D_f = ]-\frac{b}{a}, +\infty[$.
          \item Si $a<0$ : $D_f = ]-\infty, -\frac{b}{a}[$.
        \end{itemize}
  \item \textbf{Asymptote verticale :} La droite $x = -\frac{b}{a}$ est une asymptote verticale.
        \begin{itemize}
          \item Si $a>0$ : $\lim_{x\to (-\frac{b}{a})^+} f(x) = -\infty$.
          \item Si $a<0$ : $\lim_{x\to (-\frac{b}{a})^-} f(x) = -\infty$.
        \end{itemize}
  \item \textbf{Dérivée et variations :} $f'(x) = \frac{a}{ax+b}$.
        Le signe de $f'$ est le signe de $a$ (car $ax+b>0$ sur $D_f$).
        \begin{itemize}
          \item $a>0$ : $f$ strictement croissante sur $D_f$.
          \item $a<0$ : $f$ strictement décroissante sur $D_f$.
        \end{itemize}
  \item \textbf{Limites aux bornes de $D_f$ :}
        \begin{itemize}
          \item $a>0$ : $\lim_{x\to -\frac{b}{a}^+} f(x)=-\infty$, $\lim_{x\to +\infty} f(x)=+\infty$.
          \item $a<0$ : $\lim_{x\to -\infty} f(x)=+\infty$, $\lim_{x\to -\frac{b}{a}^-} f(x)=-\infty$.
        \end{itemize}
  \item \textbf{Tableau de variations :} Le dresser avec les bornes, l'asymptote, le sens de variation.
  \item \textbf{Point remarquable :} Abscisse du point où $f(x)=0$ : $ax+b=1 \Rightarrow x=\frac{1-b}{a}$.
  \item \textbf{Tracé de la courbe :} Repérer l'asymptote, le point d'annulation, la concavité ($f''(x)=-\frac{a^2}{(ax+b)^2}<0$ : courbe \emph{concave}).
\end{enumerate}
\end{methode}

\begin{exoresolu}[Étude complète : $f(x)=\ln(4-2x)$]
\textbf{Énoncé :} Étudier la fonction $f(x)=\ln(4-2x)$ et tracer sa courbe $\mathcal{C}_f$.
\textbf{Solution détaillée :}
\tcblower
\textbf{1. Domaine :} $4-2x > 0 \Leftrightarrow 2x < 4 \Leftrightarrow x < 2$. Donc $D_f = ]-\infty, 2[$.
\textbf{2. Asymptote verticale :} $x=2$. Comme $a=-2<0$, $\lim_{x\to 2^-} f(x) = -\infty$.
\textbf{3. Dérivée :} $f'(x) = \frac{-2}{4-2x} = \frac{-2}{2(2-x)} = \frac{-1}{2-x}$.
Sur $D_f$, $2-x > 0$, donc $f'(x) < 0$. $f$ est \textbf{strictement décroissante} sur $]-\infty, 2[$.
\textbf{4. Limites :}
$\lim_{x\to -\infty} \ln(4-2x) = +\infty$ (car $4-2x \to +\infty$).
$\lim_{x\to 2^-} f(x) = -\infty$.
\textbf{5. Annulation :} $f(x)=0 \Leftrightarrow 4-2x=1 \Leftrightarrow 2x=3 \Leftrightarrow x=1,5$.
Point $A(1,5\,;\,0)$.
\textbf{6. Concavité :} $f''(x) = -\frac{4}{(4-2x)^2} < 0$. Courbe concave.
\textbf{7. Tableau de variations :}
\begin{center}
\begin{tabular}{|c|ccccc|}
\hline
$x$ & $-\infty$ & & $1,5$ & & $2^-$ \\
\hline
$f'(x)$ & & $-$ & $-$ & $-$ & \\
\hline
$f(x)$ & $+\infty$ & \searrow & $0$ & \searrow & $-\infty$ \\
\hline
\end{tabular}
\end{center}
\textbf{8. Courbe :} Dans un repère, tracer l'asymptote $x=2$ (pointillés), placer $A(1,5;0)$, dessiner une branche décroissante, concave, partant de $+\infty$ à gauche, passant par $A$, et tombant vers $-\infty$ en approchant $x=2$.
\end{exoresolu}

\begin{savaistu}[Le logarithme népérien au Cameroun et dans le monde]
\begin{itemize}
  \item \textbf{Acoustique (dB) :} Le niveau sonore en décibels est $L = 10\ln(I/I_0)$ (ou $20\ln$ pour la pression). Un concert à Douala à 100 dB correspond à une intensité $I = I_0 e^{10}$.
  \item \textbf{Chimie (pH) :} $\text{pH} = -\log_{10}[\text{H}^+] = -\frac{\ln[\text{H}^+]}{\ln 10}$. L'eau pure (pH 7) a $[\text{H}^+]=10^{-7}$ mol/L.
  \item \textbf{Finance (intérêts composés continus) :} Un capital $C_0$ placé au taux $t$ pendant $n$ ans vaut $C_n = C_0 e^{tn}$. Le temps pour doubler : $n = \frac{\ln 2}{t}$. Au taux 5\% ($t=0,05$), $n \approx \frac{0,693}{0,05} \approx 13,9$ ans.
  \item \textbf{Échelle de Richter (séismes) :} Magnitude $M = \ln(A/A_0)$. Une magnitude 6 libère $e \approx 2,7$ fois plus d'énergie qu'une magnitude 5.
  \item \textbf{Algorithmique :} La complexité $\mathcal{O}(n\ln n)$ (tri fusion, tri rapide) est la meilleure possible pour trier par comparaison.
\end{itemize}
\end{savaistu}

\begin{aretenir}[Checklist avant le Bac]
Avant l'épreuve, vérifie que tu sais :
\begin{itemize}
  \item[$\square$] Donner le domaine de définition de $\ln u$ en résolvant $u>0$.
  \item[$\square$] Citer $\ln 1=0$, $\ln e=1$ et le signe de $\ln x$ selon $x$ par rapport à 1.
  \item[$\square$] Transformer $\ln(ab)$, $\ln\frac{a}{b}$, $\ln(a^n)$, $\ln\sqrt{a}$ sans erreur.
  \item[$\square$] Résoudre $\ln A=\ln B$ et $\ln A=k$ en vérifiant \textbf{systématiquement} le domaine.
  \item[$\square$] Résoudre une inéquation avec $\ln$ en conservant le sens de l'inégalité (croissante).
  \item[$\square$] Utiliser le changement d'inconnue $X=\ln x$ puis revenir par $x=e^X$.
  \item[$\square$] Citer les quatre limites usuelles, dont $\lim\limits_{x\to+\infty}\frac{\ln x}{x}=0$ et $\lim\limits_{x\to 0}\frac{\ln(1+x)}{x}=1$.
  \item[$\square$] Dériver $\ln(ax+b)$ en $\frac{a}{ax+b}$ sans oublier le facteur $a$ (règle de la chaîne).
  \item[$\square$] Trouver une primitive de $\frac{a}{bx+c}$ : $\frac{a}{b}\ln|bx+c|$ (valeur absolue !).
  \item[$\square$] Dresser le tableau de variations complet de $\ln(ax+b)$ avec ses limites aux bornes.
  \item[$\square$] Placer l'asymptote verticale $x=-\frac{b}{a}$ et tracer l'allure de la courbe (concave).
  \item[$\square$] Rédiger proprement : domaine annoncé dès la première ligne de toute résolution.
\end{itemize}
\end{aretenir}

\findecours

\end{document}
